% concatenated from main.tex
\documentclass[journal,10pt]{IEEEtran}  
\IEEEoverridecommandlockouts
%\overrideIEEEmargins     
\usepackage{amsmath,amssymb,bbm,graphicx}
\usepackage[utf8]{inputenc,xcolor}
\usepackage{subfigure,url}
%\usepackage{appendix}

% \usepackage{hyperref}
\usepackage{amsthm}
\newtheorem{definition}{\bfseries Definition}%[section]
\newtheorem{note}{\bfseries Note}%[section]
\newtheorem{proposition}{\bfseries Proposition}%[section]
\newtheorem{example}{\bfseries Example}%[section]
\newtheorem{assumption}{\it Assumption}%[section]
\newtheorem{theorem}{\bfseries Theorem}
\newtheorem{maintheorem}{\bfseries Main Theorem}
\newtheorem{corollary}{\bfseries Corollary}%[section]
\newtheorem{lemma}{\bfseries Lemma}%[section]
\newtheorem{property}{\bfseries Property}
\newtheorem{remark}{\bfseries Remark}
\newtheorem{notation}{\bfseries Notation}
\newtheorem{condition}{\bfseries Condition}
\newtheorem{problem}{\bfseries Problem}

\setlength{\marginparwidth}{2cm}
\allowdisplaybreaks[4]
\newcommand{\vect}[1]{\boldsymbol{\mathbf{#1}}}
\newcommand{\mat}[1]{\begin{bmatrix} #1 \end{bmatrix}}

\newcommand{\R}{\mathbb{R}}
\newcommand{\Z}{\mathbb{Z}} 
\newcommand{\C}{\mathcal{C}} 
\newcommand{\D}{\mathcal{D}} 
\newcommand{\Kzcbf}{K_{\mathrm{zcbf}}}
\newcommand{\la}{\lambda}
\newcommand*\diff{\mathop{}\!\mathrm{d}}
\newcommand{\pa}{\partial}
\newcommand{\sat}{\mbox{sat}}
\newcommand{\lv}{\left\Vert}
\newcommand{\rv}{\right\Vert}
%\newcommand{\d}{\mathrm{d}} 
% Comment Environment
\newif\ifdraft
% To remove comments change this to 
% \draftfalse
\drafttrue
\newcommand{\XX}[1]{\ifdraft{\color{red}[{\bf XX}: #1]}\fi}
%\newcommand{\XX}[2][]{\todo[color=violet!20,size=\scriptsize,#1]{XX: #2}}
\newcommand{\VF}[1]{\ifdraft{\color{blue}[{\bf VF}: #1]}\fi}
\setlength\parindent{0pt}
\setlength{\parskip}{2.5mm}

% \title{\LARGE \bf Convex Framework for Constrained Quadcopter Trajectory Optimization and Tracking with Continuous-Time Guarantees}

% \title{Convex Framework for Constrained Quadcopter Trajectory Optimization and Tracking with Continuous-Time Safe Guarantee}

\title{Flatness-based Quadcopter Trajectory Planning and Tracking with Continuous-time Safety Guarantees}

% \title{Safe Quadcopter Trajectory Generation and Tracking in Continuous-time by Differential Flatness}

% \title{Continuous-time Safety Guarantee for Quadcopter Trajectory Generation and Tracking based on Differential Flatness }

% \title{Flatness-based Quadcopter Trajectory Optimization and Tracking with Continuous-time Safe Guarantees}



\author{Victor Freire,  Xiangru Xu\thanks{Victor Freire and Xiangru Xu are with the Department of Mechanical Engineering, University of Wisconsin-Madison,
        Madison, WI, USA. Email: 
        {\tt\small \{victor.freiremelgizo,xiangru.xu\}@wisc.edu}}}
%\date{March 2021}


\begin{document}
\maketitle

\begin{abstract}
This work presents a convex optimization framework for the planning and tracking of quadcopter trajectories with continuous-time safety guarantees. Using B-spline basis functions and the differential flatness property of quadcopters, a second-order cone program is formulated to generate optimal trajectories that respect safe state and input constraints in the continuous-time sense. A quadratic program (QP) based on control barrier functions is proposed to guarantee bounded trajectory tracking in continuous time by filtering a nominal controller, where the QP is shown to be always feasible. Furthermore, conditions that ensure the safe tracking controller respects thrust, roll angle, and pitch angle constraints are also proposed. The effectiveness of the proposed framework is demonstrated by real-world experiments using a Crazyflie2.1 nano quadcopter. 
%(including position, linear velocity, angle, angular velocity, thrust and waypoints constraints) 
% The generated trajectories are endowed with many layers of safety from state and input constraints to obstacle avoidance. A quadratic program based on control barrier functions guarantees bounded trajectory tracking in continuous-time by filtering a nominal controller. Furthemore, we guarantee the feasibility of the quadratic program and provide conditions to ensure the filtered controller respects input constraints. We present examples using the proposed framework in real-world experiments with a Crazyflie2.1 nano quadcopter. 

The video for the experiments of Section \ref{sec:exper} is available at %\cite{victor21video}. 
\url{https://xu.me.wisc.edu/wp-content/uploads/sites/1196/2021/10/continuous-safety.mp4}.
%\XX{This is the link I created for the video. We can update the video any time, but the size should be less than 24Mb.}\VF{The full video I sent you is 47.3 Mb. I can try lower qualities or trimming some content if that is what we want. YouTube can probably host larger videos, though.}\XX{We only need to include the three examples in the video. }

% \XX{Use `` '' for quotes, not '' ''.}
\end{abstract}

\section{Introduction}



Autonomous quadcopters have been widely used in many fields such as cinematography \cite{cinema}, search and rescue \cite{ludovic2014rescue}, agriculture \cite{navia2016multispectral} and delivery \cite{tang2018aggressive}. Many of these applications are safety-critical and require rigorous satisfaction of safety constraints on the states and inputs of the quadcopters for all time. Although numerous trajectory planning/tracking methods have been proposed for quadcopters, a systematic approach that provides a formal safety guarantee in the continuous-time sense is still largely lacking.
%most of them  provide guarantees at discrete times only, without  \cite{min_snap}.
%because their underlying algorithms are based on state/input discretization and/or are solved at discrete sample times \cite{min_snap,dueri_trajectory_2017}. %The works that target safety in continuous-time tend to provide algorithms that are too conservative or too slow.

The trajectory planning problem is usually split into finding discrete waypoints that avoid obstacles, and generating a smooth curve that respects the system dynamics through the waypoints. However, a trajectory generated this way is normally only safe at the discrete waypoints without a formal safety guarantee in-between. The inter-sampling safety can be improved by increasing the number of samples at the expense of computation, which works well in practice but lacks  a formal continuous-time guarantee \cite{Dontchev_2019,min_snap,polynomial}. Other works addressing the gap between discrete and continuous-time system trajectories include \cite{dueri_trajectory_2017,bonnans2017error,dontchev1981error}; for example, \cite{dueri_trajectory_2017} predicts when safety will be violated during each inter-sample time interval, but it relies on finite violations derived from polynomial dynamics systems.  Another body of literature exploits the differential flatness property of quadcopters and the convexity property of B-spline basis functions to generate formally safe trajectories in continuous-time \cite{Constrained_trajectory,Flat_trajectory_design,stoical2016obstacle,tang2021real}. These works find a sufficiently smooth trajectory in the flat output space which can be mapped to the states and inputs of the system with algebraic transformations. Recent works along this line of research provide continuous-time safety guarantees in the form of state, thrust and angle constraints \cite{Constrained_trajectory,Flat_trajectory_design} and obstacle avoidance \cite{stoical2016obstacle,tang2021real} constraints.  The main challenge with the flatness-based method lies in the highly nonlinear mapping between state/input spaces and flat output space, which makes the consideration of safety constraints in the state/input space challenging. Furthermore, the resulting optimization problem is usually nonconvex, which renders the online trajectory re-planning computationally intractable.


% , however, lie in the highly nonlinear map from states and inputs to the flat outputs\XX{what does this mean?}, and in the high-order time dependence \XX{what does this mean?} resulting in nonconvex optimization problems, which makes the online trajectory re-planning computationally intractable.

\begin{figure}
    \centering
    \includegraphics[width=8.5cm]{block_diagram_summary.PNG}
    \caption{Configuration of the proposed convex optimization-based control framework for quadcopters with continuous-time safety guarantees. The safe trajectory planning is based on solving a second-order cone program, and the safe trajectory tracking is based on solving a quadratic program.}
    \label{fig:block_diagram_summary}
\end{figure}

% The resolution of the continuous-time optimal control problem remains an open research question and the approaches may be broadly divided into indirect and direct methods \cite{rao_survey}: Indirect methods find optimality conditions using calculus of variations to formulate a boundary value problem, which is solved with discrete solvers. Furthermore, computing the gradients and Hamiltonian conditions is not a trivial task. Direct methods, on the other hand, transcribe (discretize) the infinite-dimensional optimal control problem into a nonlinear optimization problem for which solution approaches are well-established. In either case, continuous-time guarantees of the optimal-control solution are seldom provided. Discrete MPC and its variations are commonly-used direct methods that consider state and input constraints. However, They provide no guarantees in-between sample times. This inter-sampling uncertainty was addressed, for example, in \cite{dueri_trajectory_2017}. They predict when each constraint will become active from gradient information to address the violations before they happen. However, their approach relies on finite constraint activations derived from considering polynomial dynamics systems. The usual approach to improve inter-sampling safety is to reduce the size of the discretization time step at the expense of computation. This works well in practice, and works like \cite{Dontchev_2019} focus on speeding-up the convergence of the Newton method for exactly this purpose. However, they are still unable to provide formal continuous-time guarantees. \XX{Error bound between optimal control and MPC due to discretization (papers below commented out). More paper (except for \cite{dueri_trajectory_2017}) that try to generate trajectory with continuous time constraint satisfaction; e.g, see its references. \cite{yu_minimum_2014} formulates min-jerk with B-splines. \cite{usenko_real-time_2017} considers a general min-(velocity, acc, jerk, snap) objective function with B-splines.} 



%The control of quadcopter systems is a well-studied problem in the nonlinear control community \cite{lee2009feedback,mokhtari2005robust, bouabdallah2004design, bouabdallah2004pid}. Furthermore, recent works show improved experimental performance for aggressive flights:  \cite{tal2020accurate} tracks high-order derivatives to reject aerodynamic effects present in aggressive flights, \cite{foehn2018onboard} considers a state-dependent linearization of the quadcopter dynamics to apply LQR techniques when operating far from the equilibrium point, and \cite{FMPC} employs flatness results to feedforward-linearize the MPC problem for optimal control. These controllers are usually endowed with stability (Lyapunov functions) and performance guarantees (optimal control). Safety, however, takes different forms: Controllers based on reinforcement learning have recently gained popularity \cite{lupashin2010simple,singla2019memory}, and the concept of ``safe-learning'' was introduced in \cite{berkenkamp2017safe}. Another considered form of safety is the forward invariance of state space subsets, explored through control barrier functions (CBFs) \cite{CBF_paper} as in \cite{cbf_pos,IBVS} for quadcopters. One recurring problem with hardware-implementation of these controllers is that they usually operate in discrete-time while the system dynamics evolve in continuous-time; Quadcopters are ``sampled-data systems''. This limitation is studied explicitly, for example, in \cite{mitchell2015improved,cortez2019control}. A more common approach, however, to close the gap between controllers operating in discrete-time and systems evolving in continuous-time, is the cascading of algorithms with different sampling rates. For example, \cite{rosolia2020unified} considers three layers of control of increasing complexity such that the time constants for relevant system dynamics are slower than the algorithm layer controlling them.

Recently, there are also many works that aim to improve the tracking performance of quadcopters, particularly for aggressive flights \cite{invernizzi2020comparison,tal2020accurate,foehn2018onboard,FMPC}. 
For example,  \cite{tal2020accurate} tracks high-order derivatives to reject aerodynamic effects present in aggressive flights, \cite{foehn2018onboard} considers a state-dependent linearization of the quadcopter dynamics and applies LQR techniques when the system is operated far from the equilibrium point, and \cite{FMPC} employs flatness results to feedforward-linearize the MPC problem for optimal control. These controllers are usually endowed with stability and/or performance guarantees, but not safety.  Safe controllers based on reinforcement learning and Gaussian processes have been investigated for quadcopters in \cite{lupashin2010simple,singla2019memory,berkenkamp2017safe}. Another line of research that can provide formal safety guarantee in the sense of forward invariance of a given safe set in continuous time is based on control barrier functions (CBFs); see  \cite{ames2016control,cbf_pos,IBVS,rosolia2020unified} for more details. 


% \XX{trajectory tracking algorithms with continuous satisfaction. 
% % paper by Ezra Tal} 

% \XX{A paragraph about trajectory generation algorithms with continuous satisfaction. \\
% Something you should know about optimal control: 1) indirect methods (Pontryagin maximum principle), necessary condition expressed in continuous time, but the boundary value problem needs to be solved by discrete solvers; 2) direct methods, transform into nonlinear program - single shooting (discretize controls), collocation (discretize both controls and states), multiple shooting. \\
% Error bound between optimal control and MPC due to discretization (papers below commented out).\\
% % cite a couple of papers, summarize this line of research
% %[1] Error Estimates for the Euler Discretization of an Optimal Control Problem with First-Order State Constraints
% %[2] Error estimates for a discrete approximation to constrained control problems
% %[3]  Inexact Newton-Kantorovich methods for constrained nonlinear model predictive control
% More paper (except for \cite{dueri_trajectory_2017}) that try to generate trajectory with continuous time constraint satisfaction; e.g, see its references. \\
% B-spline papers for traj? }


% \XX{A few sentences about trajectory tracking algorithms with continuous satisfaction. \\
% paper by Ezra Tal?\\
% CBF papers}

% \XX{Intro is the place to lay the backdrop and show the motivation.  Need more papers in the intro to show we know the literature. The current intro is pale. Too few references - find papers from references or citations. \underline{Read how other (good) papers write intro.} }



This paper proposes a convex optimization-based framework that \emph{guarantees continuous-time safety satisfaction} of quadcopters for both trajectory planning and tracking (see Figure \ref{fig:block_diagram_summary}). The contributions of this paper are summarized as follows:
% \XX{Please change the following paragraph to a summary of contributions.}
\begin{itemize}
    \item We propose sufficient conditions for the $r$-th derivative of clamped B-spline curves to be included by a convex set. 
    %Based on the differential flatness of the quadcopter, we express/relax a set of state and input safety constraints (including position, linear velocity, attitude, thrust, angular velocity and waypoint constraints) into respective constraints in terms of the B-spline curve's control points.
    % Based on the differential flatness of the quadcopter and B-splines for the flat output parameterization, we propose  sufficient conditions for the $r$-th derivative of clamped B-spline curves to be included by a convex set, and express/relax a set of state and input safety constraints (including position, linear velocity, angular, thrust, angular velocity and waypoint constraints) into respective constraints that are convex in terms of the virtual control points.
    \item We present a second-order cone program (SOCP) framework to solve the safe trajectory planning problem, which includes position, linear velocity, angle, angular velocity, thrust and waypoints constraints, such that the reference trajectory generated satisfies the safety constraints rigorously in the continuous-time sense. 
    %, instead of only at discrete sampling times. 
    Obstacle avoidance is also considered within the SOCP framework.
    \item We propose a CBF-QP-based trajectory tracking method for quadcopters such that the real trajectory is in a prescribed tube of the nominal trajectory in the continuous-time sense. We present sufficient conditions to guarantee feasibility of the CBF-QP. Furthermore, we provide conditions to ensure the CBF-QP-based tracking controller respects thrust, roll angle and pitch angle constraints.
    %. Based on the safe reference trajectory, we use CBFs as a safety filter in the trajectory tracking to ensure the real trajectory is kept within a prescribed tube of the nominal trajectory in continuous time.
    % \item We present sufficient conditions for these CBFs to have the control sharing property such that the corresponding convex quadratic program (QP) is always feasible, and for the trajectory such that solutions of the QP respect input constraints.
    \item We demonstrate the proposed optimization-based algorithms through multiple experiments using a Crazyflie2.1 nano-quadcopter.
\end{itemize}

The remainder of the paper is organized as follows: Section \ref{Preliminaries} describes the quadcopter dynamics model,  differential flatness, B-spline curves, CBFs and second-order cone constraints; Section \ref{sec:Constraint} provides sufficient conditions for continuous-time set inclusion constraint satisfaction using clamped B-spline curves; Section \ref{sec:traj} formulates a SOCP for safe trajectory planning; Section \ref{sec:track} describes a convex CBF-based QP for safe trajectory tracking; Section \ref{sec:exper} demonstrates the experiments and finally, Section \ref{sec:Conclusions} concludes the paper.




\section{Preliminaries } \label{Preliminaries}
\subsection{Quadcopter Dynamics \& Differential Flatness} \label{Dynamics}
The equations of motion of a quadcopter used in this work are given as follows \cite{min_snap}: 
\begin{subequations} \label{dyn_model}
\begin{align}
    \ddot{\vect{r}} &\!=\! T\vect{z}_B -g\vect{z}_W, \label{dyn_model_acc}\\
    \dot{\vect{\xi}} &\!=\! N^{-1}(\vect{\xi})\vect{\omega},\\
    N(\vect{\xi}) &\!=\! \mat{1 & 0 & -\sin\theta\\
                         0 & \cos\phi & \sin\phi\cos\theta\\
                         0 & -\sin\phi & \cos\phi\cos\theta},\\
    R_B^W &\!=\! \mat{
c\theta c\psi & s\phi s\theta c\psi - c\phi s\psi & s\phi s\psi + c\phi s\theta c\psi\\
c\theta s\psi & s\phi s\theta s\psi + c\phi c\psi & c\phi s\theta s\psi - s\phi c\psi\\
-s\theta & s\phi c\theta & c\phi c\theta}\\
&= \mat{\vect{x}_B & \vect{y}_B & \vect{z}_B},\nonumber
\end{align}
\end{subequations}
where
%$\vect{z}_W = \mat{0 & 0 & 1}^T$, and  
\begin{align*}
  \vect{z}_W =\mat{0\\0\\1},\quad\vect{r}=\mat{x\\y\\z},\quad  \vect{\xi}=\mat{\phi \\\theta \\\psi},\quad \vect{\omega} =\mat{p \\q \\r},
\end{align*}
are the world-frame's z-axis, the position vector, the Z-Y-X Euler angle vector (see Figure \ref{fig:Coord_Frames}), and the angular velocity vector, respectively, $T$ is the normalized thrust, and $s,c$ stand for $\sin$ and $\cos$ respectively. The considered state and input vectors for a quadcopter are expressed, respectively, as:
\begin{subequations} \label{state_input}
\begin{align}
  \vect{x} &= \mat{x&y&z&\phi&\theta&\psi&\dot{x}&\dot{y}&\dot{z}}^T,\\
  \vect{u} &= \mat{T&p&q&r}^T.
\end{align}
\end{subequations}

\begin{figure}
    \centering
    \includegraphics[width=7cm]{Coord_Frames.png}
    \caption{Inertial (or world) coordinate frame $W$ (black) and body coordinate frame $B$ (blue). The intermediate  coordinate frame $C$ (red) is also shown. Crazyflie2.1 figure from \cite{PX4}.}
    \label{fig:Coord_Frames}
\end{figure}

The quadcopter model shown in \eqref{dyn_model} is known to be differentially flat \cite{min_snap,fliess1995flatness}. 
By choosing flat outputs as:
$$
\vect{\sigma} = \mat{x & y & z & \psi}^T,
$$ 
the state and input can be expressed as a function of $\vect{\sigma}$ and its derivatives for some $\Phi$: 
%$    (\vect{x},\vect{u})=\Phi(\vect{\sigma},\dot{\vect{\sigma}},\ddot{\vect{\sigma}},\dddot{\vect{\sigma}})$, 
\begin{equation}
    (\vect{x},\vect{u})=\Phi(\vect{\sigma},\dot{\vect{\sigma}},\ddot{\vect{\sigma}},\dddot{\vect{\sigma}}),
    \label{flat_map}
\end{equation}
where the expressions of $\Phi$ are given as \cite{min_snap,zhou2014vector}:
\begin{subequations}
\begin{align}
    \phi &= \sin^{-1}\Big(\frac{\vect{z}_W^T\vect{y}_B}{\cos(\sin^{-1}(\vect{z}_W^T\vect{x}_B))} \Big),\\
    \theta &= -\sin^{-1}(\vect{z}_W^T\vect{x}_B),\\
    T &= \sqrt{\ddot{x}^2 + \ddot{y}^2 + (\ddot{z}+g)^2},\label{flat_T_map}\\
    p &= -\vect{y}_B \cdot \vect{h}_{\omega},\label{flat_p_map}\\
    q &= \vect{x}_B \cdot \vect{h}_{\omega},\label{flat_q_map}\\
    r &= \vect{z}_B \cdot \dot{\psi}\vect{z}_W,
\end{align}
\end{subequations}
with intermediate quantities:
\begin{align*}
    \vect{T} &= \mat{\ddot{x} & \ddot{y} & \ddot{z} + g }^T,\\
    \vect{y}_C &= \mat{-\sin\psi & \cos\psi & 0 }^T,\\
    \vect{z}_B &= \frac{\vect{T}}{\|\vect{T}\|_2},\\
    \vect{x}_B &= \frac{\vect{y}_C \times \vect{z}_B}{\|\vect{y}_C \times \vect{z}_B\|_2},\\
    \vect{y}_B &= \vect{z}_B\times\vect{x}_B,\\
    \vect{h}_{\omega} &= \frac{1}{\|\vect{T}\|_2}(\dddot{\vect{r}}-(\vect{z}_B \cdot \dddot{\vect{r}})\vect{z}_B).
\end{align*}
%where the construction of $\Phi   $ can be found in  \cite{min_snap,zhou2014vector}. 

Generating a trajectory for differentially flat systems reduces to finding a sufficiently smooth flat output trajectory \cite{murray1998real}. In the case of system \eqref{dyn_model}, the flat trajectory $\vect{\sigma}(t)$ needs to be at least thrice-differentiable.
%it suffices to find a thrice-differentiable flat trajectory $\vect{\sigma}(t)$. 
%\XX{The objective function in SOCP involves $r^{(4)}$?} \VF{We require $r^{(3)}$ to be continuous ($d=4$) for $\Phi$. The objective function \eqref{min_snap_obj} involves integrating $r^{(4)}$ which is discontinuous for ($d=4$) and continuous for ($d=5$). Since I approximate the integral with a discrete sum, $d=4$ is not a problem for the code.}

% We now expose the geometric meaning of two results from literature on quadcopter differential flatness \cite{effective_ang,raff}. This geometric interpretation will be key for the forthcoming results:


\subsection{B-Spline Curves} \label{B-Splines}
B-splines are widely used to parameterize trajectories because of their nice properties and ease of use. For example, the smoothness of B-splines makes the differentiability constraint of the flat trajectory $\vect{\sigma}(t)$ easy to satisfy, the convexity property allows one to explore state and input constraints in continuous time, and the locality property is useful for re-planning.

%\begin{definition}\cite{deboor}
Given a \emph{knot vector} $\vect{\tau} = (\tau_0, \ldots,\tau_v)^T$ satisfying $\tau_i \leq \tau_{i+1}$ where $i=0,\ldots,v-1$, a $d$-th degree \emph{B-spline basis} with  $d\in\Z_{> 0}$ is computed recursively as \cite{deboor}:
\begin{align*}
    \lambda_{i,0}(t) &= \begin{cases} 
      1, & \tau_i \leq t <\tau_{i+1}, \\
      0, & \text{otherwise}, \\
   \end{cases}\\
   \lambda_{i,d}(t)&=\frac{t-\tau_i}{\tau_{i+d}-\tau_i}\lambda_{i,d-1}(t)+ \frac{\tau_{i+d+1}-t}{\tau_{i+d+1}-\tau_{i+1}}\lambda_{i+1,d-1}(t).
\end{align*}
In this work, we consider the \emph{clamped, uniform B-spline basis}, which is defined over knot vectors satisfying:
\begin{subequations} \label{clamped_uniform}
\begin{align}
    \text{(clamped)} \quad &\tau_0=\ldots=\tau_d, \quad \tau_{v-d}=\ldots=\tau_{v},\label{clamped}\\
    \text{(uniform)}\quad &\tau_{d+1}-\tau_{d}=\ldots =\tau_{N+1}-\tau_{N},\label{uniform}
\end{align}
\end{subequations}
where $N = v-d-1$. A $d$-th degree \emph{B-spline curve} $\vect{s}(t)$ is a $m$-dimensional parametric curve  built by linearly combining \emph{control points} $\vect{P}_i\in\mathbb{R}^m (i=0,\ldots,N)$ and B-spline basis functions of the same degree. Noting $\vect{s}^{(0)}(t) = \vect{s}(t)$, we generate a B-spline curve and its $r$-th order derivative by:
\begin{equation}
    \vect{s}^{(r)}(t) = \sum_{i=0}^{N}\vect{P}_i\vect{b}_{r,i+1}^T\vect{\Lambda}_{d-r}(t) = PB_r\vect{\Lambda}_{d-r}(t),
    \label{b-spline_curve}
\end{equation}
where the control points are grouped into a matrix
\begin{equation} \label{defP}
    P = \mat{\vect{P}_0&\dots& \vect{P}_N} \in\R^{m\times (N+1)},
\end{equation}
the basis functions are grouped into a vector
\begin{equation}
    \vect{\Lambda}_{d-r}(t) = \left[\lambda_{0,d-r}(t),\dots ,\lambda_{N+r,d-r}(t)\right]^T\in\R^{N+r+1},
\end{equation}
and $\vect{b}_{r,j}^T$ is the $j$-th row of a time-invariant matrix $B_{r}$ whose construction is given in Appendix \ref{B_mat}. %(see \cite{fajar_thesis} for more detail). %\VF{I also think it is redundant to specify the dimensions of $P$ and $\vect{\Lambda}_{d-r}(t)$}.

% Without loss of generality\VF{I think we do lose generality. The derivatives depend a lot on the size of the knot intervals. Clamped, uniform B-splines are a special case of B-splines.}, we assume the segments $[\tau_i,\tau_{i+1})$ are nonempty and uniform when $d\leq i\leq N = v-d-1$.

% \VF{I believe the following to be the most clear and concise way to define B-spline curves and their derivatives. The reason is the special structure of vector $\vect{\Lambda_{d-r}(t)}$, which changes dimensions based on $r$. I commented out your modified version.}


% A \emph{B-spline curve} $\vect{s}(t)$ of order $d$ is a $m$-dimensional parametric curve generated from B-spline bases of the same order by:
% \begin{equation} \label{b_spline_curve}
%     \vect{s}(t) =  P\vect{\Lambda}_{d}(t) 
% \end{equation}
% where 
% \begin{equation}
% P = \mat{\vect{P}_0&\dots& \vect{P}_N} \in\R^{m\times N+1} \label{defP}
% \end{equation}
% with $\vect{P}_i\in\R^{m} (i=0,...,N)$ the control points, and
% \begin{align}
% \vect{\Lambda}_{d-r}(t) &= \mat{\lambda_{0,d-r}(t)&\dots &\lambda_{N+r,d-r}(t)}^T\in\R^{?\times ?}.
% \end{align}
% The $r$-th order derivatives $\vect{r}^{(r)}(t)$ with $0\leq r \leq d$ is
% \begin{equation}
%     \vect{r}^{(r)}(t) = \sum_{i=0}^{N}\vect{P}_i\vect{b}_{r,i+1}^T\vect{\Lambda}_{d-r}(t) = PB_r\vect{\Lambda}_{d-r}(t)
%     \label{b-spline_curve_der}
% \end{equation}
% where the vector $\vect{b}_{r,j}^T$ is the $j$-th row of the matrix $B_{r}$, whose construction  is discussed in Appendix \ref{B_mat}.
%\end{definition}
%Since $B_{0} = I_{N+1}$, it follows that $\vect{r}^{(0)}(t) = \vect{r}(t)= P\vect{\Lambda}_{d}(t) $. 

%\XX{what does this equation mean? I think you are just defining $\vect{r}(t)$ to be $\vect{r}^{(0)}(t)$?}  
%Denote $P^{(r)} = PB_{r}$. 
% call the columns of $PB_{r}$ as the \emph{$r$-th order virtual control points}; that is, 
% \begin{equation}\label{virtual_cp_eq}
%     P^{(r)} = PB_{r} = \mat{\vect{P}_0^{(r)} & \dots & \vect{P}_{N+r}^{(r)}}
% \end{equation}


\begin{definition}\label{virtual_control_points_old} The columns of $P^{(r)} \triangleq PB_{r}$ are called the \emph{$r$-th order virtual control points} (VCPs) of $\vect{s}^{(r)}(t)$ and denoted as $\vect{P}_i^{(r)}$ where $i=0,1,\ldots,N+r$, i.e., 
\begin{equation}\label{virtual_cp_eq}
    P^{(r)} = PB_{r} = \mat{\vect{P}_0^{(r)} & \dots & \vect{P}_{N+r}^{(r)}}.
\end{equation}
\end{definition}
Note that $P^{(0)} = P$ and that the first and last $r$ columns of $P^{(r)}$ are zero vectors because of the form of $B_{r}$. 

%Properties of B-splines are discussed in detail in \cite{deboor,fajar_thesis}. 
Some properties of B-splines that will be used in the following sections are listed below (see \cite{deboor,fajar_thesis} for more details). 
\begin{itemize}
    \item[P1)] \textit{Continuity}. Given a clamped knot vector satisfying \eqref{clamped},  $\vect{s}(t)$ is $C^{\infty}$ for any  $t\notin \vect{\tau}$ and $C^{d-1}$ for any $t\in\vect{\tau}$.
    \item[P2)] \textit{Convexity}. The B-Spline basis functions have the ``partition of unity'' property, meaning that: 
    %(i) $\lambda_{i,d}(t) \geq 0$; (ii) $\sum_{i=0}^{N}\lambda_{i,d}(t)= 1, \ \forall t\in[\tau_0,\tau_v)$.
    \begin{itemize}
        \item[(i)] $\lambda_{i,d}(t) \geq 0$;
        \item[(ii)] $\sum_{i=0}^{N}\lambda_{i,d}(t)= 1, \ \forall t\in[\tau_0,\tau_v)$.
    \end{itemize}
    \item[P3)] \textit{Local Support}. Any B-spline basis is only nonzero over a local set of knots: $\lambda_{i,d}(t)\neq 0$ iff $t\in[\tau_i,\tau_{i+d+1})$.
    \item[P4)] \textit{Strong Convexity}. Segments of the B-spline curve are contained within the convex hull of a limited set of control points. Combining each segment we obtain 
    %$ \vect{s}(t) \in \bigcup_{i=d}^{N}\text{Conv}\{\vect{P}_{i-d},\dots,\vect{P}_i\}.$
    \begin{equation*}
        \vect{s}(t) \in \bigcup_{i=d}^{N}\text{Conv}\{\vect{P}_{i-d},\dots,\vect{P}_i\}.
    \end{equation*}
\end{itemize} 
% \XX{I think the standard notation for convex hull is \text{Conv}.}

% \begin{itemize}
%     \item[P1)] \textit{Continuity}. Given a clamped and uniform knot vector satisfying \eqref{clamped_uniform},  $\vect{s}(t)$ is $C^{\infty}$ for any  $t\notin \vect{\tau}$ and $C^{d-1}$ for any $t\in\vect{\tau}$.
%     \item[P2)] \textit{Convexity}. The B-Spline basis functions have ``partition of unity'' property, meaning that: 
%     %(i) $\lambda_{i,d}(t) \geq 0$; (ii) $\sum_{i=0}^{N}\lambda_{i,d}(t)= 1, \ \forall t\in[\tau_0,\tau_v)$.
%     \begin{itemize}
%         \item[(i)] $\lambda_{i,d}(t) \geq 0$;
%         \item[(ii)] $\sum_{i=0}^{N}\lambda_{i,d}(t)= 1, \ \forall t\in[\tau_0,\tau_v)$.
%     \end{itemize}
%     \item[P3)] \textit{Local Support}. Any B-spline basis is only nonzero over a local set of knots: $\lambda_{i,d}(t)\neq 0$ iff $t\in[\tau_i,\tau_{i+d+1})$.
%     \item[P4)] \textit{Strong Convexity}. Segments of the B-spline curve are contained within the convex hull of a limited set of control points. Combining each segment we obtain 
%     %$ \vect{s}(t) \in \bigcup_{i=d}^{N}\text{Conv}\{\vect{P}_{i-d},\dots,\vect{P}_i\}.$
%     \begin{equation*}
%         \vect{s}(t) \in \bigcup_{i=d}^{N}\text{Conv Hull}\{\vect{P}_{i-d},\dots,\vect{P}_i\}.
%     \end{equation*}
% \end{itemize} 

\subsection{Control Barrier Functions}
% Consider an affine control system 
% \begin{eqnarray}
% \label{eqnsys}
% \dot{x} = f(x) + g(x) u
% \end{eqnarray}
% where $x \in \R^n$, $u \in U \subset \R^m$, $f:\R^n\rightarrow \R^n$ and $g:\R^n\rightarrow \R^m$ are locally Lipschitz continuous. A set $\mathcal{S}$ is called forward controlled invariant with respect to system \eqref{eqnsys} if for every $x_0 \in \mathcal{S}$, there exists a control signal $u(t)$ such that $x(t;t_0,x_0) \in \mathcal{S}$ for all $t\geq t_0$. %In this paper, it will be assumed that $U=\R^m$ (i.e., there is no constraint on the input).
%{\color{blue}It is assumed that $f(0)=0$.}

The concept of (zeroing)  control barrier functions (CBFs) was first proposed in \cite{Xu2015ADHS}, and was proven to be a powerful control design method to ensure the safety (in the sense of controlled invariance) of control affine systems \cite{ames2016control,xu2018correctness}. A provably safe control law is generated by solving a CBF-QP online \cite{Xu2015ADHS}. In this work, we will use time-varying CBFs with relative degree 2 and their control sharing property \cite{xu2018constrained}. %The definition of CBFs with relative degree $r$ is given as follows.

% \begin{definition}\label{dfn:newcbf}\cite{Xu2015ADHS}
% Consider control system \eqref{eqnsys} and a set $\mathcal{C} \subset \R^n$ defined by 
% \begin{equation}
%     \mathcal{C} = \{ x \in \R^n : h(x) \geq 0\}
% \end{equation}
% for a continuously differentiable function $h: \R^n \to \R$ where $h$ has relative degree 1. The function $h$ is called a zeroing control barrier function (CBF) defined on a  set $\mathcal{D}$ with $\C\subseteq\mathcal{D}\subset \R^n$, if  there exists a constant $\gamma>0$ such that
% %\footnote{A more general definition for the (zeroing) CBFs that involves extended class $\mathcal{K}$ functions can be found in Def. 6 in \cite{Xu2015ADHS}.}
% \begin{align}\label{ineqZCBF}
% & \sup_{u \in U}  \left[ L_f h(x) + L_g h(x) u + \gamma h(x)\right] \geq 0,\;\forall x \in \mathcal{D},
% \end{align}
% where $L_fh(x)=\frac{\partial h}{\partial x}f(x)$ and $L_gh(x)=\frac{\partial h}{\partial x}g(x)$ are Lie derivatives.
% 	%	The ZCBF $h$ is said to be locally Lipschitz continuous if $\alpha$ and the derivative of $h$ are both locally Lipschitz continuous.
% \end{definition}

% %If $U=\R^m$ and $L_gh(x) \ne 0$ for $x\in\mathcal{D}$, then the function $h$ is always a ZCBF.

% Given a CBF $h$ of relative degree 1, the set of all control values that satisfy \eqref{ineqZCBF} for all $x\in\D$ is defined as $\Kzcbf(x) =  \{ u \in U : L_f h(x) + L_g h(x) u + \gamma h(x) \geq 0\}.$ 
% %Given a CBF $h_0$, for all $x\in\D$, define the set of all control values that satisfy \eqref{ineqZCBF}:
% % \begin{equation}\label{zcbfinputset}
% % \Kzcbf(x) =  \{ u \in U : L_f h_0(x) + L_g h_0(x) u + \gamma h_0(x) \geq 0\}.
% % \end{equation}
% It was proven in \cite{Xu2015ADHS} that any Lipschitz continuous controller $u: \mathcal{D} \to U$ such that $u(x) \in \Kzcbf(x)$ for every $x\in\mathcal{D}$ will guarantee the forward invariance of $\mathcal{C}$. The provably safe control is obtained by solving online a QP that include the control barrier condition as its constraint.

% \begin{definition}\cite{xu2018constrained}\label{dfn:cbfhigh}
% Given a control system \eqref{eqnsys}, a $C^r$ function $h(x): \R\times \R^n \to \R$ with a relative degree $r$ is called a zeroing control barrier function of order $r$ if  there exists a column vector ${\bf a}\in\R^r$ such that $\forall x\in\R^n$,
% \begin{align}\label{ineq:ZCBF2}
% 	& \sup_{u \in U}  [L_g L_f^{r-1}h(x)u +L_f^rh(x)+{\bf a}'\eta(x) ] \geq 0,
% \end{align}
% where $\eta(x)=(L_f^{r-1}h, L_f^{r-2}h,...,h)'\in\R^r$, and ${\bf a}=(a_1,...,a_r)'\in\R^r$ is chosen such that the roots of $p_0^r(\la)=\la^r+a_1\la^{r-1}+...+a_{r-1}\la+a_r$ are all negative.
% \end{definition}
% Define functions $s_k(x(t))$ for $k=0,1,...,r$ as follows:
% \begin{align}\label{liftedh}
% s_0(x(t))&=h(x(t)),\;s_{k}(x(t))=(\frac{\diff}{\diff t}+\la_k)s_{k-1},1\le k\le r.
% \end{align}

Consider a time-varying affine control system
\begin{align}\label{eqn:controlsys}
\dot{\vect{x}} = f(t,\vect{x}) + g(t,\vect{x})u,
\end{align}
where $\vect{x} \in \R^n$, $u \in U\subset \R$ and $f:\R\times \R^n \to \R^n$, $g:\R\times \R^n \to \R^n$ are piecewise continuous in $t$ and locally Lipschitz in $\vect{x}$. 
Given a sufficiently smooth time-varying function $h(t,\vect{x}):\R\times \R^n \to \R$, we define the modified Lie derivative of $h(t,\vect{x})$ along $f$ as $\bar L_f^ih:=(\frac{\pa }{\pa t}+L_f)^ih$  where $i$ is a non-negative integer. 

\begin{definition}\cite{xu2018constrained}\label{dfn:cbfhigh}
Given a control system \eqref{eqn:controlsys}, a $C^r$ function $h(t,\vect{x}): \R\times \R^n \to \R$ with a relative degree $r$ is called a CBF if  there exists a column vector ${\bf a}=(a_1,\ldots,a_r)^T\in\R^r$ such that for all $\vect{x}\in\R^n$, $t\geq 0$,
\begin{align}\label{ineq:ZCBF2}
	& \sup_{u \in U}  [L_g\bar L_f^{r-1}h(t,\vect{x})u +\bar L_f^rh(t,\vect{x})+{\vect{a}}^T\eta(t,\vect{x}) ] \geq 0,
\end{align}
where $\eta(t,\vect{x})=(\bar L_f^{r-1}h,\bar L_f^{r-2}h,\ldots,h)^T\in\R^r$, and the roots of $p_0^r(\la)=\la^r+a_1\la^{r-1}+...+a_{r-1}\la+a_r$ are negative reals $-\la_1,...,-\la_r<0$.
\end{definition}

Define a series of functions:%\VF{Redefining to avoid confusion with B-spline $\vect{s}^{(r)}(t)$}
\begin{align}\label{liftedh}
\nu_0(t,\vect{x})&=h(t,\vect{x}),\;\nu_{k}=\Big(\frac{\diff}{\diff t}+\la_k\Big)\nu_{k-1},\;1\le k\le r.
\end{align}
For any $t\geq 0,\vect{x}\in \R^n$, define the set of admissible inputs satisfying the CBF condition as $\mathcal{K}_h(t,\vect{x}):=\{u\in U|L_g\bar L_f^{r-1}h(t,\vect{x})u +\bar L_f^rh(t,\vect{x})+\vect{a}^T\eta(t,\vect{x}) \ge 0\}$. 

\begin{lemma}\cite{xu2018constrained}\label{thm:zcbf}
Given a control system \eqref{eqn:controlsys}, if 1) $h(t,\vect{x})$ with a relative degree $r$ is a CBF such that \eqref{ineq:ZCBF2} holds and the roots of $p_0^r(\la)=\la^r+a_1\la^{r-1}+...+a_{r-1}\la+a_r$ are  $-\la_1,...,-\la_r<0$, 2) $\nu_i$ defined in \eqref{liftedh} satisfy $\nu_i(0,\vect{x}_0)\ge 0$ for $i=0,1,...,r-1$, then any controller $u(t,\vect{x}) \in \mathcal{K}_h(t,\vect{x})$ that is Lipschitz in $\vect{x}$ will render $h(t,\vect{x})\ge 0$ for any $t\geq 0$.
\end{lemma}

Given a control system \eqref{eqn:controlsys} and $q (q\geq 2)$ CBFs $h_i(t,\vect{x}),i=1,...,q$, the \emph{control-sharing property} of CBFs was proposed in \cite{xu2018constrained} to ensure the feasibility of the CBF-QP that includes multiple CBF constraints. We refer the reader to Definition 2 in \cite{xu2018constrained} for more detail.
% \begin{definition}\cite{xu2018constrained}\label{dfn:mulcbfs}
% Consider control system \eqref{eqn:controlsys} and $q(q\ge 2)$ CBFs $h_i(t,\vect{x}): \R\times \R^n \to \R$, $i=1,...,q$, where $h_i$ has a well-defined (global, uniform) relative degree $r_i$, that is, $L_g\bar L_f^{r-1}h(t,\vect{x})\neq 0$, $L_g\bar L_f^{i}h(t,\vect{x})=0,\;i=0,1,...,r-2$, for any $t\in \R$,  $x\in\R^n$. Suppose that \eqref{ineq:ZCBF2} holds with $\eta_i(t,\vect{x})=(\bar L_f^{r_i-1}h_i,\bar L_f^{r_i-2}h_i,...,h_i)$ for some
% ${\bf a}_i=(a_1^i,\ldots,a_{r_i}^i)$ where the roots of $p_0^{r_i}(\lambda)$ are all negative. 
% The CBFs $h_1,...,h_q$ are said to have the \emph{control-sharing property} if for any $t\geq 0,\vect{x}\in \R^n$, there exists control $\vect{u}\in \R$ such that the following condition holds for any $i=1,...,q$: 
% \begin{align}
% &L_g\bar L_f^{r_i-1}h_i(t,\vect{x})u +\bar L_f^{r_i}h_i(t,\vect{x})+{\bf a}_i^T\eta_i(t,\vect{x})\ge 0. \label{sharecon}
% \end{align}
% \end{definition}

% Given a CBF $h$ of relative degree $r$, if  $s_k(x(0))\geq 0$ for $k=0,1,...,r$, where $s_k$ are given in \eqref{liftedh}, then any Lipschitz continuous controller $u(x)\in \{u\in U \mid L_g L_f^{r-1}h(x)u +L_f^rh(x)+{\bf a}'\eta(x)  \geq 0,x\in \mathcal{D}\}$ for every $x\in\mathcal{D}$ will guarantee the forward invariance of $\mathcal{C}$ \cite{xu2018constrained}.

\subsection{Second-order Cone Constraints} \label{SOC_OPT} 
A \emph{second-order cone program} (SOCP) is a type of convex optimization problem of the following form \cite{lobo1998applications,alizadeh2003second,boyd2004convex}: 
\begin{align} \label{SOCP_form}
    \min \quad & \vect{f}^T\vect{x}\\
    \text{s.t.}\quad &\|A_i\vect{x}+\vect{b}_i\|_2\leq \vect{c}_i^T\vect{x}+d_i, \quad i=1,\ldots,m,\nonumber
%    &F\vect{x} = \vect{g},\nonumber
\end{align}
where $\vect{x}\in\mathbb{R}^n$ is the optimization variable, and the problem parameters are $\vect{f}\in\R^n,A_i\in\mathbb{R}^{(n_i-1)\times n},\vect{b}_i\in\R^{n_i-1},\vect{c}_i\in\R^{n}$ 
% \VF{Why $n_i-1?$ I think just $n_i$ is sufficient.}\XX{See e.g., \cite{lobo1998applications}.} 
and $d_i\in\R$. Constraints of the form:
\begin{equation} \label{SOC}
    \|A\vect{x}+\vect{b}\|_2\leq \vect{c}^T\vect{x}+d
\end{equation}
are called \emph{second-order cone} (SOC) constraints  of dimension $n$. Many commonly-seen convex constraints can be formulated as SOC; for example, ellipsoidal constraints ($\vect{x}^TQ\vect{x} \leq r,\ Q\succeq0$) and polyhedral constraints ($A\vect{x}\leq \vect{b}$) \cite{lobo1998applications,alizadeh2003second,boyd2004convex}. SOCP encompasses many types of important convex programs as special cases, such as linear programs, convex QPs and convex quadratically constrained QPs. SOCP can be solved in polynomial time via the interior-point
method and off-the-shelf solvers such as MOSEK exist \cite{potra2000interior,mosek}.

% \subsection{Problem Statement}
% The problem that will be studied in this paper is the following:
% \XX{Finish it}

\section{Continuous-Time Set Inclusion of B-spline Curves \& Their Derivatives}\label{sec:Constraint}
In this section, continuous-time convex set inclusion of B-spline curves and their derivatives will be considered. These results provide an intuitive description of the strong convexity property of B-spline curves and their derivatives, which is the key to enforcing the continuous-time state and input constraint satisfaction of quadcopter trajectories.

% constraints. useful to convert the input and state constraints of quadcopters into SOC constraints in Section \ref{sec:traj}. 
% \begin{definition}[$r$-th order Virtual Control Points] \label{virtual_control_points}
% Given a B-spline curve $\vect{r}(t)$ defined as in \eqref{b-spline_curve} with $r=0$, we group its $r$-th order virtual control points into matrix:
% \begin{equation}\label{virtual_cp_eq}
%     P^{(r)} = PB_{r} = \mat{\vect{P}_0^{(r)} & \dots & \vect{P}_{N+r}^{(r)}}
% \end{equation}
% \end{definition}


%The following Lemma is inspired by Lemma 1 in \cite{Constrained_trajectory} and provides an intuitive description of the strong convexity properties of B-Spline curves and their derivatives.
%within the context of Definition \ref{virtual_control_points}. 
% This result will be key to guarantee continuous-time verification of constraints.
\begin{lemma} \label{smooth_strong_convexity}
The $r$-th derivative of a clamped B-spline curve defined as in \eqref{b-spline_curve} is contained within the convex hull of its $r$-th order virtual control points such that:
\begin{equation}
    \vect{s}^{(r)}(t) \in \text{Conv}\{\vect{P}_{0}^{(r)},\dots,\vect{P}_{N+r}^{(r)}\}, \quad t\in [\tau_0,\tau_v).
    \label{convex}
\end{equation} 
% \VF{To be precise, we need to specify time interval $t\in[\tau_0,\tau_v)$ right-side-open because of the recursive definition of B-spline bases (see above \eqref{clamped_uniform}).}
Furthermore, for $d\leq i\leq N$, %we have that:
\begin{equation}        
    \vect{s}^{(r)}(t) \in \text{Conv}\{\vect{P}_{i-d+r}^{(r)},\dots,\vect{P}_{i}^{(r)}\}, \quad  t\in[\tau_i,\tau_{i+1}).
    \label{strong_convex}
\end{equation}
By \eqref{strong_convex}, condition \eqref{convex} can be reduced to:
\begin{equation}
    \vect{s}^{(r)}(t) \in \text{Conv} \{\vect{P}_{r}^{(r)},\dots,\vect{P}_N^{(r)}\}, \quad t\in[\tau_0,\tau_v).
    \label{smooth_strong_convex}
\end{equation}
\end{lemma}
\begin{proof}
Applying 
%Definition \ref{virtual_control_points}
the definition of VCP to \eqref{b-spline_curve} we obtain:
\begin{equation}
    \vect{s}^{(r)}(t) = P^{(r)}\vect{\Lambda}_{d-r}(t) = \sum_{i=0}^{N+r} \lambda_{i,d-r}(t)\vect{P}^{(r)}_i.
    \label{lin_comb}
\end{equation}
Applying P2) of B-splines shown in Section \ref{B-Splines}, we determine that $\vect{s}^{(r)}(t)$ is a convex combination of $r$-th order VCPs, which is equivalent to \eqref{convex}. Using P3) of B-splines we have %$\sum_{j=i-d+r}^{i}\lambda_{j,d-r}(t)=1,  t\in[\tau_i,\tau_{i+1}), d\leq i \leq N.$ 
\begin{equation*}
    \sum_{j=i-d+r}^{i}\lambda_{j,d-r}(t)=1,\quad  t\in[\tau_i,\tau_{i+1}), \quad d\leq i \leq N.
\end{equation*}
Leveraging this result, \eqref{lin_comb} reduces to:
\begin{equation*}
    \vect{s}^{(r)}(t) = \sum_{j=i-d+r}^{i} \lambda_{j,d-r}(t)\vect{P}^{(r)}_j, \forall t\in [\tau_i,\tau_{i+1}), d\leq i\leq N,
\end{equation*}
which is equivalent to \eqref{strong_convex}. Taking the union of all nonempty knot segments in the knot vector $\vect{\tau}$ results in \eqref{smooth_strong_convex}. %This completes the proof.
\end{proof}


\begin{proposition} \label{inclusion_prop}
Given a convex set $\mathcal{S}$ and the $r$-th derivative of a clamped B-spline curve $\vect{s}^{(r)}(t)$ defined as in \eqref{b-spline_curve}, if
\begin{equation}
     %\mbox{i)}\quad \quad
     \vect{P}_{j}^{(r)} \in \mathcal{S}, \quad j=r,\ldots,N
    \label{inclusion_global}
\end{equation}
holds, then $\vect{s}^{(r)}(t) \in \mathcal{S}, \ t\in[\tau_0,\tau_v)$. If
\begin{equation}
    %\mbox{ii)}\quad\quad
    \vect{P}^{(r)}_j \in S, \quad j=i-d+r,\dots,i, \quad i\in\{d,\ldots,N\}
    \label{inclusion_local}
\end{equation}
holds, then $\vect{s}^{(r)}(t)\in \mathcal{S}, \ t\in[\tau_i,\tau_{i+1})$.
\end{proposition}

% \begin{proposition} \label{inclusion_prop}
% Given a convex set $S$ and the $r$-th derivative of a clamped B-spline curve defined as in \eqref{b-spline_curve}, a sufficient condition for $\vect{s}^{(r)}(t) \in S$ for all
% \begin{itemize}
%     \item[i)] $t\in[\tau_0,\tau_v)$
% is that the following holds:
% \begin{equation}
%     \vect{P}_{j}^{(r)} \in S, \quad j=r,\ldots,N,
%     \label{inclusion_global}
% \end{equation}
% \item[ii)] $t\in[\tau_i,\tau_{i+1})$ with $i=d,\ldots,N$, is that the following holds:
% \begin{equation}
%     \vect{P}^{(r)}_j \in S, \quad j=i-d+r,\dots,i.
%     \label{inclusion_local}
% \end{equation}
% \end{itemize}
% \end{proposition}
\begin{proof}
     From \eqref{smooth_strong_convex} and the convexity property, it follows that containing $\vect{s}^{(r)}(t)$ in a convex set $\mathcal{S}$ for $t\in[\tau_0,\tau_v)$ can be achieved by containing $\vect{P}_j^{(r)}$ in $\mathcal{S}$ for $j=r,\ldots,N$. This results in \eqref{inclusion_global}. Furthermore, containing $\vect{s}^{(r)}(t)$ in $\mathcal{S}$ for a nonempty time interval $[\tau_i,\tau_{i+1})$ reduces to containing $d-r+1$ $r$th order VCPs in $\mathcal{S}$ (i.e., those that form the convex hull in \eqref{strong_convex}). This is equivalent with \eqref{inclusion_local}.
\end{proof}

\begin{remark}
Proposition \ref{inclusion_prop} above generalizes Proposition 1 in \cite{Constrained_trajectory} from an interval set inclusion to a more general convex set inclusion. This will be key when considering spherical and conic sets to ensure the safety of quadcopters as will be shown in the following sections.
\end{remark}

\section{Trajectory Planning with Continuous-Time Safety Guarantee} \label{sec:traj}
This section presents a SOCP framework to solve a safe trajectory planning problem that includes various state and input safety constraints, such that the reference trajectory  satisfies the safety constraints rigorously in the continuous-time sense. 
% This section investigates how state and input safety constraints can be converted into convex constraints in terms of a B-spline curve's control points, and how a SOCP can be formulated to generate a reference trajectory such that all the safety constraints are satisfied in continuous-time. 
%safe trajectories in continuous-time can be generated by solving a SOCP. %Proposition \ref{inclusion_prop} is used to generate trajectories that satisfy state and input constraints in continuous time. 
Since the flat output $\psi$ is irrelevant to the state and input constraints considered in this paper, we set it to zero and limit ourselves to finding a sufficiently smooth, position flat output trajectory defined as
$$
\vect{r}(t) = \mat{x(t) & y(t) & z(t)}^T.
$$ 

\subsection{Encoding Safety Specifications as SOC Constraints}
Consider the following state and input safety constraints that have to be satisfied rigorously in continuous-time:
\begin{subequations} \label{state_input_const}
\begin{align}
  \vect{x}(t) &\in\mathcal{X} = \mathcal{X}_{p} \cap \mathcal{X}_{v} \cap\mathcal{X}_{\xi},\\
  \vect{u}(t) &\in \mathcal{U} = \mathcal{U}_T\cap\mathcal{U}_{\omega},
\end{align}
\end{subequations}
where the state $\vect{x}(t)$ and input $\vect{u}(t)$ are defined in \eqref{state_input}, and $\mathcal{X}_{p},\mathcal{X}_{v},\mathcal{X}_{\xi},\mathcal{U}_T,\mathcal{U}_{\omega}$ are the constraint sets for the position, linear velocity, roll and pitch angles, total thrust, and angular velocity, respectively. For notation simplicity we first consider \eqref{state_input_const} over the entire time interval $t\in[\tau_0,\tau_v)$. The procedure will be to formulate each constraint in flat space and then use Proposition \ref{inclusion_prop} to derive finitely many sufficient conditions for \eqref{state_input_const}. In Section \ref{subavoidance}, we will discuss how to use Proposition \ref{inclusion_prop} to achieve obstacle avoidance by satisfying the constraints shown in \eqref{state_input_const} on sub-intervals instead. Apart from the continuous-time safety constraints, we will also consider waypoint constraints that have to be satisfied at certain discrete time instances.
% \begin{align}
%   \vect{x}(t_i) &\in\mathcal{W}_i,\quad i=1,...,n_{wp}.
% \end{align}
% Recall that second-order cone (SOC) constraints have the form  $\|A\vect{x}+\vect{b}\|_2\leq \vect{c}^T\vect{x}+d$ where  $A$ is a matrix and $\vect{b},\vect{c},d$ are vectors of appropriate dimensions \cite{boyd2004convex,alizadeh2003second}. Ellipsoidal $\vect{x}^TQ\vect{x} \leq r,\ Q\succeq0$ and polyhedral $A\vect{x}\leq b$ constraints are special cases of SOC constraints. 


%begins by setting up the optimization problem without safety constraints, and then illustrates




% \begin{remark}
% The optimization formulation shown in \eqref{opt_u} is based on B-spline parameterization of the flat trajectory, which is more convenient than the parameterization using piecewise polynomials as in  \cite{min_snap}. By property P1), the smoothness requirement from map \eqref{flat_map} is satisfied automatically by choosing $d\geq4$, without requiring additional smoothness constraints as in \cite{min_snap} or objective function reformulation as in \cite{polynomial}. 

%\XX{Any existing paper that use B-spline to formulate min-snap objective function as ours?}\VF{\cite{yu_minimum_2014} formulates min-jerk with B-splines. \cite{usenko_real-time_2017} considers a general min-(velocity, acc, jerk, snap) objective function with B-splines.} 

% a Quadratic Program (QP) and differs from the formulations like (10) in \cite{min_snap} in the use of B-spline functions as the basis to the flat output. Note that, by property P1), the smoothness requirement from map \eqref{flat_map} is satisfied automatically by choosing $d\geq4$ and we require no additional constraints like in \cite{min_snap} or objective function reformulation as in \cite{polynomial}.
%\XX{I am sure there are some paper use B-spline to formulate min-snap objective function. No need to emphasize the difference with  \cite{min_snap}. } 
%\end{remark}




%\subsection{Encoding Safety Constraints in Continuous Time} \label{State and Input Constraints}


% \XX{PLEASE revise this subsection carefully. A few examples: calling $\mathcal{X}_r$ the position constraint is misleading as it actually involves velocity and acceleration; the subscript $r$ is misleading as well, as $\mathcal{X}_r$ is the position constraint set when $r=1$, velocity constraint set when $r=2$, etc, then what does $\vect{x}(t)\in\mathcal{X}_r$ mean in  \eqref{state_input_const}?}\VF{I moved Waypoint constraints to the previous section. The next section focuses on constraining only the states and inputs in continuous-time. I separated position and velocity constraints.}

%\XX{Call $\vect{x}\in\mathcal{X}_r$ the position constraint is misleading as it does not only involve position.}
\subsubsection{Position Constraints $\vect{x}\in\mathcal{X}_p$}
The position constraint set $\mathcal{X}_p$ encodes limited flight volumes as follows:
\begin{equation}\label{const_pos}
    \mathcal{X}_{p} = \{\vect{x}\in\mathbb{R}^9\mid\vect{r} \in \mathcal{K}_p\},
\end{equation}
where $\mathcal{K}_p$ is a given SOC.
% as in \eqref{SOC}. 
%\XX{For example, when $K_p$ is chosen as...Note that polytopes are a special case of SOC...}\VF{In Section \ref{SOC_OPT} I already mention that Polytopes and Ellipsoids are special cases of $K_p$.} 
Proposition \ref{inclusion_prop} immediately provides sufficient conditions for $\vect{x}(t)\in\mathcal{X}_p,\ t\in[\tau_0,\tau_v)$:
\begin{equation}\label{cp_pos}
    \vect{P}_{j} \in \mathcal{K}_p, \quad j=0,\dots,N.
\end{equation}
Note that we restrict ourselves to SOC sets to preserve the overall SOCP structure that will arise from the rest of the constraints. %Note also that polytopes are a special case of SOC.

\subsubsection{Linear Velocity Constraints $\vect{x}\in\mathcal{X}_v$} Limiting flight velocities can  reduce the risk of injury or equipment damage. Consider the linear velocity constraint set defined as:
\begin{equation}\label{const_vel}
    \mathcal{X}_{v} = \{\vect{x}\in\mathbb{R}^{9}:\|\dot{\vect{r}}\|_2\leq \overline{v}\},
\end{equation} 
% \VF{I think having $\vect{r}(t)$ the $t$ time dependency in the set definition is incorrect? 
% \begin{equation} 
% % \label{const_vel}
%     \mathcal{X}_{v} = \{\vect{x}\in\mathbb{R}^{9}\mid\|\dot{\vect{r}}\|_2\leq \overline{v}\},
% \end{equation}
% }
where $\overline{v}\geq0$ denotes the allowable bound of the maximal speed. Using Proposition 1, the constraint $\vect{x}(t)\in\mathcal{X}_v, \ t\in[\tau_0,\tau_v)$ can be formulated as a constraint on the $1$-st order VCPs: 
$$
\|\vect{P}_j^{(1)}\|_2 \leq \overline{v},\ j=1,\ldots,N,
$$ 
which can be expressed equivalently as a SOC constraint on the control points:
\begin{equation} \label{cp_vel}
    \bigg\|\sum_{i=0}^N\vect{P}_ib_{1,i+1,j+1}\bigg\|_2\leq \overline{v}, \quad j=1,\ldots,N,
    % |\vect{P}_j^{(1)}| \leq \overline{v},\quad j=1,\ldots,N.
\end{equation}
where $b_{r,i,j}$ is the $(i,j)$-th entry of matrix $B_r$.

% [copied] Note that the $j$-th $r$-th order VCP is a linear combination of the decision variables:
% \begin{equation}
%     tj)$-th entry of matrix $B_r$.

% \XX{Explain how \eqref{cp_vel} is a SOC constraint on $\vect{P}$. Similarly for all the following constraints for VCPs.}

\subsubsection{Angular Constraints $\vect{x}\in\mathcal{X}_{\xi}$}  
Providing bounds for the roll and pitch angles can enhance system stability and improve the accuracy of the dynamics model \eqref{dyn_model} by limiting aerodynamic effects. This is especially relevant when implementing controllers that leverage linearizations (e.g., LQR or linear MPC) or small-angle approximations. Consider the angular constraints set defined as follows:
\begin{equation} \label{const_ang}
    \mathcal{X}_{\xi}=\{\vect{x}\in\mathbb{R}^9:|\phi|,|\theta| \leq \epsilon\},
\end{equation}
where $\epsilon< \pi/2$ is a given bound for the roll and pitch angles. We begin by exposing the geometric meaning of a result from literature on quadcopter differential flatness. % \cite{effective_ang}: 
\begin{lemma}\label{epsilon_cone_lemma} 
%\VF{Changed this a bit to be in SOC format because it comes after SOC section}  
Given a positive scalar $0<\epsilon<\pi/2$, and the following SOC:
\begin{align}\label{ang_bounds_cone}
    \mathcal{K}_{\epsilon} &\triangleq \{\vect{p}\in\mathbb{R}^3:\|A_{\epsilon}\vect{p}\|_2\leq p_3+g\},\\
    A_{\epsilon} &=|\cot\epsilon|\mat{\vect{x}_W & \vect{y}_W & \vect{0}_{3\times 1}},\nonumber
\end{align}
if the acceleration signal $\ddot{\vect{r}}(t)\in\mathcal{K}_{\epsilon}$,
% satisfies
% \begin{equation}\label{ang_bounds_cone}
%     \ddot{x}^2(t)\cos^2\epsilon+\ddot{y}^2(t)\cos^2\epsilon-\sin^2\epsilon(\ddot{z}(t)+g)^2\leq 0,\;\forall t,
% \end{equation}
then $|\phi(t)|,|\theta(t)|\leq \epsilon$ regardless of the yaw angle $\psi(t)$.
\end{lemma}
\begin{proof} 
Let $\ddot{\vect{r}}(t)\in\mathcal{K}_{\epsilon}$ and expand the SOC \eqref{ang_bounds_cone} to arrive at:
\begin{equation}
    \ddot{x}^2(t)\cot^{2}\epsilon + \ddot{y}^2(t)\cot^2\epsilon \leq (\ddot{z}(t)+g)^2.
\end{equation}
% which represents inclusion in a conic surface in $\ddot{\vect{r}}$ coordinates with apex at $(0,0,-g)$ opening along positive $\ddot{z}$.
% \begin{equation}
%     \ddot{x}^2(t)\cos^2\epsilon+\ddot{y}^2(t)\cos^2\epsilon-\sin^2\epsilon(\ddot{z}(t)+g)^2\leq 0.
% \end{equation}
Let $k_1(t)\triangleq \ddot{x}(t)$, $k_2(t)\triangleq \ddot{y}(t)$ and $k_3(t) \triangleq \ddot{z}(t)+g$. Dividing the inequality above by $\cot^2{\epsilon}\neq 0$ yields:
\begin{equation*}
    \frac{k_1^2(t)+k_2^2(t)}{k_3^2(t)}\leq \tan^2\epsilon.
\end{equation*}
Then, the conclusion follows from  Proposition 2 in \cite{effective_ang}.
\end{proof}
% \XX{Change the following into the form of a lemma and a proof. Also, in each case below, show in the proof which type of SOC the constraint derived is. }
Note that \eqref{ang_bounds_cone} describes the interior of a conic surface in $\ddot{\vect{r}}$ coordinates with apex at $(0,0,-g)$. We have the following result for angular constraints satisfaction.
\begin{lemma}
Consider B-spline curve $\vect{r}(t)$ defined as in \eqref{b-spline_curve}. If
\begin{align} \label{cp_ang}
    \bigg\|A_{\epsilon}\sum_{i=0}^N\vect{P}_ib_{2,i+1,j+1}\bigg\|_2 \leq g + \sum_{i=0}^N\vect{z}_W^T\vect{P}_ib_{2,i+1,j+1}
\end{align}
holds for $j=2,\ldots,N$ with $A_{\epsilon}$ given as in \eqref{ang_bounds_cone}, then $\vect{x}(t)\in\mathcal{X}_{\xi}, \forall  t\in[\tau_0,\tau_v)$.
\end{lemma}
% \begin{lemma}
% Sufficient conditions for $\vect{x}(t)\in\mathcal{X}_{\xi}, \ t\in[\tau_0,\tau_v)$ are:
% \begin{align} \label{cp_ang}
%     \bigg\|A_{\epsilon}\sum_{i=0}^N\vect{P}_ib_{2,i+1,j+1}\bigg\|_2 \leq g + \sum_{i=0}^N\vect{z}_W^T\vect{P}_ib_{2,i+1,j+1},\\
%   \mbox{with}\; A_{\epsilon}=\sqrt{\frac{1-\sin^2\epsilon}{\sin^2\epsilon}}\mat{1 & 0 & 0\\ 0 & 1 & 0 \\ 0 & 0 & 0}, \; j=2,\dots,N.\nonumber
% \end{align}
% \end{lemma}
% \VF{If I remove the sums, the formulation would be:
% \begin{align}
%     \bigg\|A_{\epsilon}P\vect{B}_{2,j}\bigg\|_2 \leq g + P\vect{B}_{2,j}, \quad j=2,\dots,N,\\
%     A_{\epsilon}=\sqrt{\frac{1-\sin^2\epsilon}{\sin^2\epsilon}}\mat{1 & 0 & 0\\ 0 & 1 & 0 \\ 0 & 0 & 0},\nonumber
% \end{align} Where $\vect{B}_{r,j}$ is the $j$-th column of $B_r$. However $P$ is not the decision variable exactly, but instead its columns $\vect{P}_i$.}

\begin{proof} Applying the definition of $2$nd order VCPs given in \eqref{virtual_cp_eq} to condition \eqref{cp_ang} yields:
\begin{equation} \label{vcp_ang}
\|A_{\epsilon}\vect{P}_j^{(2)}\|_2 \leq \vect{z}_W^T\vect{P}_j^{(2)}+g, \quad j=2,\dots,N.
\end{equation}
By Proposition \ref{inclusion_prop}, we have:
\begin{equation}\label{lem3ineq2}
    \|A_{\epsilon}\vect{r}^{(2)}(t)\|_2\leq \vect{z}_W^T\vect{r}^{(2)}(t) + g, \quad t\in[\tau_0,\tau_v).
\end{equation}
Since condition \eqref{lem3ineq2} is equivalent to $\ddot{\vect{r}}(t)\in\mathcal{K}_{\epsilon}$ with $\mathcal{K}_{\epsilon}$ given as in \eqref{ang_bounds_cone}, it follows from Lemma \ref{epsilon_cone_lemma} that $|\phi(t)|,|\theta(t)|\leq \epsilon$, $\forall t\in[\tau_0,\tau_v)$, which means that $\vect{x}(t)\in\mathcal{X}_{\xi}, \forall  t\in[\tau_0,\tau_v)$.
% \begin{equation}
%     \ddot{x}^2(t)\cos^2\epsilon+\ddot{y}^2(t)\cos^2\epsilon-\sin^2\epsilon(\ddot{z}(t)+g)^2\leq 0.
% \end{equation}
% Finally, we conclude from Lemma \ref{epsilon_cone_lemma} that the roll and pitch are bounded by this cone independently of the yaw angle, i.e., $|\phi(t)|,|\theta(t)|\leq \epsilon, \ t\in[\tau_0,\tau_v)$.
%\XX{It is not a direct result of Proposition 2 in \cite{effective_ang}. Give the detailed proof here to explain how the conclusion $|\phi(t)|,|\theta(t)|\leq \epsilon, \ t\in[\tau_0,\tau_v)$ can be derived.}\VF{I added a Lemma about this in the preliminaries because I use this cone bound also in a later section when I prove input saturation safeguard.}
% Finally, we recover from Proposition 2 in \cite{effective_ang} that the roll and pitch are bounded by this cone independently of the yaw angle to conclude $|\phi(t)|,|\theta(t)|\leq \epsilon, \ t\in[\tau_0,\tau_v)$.
\end{proof} 

% \begin{proof} \VF{Old backwards reasoning}
% In \cite{effective_ang}, it was shown that the roll and pitch angles are bounded independently of the yaw angle leveraging the differential flatness map. By this result, the constraints $|\phi|,|\theta| \leq \epsilon$ reduce to verifying:
% \begin{equation}\label{ang_bounds_cone}
%     (1-\sin^2\epsilon)\ddot{x}^2+(1-\sin^2\epsilon)\ddot{y}^2-\sin^2\epsilon(\ddot{z}+g)^2\leq 0.
% \end{equation}
% Note that \eqref{ang_bounds_cone} describes the region inside a conic surface in $\ddot{\vect{r}}$ coordinates with axis aligned with the $\ddot{z}$ coordinate and apex at $\mat{0&0&-g}^T$. By Proposition \ref{inclusion_prop},  \eqref{ang_bounds_cone} holds by enforcing the $2$-nd order VCPs to lie within the following  cone:
% \begin{equation} \label{vcp_ang}
% \|A_{\epsilon}\vect{P}_j^{(2)}\|_2 \leq \vect{z}_W^T\vect{P}_j^{(2)}+g, \quad j=2,\dots,N,
% \end{equation}
% with $A_\epsilon$ defined as in the Lemma. Applying the definition of $\vect{P}_j^{(2)}$ results in \eqref{cp_ang}.
% \end{proof}


\subsubsection{Thrust Constraints $\vect{u}\in\mathcal{U}_{T}$} 
Upper-bounding the total thrust is needed to prevent actuator saturation, while lower-bounding the total thrust is common in aerospace systems, e.g., to facilitate  soft-landing  \cite{accikmecse2011lossless}. Consider the thrust constraints  set as follows:
\begin{equation}\label{const_T}
    \mathcal{U}_T = \{\vect{u}\in\mathbb{R}^4\mid\underline{T}\leq T \leq \overline{T}\},
\end{equation}
where $\underline{T},\overline{T}$ are given constants satisfying $0\leq \underline{T}\leq g\leq \overline{T}$.
\begin{lemma}\label{lemthrust} 
Consider B-spline curve $\vect{r}(t)$ defined as in \eqref{b-spline_curve}. If
\begin{subequations}\label{cp_thrust}
\begin{align}
    \bigg\|\sum_{i=0}^N\vect{P}_ib_{2,i+1,j+1} + g\vect{z}_W\bigg\|_2\leq \overline{T}, \quad j=2,\ldots,N, \label{cp_Tmax}\\
    \sum_{i=0}^N\vect{z}_W^T\vect{P}_ib_{2,i+1,j+1}\geq \underline{T}-g, \quad j= 2,\ldots,N, \label{cp_Tmin}
\end{align}
\end{subequations}
hold, then $\vect{u}(t) \in \mathcal{U}_T, \; \forall t\in[\tau_0,\tau_v)$.
\end{lemma}

% \begin{lemma}\label{lemthrust}
% Sufficient conditions for $\vect{u}(t) \in \mathcal{U}_T, \ t\in[\tau_0,\tau_v)$ are:
% \begin{subequations}\label{cp_thrust}
% \begin{align}
%     \bigg\|\sum_{i=0}^N\vect{P}_ib_{2,i+1,j+1} + g\vect{z}_W\bigg\|_2\leq \overline{T}, \quad j=2,\ldots,N, \label{cp_Tmax}\\
%     \sum_{i=0}^N\vect{z}_W^T\vect{P}_ib_{2,i+1,j_1}\geq \underline{T}-g, \quad j= 2,\ldots,N. \label{cp_Tmin}
% \end{align}
% \end{subequations}
% \end{lemma}

\begin{proof}
%\XX{No need to say ``Suppose  \eqref{cp_thrust} holds'' as it is the lemma condition.}\VF{I addressed this in future proofs.} 
Applying the definition of $2$nd order VCPs \eqref{virtual_cp_eq} to \eqref{cp_Tmax} yields:
\begin{equation} \label{vcp_Tmax}
    \|\vect{P}_j^{(2)}+g\vect{z}_W\|_2\leq \overline{T}, \quad j = 2,\dots,N.
\end{equation}
By Proposition \ref{inclusion_prop}, and expanding the norm shown in \eqref{vcp_Tmax}, we have:
\begin{equation} \label{T_sphere}
    \sqrt{\ddot{x}^2(t) + \ddot{y}^2(t)+\big(\ddot{z}(t)+g\big)^2}\leq \overline{T}, \quad t\in[\tau_0,\tau_v),
\end{equation}
which constrains the acceleration to lie within a sphere in $\ddot{\vect{r}}$ coordinates centered at $(0,0,-g)$. Note that by the flatness map $\eqref{flat_T_map}$, \eqref{T_sphere} is equivalent to $T(t)\leq \overline{T},\ t\in[\tau_0,\tau_v)$. We now apply the same initial steps to \eqref{cp_Tmin} to get
\begin{equation}\label{vcp_Tmin}
    \vect{P}_j^{(2)}\geq \underline{T}-g, \quad j=2,\ldots,N,
\end{equation}
and applying Proposition 1 to get
\begin{equation}
    \ddot{z}(t)\geq \underline{T}-g, \quad t\in[\tau_0,\tau_v),
\end{equation}
which lower-bounds the $z$-component of acceleration. Therefore,
\begin{equation}
    \underline{T}\leq \sqrt{\ddot{x}^2(t) + \ddot{y}^2(t)+\big(\ddot{z}(t)+g\big)^2}=T(t)
\end{equation}
holds over $t\in[\tau_0,\tau_v)$. This completes the proof.
\end{proof}

% \begin{proof}\VF{Old backwards reasoning} From the flatness map \eqref{flat_T_map}, 
% % $\Phi$ (see e.g., \cite{min_snap,zhou2014vector}), %\XX{where did we explain the explicit form of the flatness map?},
% the upper-bound part of \eqref{const_T} is equivalent to:
% \begin{equation} \label{T_sphere}
%     \ddot{x}^2 + \ddot{y}^2+\big(\ddot{z}+g\big)^2\leq \overline{T}^2,
% \end{equation}
% which describes a convex spherical region in $\ddot{\vect{r}}$ coordinates with the center shifted $g$ units in the negative $\ddot{z}$ direction. By Proposition \ref{inclusion_prop},  \eqref{T_sphere} can be written equivalently as:
% \begin{equation} \label{vcp_Tmax}
%     \|\vect{P}_j^{(2)}+g\vect{z}_W\|_2\leq \overline{T}, \quad j = 2,\dots,N.
% \end{equation}
% By \eqref{virtual_cp_eq}, condition \eqref{vcp_Tmax} is equivalent to \eqref{cp_Tmax}.
% Next, we consider the constraint $\underline{T}\leq \ddot{x}^2 + \ddot{y}^2+\big(\ddot{z}+g\big)^2$. A sufficient condition satisfying this inequality is $\ddot{z}(t)\geq \underline{T}-g$ for all $t$. By Proposition \ref{inclusion_prop}, this sufficient condition is guaranteed by a single half-space constraint on the $2$-nd order VCPs: 
% $$
% \vect{z}_W^T\vect{P}_{j}^{(2)}\geq \underline{T}-g, j = 2,\dots,N,$$ 
% which is equivalent to \eqref{cp_Tmin} by \eqref{virtual_cp_eq}. %This completes the proof.
% % \begin{equation}\label{vcp_Tmin}
% %     \vect{z}_W^T\vect{P}_{j}^{(2)}\geq \underline{T}-g, \quad j = 2,\dots,N.
% % \end{equation}
% \end{proof}

\begin{remark}
%Adding a lower bound to the inequality shown in \eqref{T_sphere} makes it nonconvex. 
When the inequality $\underline{T}\leq T(t)$ is expressed in terms of the acceleration $\ddot{\vect{r}}(t)$ by \eqref{flat_T_map}, it signifies exclusion from a spherical region. As a result, it is a nonconvex constraint in terms of the B-spline control points. 
% The constraint $\underline{T}\leq T \leq \overline{T}$ is nonconvex in terms of control points
%\XX{Is this what you meant?}\VF{I specified with a longer sentence.}
Such a constraint was converted to a mixed-integer constraint in \cite{Constrained_trajectory}, and its conservative relaxation  was also considered in \cite{mueller2013model}. The spherical region \eqref{T_sphere} is bigger than typical quadcopter operation requires: the lower half ($\ddot{z} < -g$) of sphere \eqref{T_sphere} is accessed only when the quadcopter is ``upside-down'' ($|\phi|\geq \pi/2$ or $|\theta| \geq \pi/2$) and this is only required in very aggressive flights. The proof of Lemma \ref{lemthrust} reveals that some  generally unused feasible region is  sacrificed by lower-bounding $\ddot{z}(t)$ which results in a lower-bound for $T(t)$. 
\end{remark}

\subsubsection{Angular Velocity Constraints $\vect{u}\in\mathcal{U}_{\omega}$} 
Angular velocities are the inputs to the quadcopter model considered and are usually subject to saturation. Limits in angular velocities are also desirable to ensure integrity of gyroscope data. Consider the angular velocity constraints set defined as follows:
\begin{equation}\label{const_dang}
    \mathcal{U}_{\omega} = \{\vect{u}\in\mathbb{R}^4:|p|,|q| \leq \overline{\omega}\},
\end{equation}
where $\overline{\omega}$ is a given bound for the angular velocities. Recall that $r(t)=0$ because of the assumption $\psi(t) = 0$. 
\begin{lemma}\label{lemangular}
Consider B-spline curve $\vect{r}(t)$ defined as in \eqref{b-spline_curve}. For any vector  $\vect{\zeta}=(\zeta_1,...,\zeta_{N-d+1})^T$ whose entries are positive constants, if the following conditions:
\begin{subequations} \label{cp_omega}
\begin{align}
    \sum_{i=0}^N\vect{z}_W^T\vect{P}_ib_{2,i+1,j+1} &\geq \zeta_{l-d+1}-g,\quad j = l-d+2,\ldots,l, \label{cp_omega_Tmin}\\
    \bigg\|\sum_{i=0}^N\vect{P}_ib_{3,i+1,j+1}\bigg\|_2 &\leq \overline{\omega}\ \zeta_{l-d+1},\quad j=l-d+3,\ldots,l,\label{cp_omega_jerk}
\end{align}
\end{subequations}
hold for $l=d,\ldots,N$, then $\vect{u}(t)\in\mathcal{U}_{\omega},\; \forall t\in[\tau_0,\tau_v)$.
\end{lemma}
% $\underline{\vect{T}}\in\mathbb{R}^{N-d+1}$ satisfying:
% \begin{equation} \label{cp_omega_Tvec}
%     0\leq \underline{\vect{T}} = \{\underline{T}_k\}, \quad k = 1,\dots,N-d+1,
% \end{equation}\end{lemma}
\begin{proof}
% Suppose \eqref{cp_omega} holds for some $l\in\{d,d+1,\ldots,N\}$ and consider vector $\vect{\zeta}\in\mathbb{R}^{N-d+1}$ with positive entries. We can then apply the definition of $2$-nd order VCPs to \eqref{cp_omega_Tmin} and of $3$-rd order VCPs to \eqref{cp_omega_jerk} to obtain, respectively:
Applying the definitions of $2$nd order VCPs to \eqref{cp_omega_Tmin} and $3$rd order VCPs to \eqref{cp_omega_jerk} yields:
\begin{subequations}
\begin{align}
    \vect{z}_W^T\vect{P}_{j}^{(2)}\geq \zeta_{l-d+1}-g,  \quad j = l-d+2,\dots,l,\label{vcp_omega_Tmin}\\
    \|\vect{P}^{(3)}_j\|_2\leq \zeta_{l-d+1}\overline{\omega},\quad j = l-d+3,\dots,l.\label{vcp_omega_jerk}
\end{align}
\end{subequations}
Applying now \eqref{inclusion_local} from Proposition \ref{inclusion_prop} allows us to establish:
\begin{equation}\label{dang_bounds_lax_local}
T(t) \geq \zeta_{l-d+1},\; \|\dddot{\vect{r}}(t)\|_2 \leq \overline{\omega}\ \zeta_{l-d+1}, \quad t\in[\tau_l,\tau_{l+1}),
\end{equation}
which lower-bounds the thrust and upper-bounds the magnitude of the jerk over the $l$-th knot sub-interval. Recalling that the entries of $\vect{\zeta}$ are positive, we can combine inequalities to write:
\begin{equation*}
    \frac{\|\dddot{\vect{r}}(t)\|_2}{T(t)}\leq\frac{\|\dddot{\vect{r}}(t)\|_2}{\zeta_{l-d+1}}\leq \overline{\omega},\quad t\in[\tau_l,\tau_{l+1}). 
\end{equation*}
Now examine the flat maps \eqref{flat_p_map} and \eqref{flat_q_map} and take the absolute value of both sides to arrive at:
\begin{equation*}
    |p| = |-\vect{y}_B\cdot\vect{h}_{\omega}|, \quad
    |q| = |\vect{x}_B\cdot\vect{h}_{\omega}|.
\end{equation*}
Since $\vect{x}_B$ and $\vect{y}_B$ are unit vectors, we can write $|p|,|q|\leq \|\vect{h}_{\omega}\|_2$ by the definition of the dot product. Furthermore, examining $\vect{h}_{\omega}$ we note that the numerator is a projection of the jerk to a plane perpendicular to the $\ddot{z}$-axis which only reduces its norm, and that the denominator is exactly the thrust $T$. In particular, the following holds:
\begin{equation*}
    |p(t)|,|q(t)|\leq \|\vect{h}_{\omega}(t)\|_2\leq \frac{\|\dddot{\vect{r}}(t)\|_2}{T(t)}\leq \overline{\omega}, \quad  t\in[\tau_l,\tau_{l+1}).
\end{equation*}
Finally, taking $l=d,\ldots,N$ and letting $\vect{\tau}$ be a clamped knot vector satisfying \eqref{clamped} allows to establish $|p(t)|,|q(t)|\leq \overline{\omega}$ over the whole interval $[\tau_0,\tau_v)$. This completes the proof.
%\XX{Explain the detailed steps how the following conclusion can be derived, instead of just mentioning \cite{raff}.} \VF{I added the derivation for bounded angular velocities.}
% We now recover from \cite{raff} (Section IV-C) that the angular velocities $p,q$ are upper-bounded by the ratio of the jerk's magnitude and the thrust regardless of the yaw angle:
% \begin{equation}\label{lem5ineq}
%     |p(t)|,|q(t)|\leq \frac{\|\dddot{\vect{r}}(t)\|_2}{T(t)}\leq \overline{\omega},\quad t\in[\tau_l,\tau_{l+1}).
% \end{equation}
% Finally, the clamped knot vector \eqref{clamped} allows to establish \eqref{lem5ineq} over interval $[\tau_0,\tau_v)$ and for $l=d\ldots,N$. This completes the proof.\XX{Check the rewording.}
% Finally, taking $l=d\ldots,N$ and recalling that we consider a clamped knot vector \eqref{clamped} allows to establish the previous result over interval $[\tau_0,\tau_v)$ instead, concluding the proof.
\end{proof}
% \begin{proof} \VF{old backwards}By Section IV (C) of \cite{raff}, the angular velocities $p,q$ are bounded in flat space via the following inequality:
% \begin{equation}
%     |p(t)|, |q(t)| \leq \frac{\|\dddot{\vect{r}}(t)\|_2}{\|\ddot{\vect{r}}(t)+g\vect{z}_W\|_2} = \frac{\|\dddot{\vect{r}}(t)\|_2}{T(t)}.%\leq \overline{\omega},
% \end{equation}
% It is clear that if
% \begin{equation}\label{dang_bounds_lax_global}
% T(t) \geq k,\quad  \|\dddot{\vect{r}}(t)\|_2 \leq k\overline{\omega},\quad t\in[\tau_0,\tau_{v})  
% \end{equation}
% holds for any positive scalar $k\in\mathbb{R}$, then $\vect{u}(t)\in\mathcal{U}_{\omega},\ t\in[\tau_0,\tau_v)$. We now leverage the locality properties of B-splines to formulate the segment-wise equivalent:
% \begin{equation}\label{dang_bounds_lax_local}
% T(t) \geq \zeta_{l-d+1},\quad \|\dddot{\vect{r}}(t)\|_2 \leq \zeta_{l-d+1}\overline{\omega}, \quad t\in[\tau_l,\tau_{l+1}),
% \end{equation}
% which will allow us to maximize the feasible region while maintaining convexity by introducing $\vect{\zeta}$ as an optimization variable in the next section.
% % \begin{equation} \label{dang_bounds_lax_local}
% %     T(t) \geq \underline{T}_{l-d+1},\ \|\dddot{\vect{r}}(t)\|_2 \leq \underline{T}_{l-d+1}\overline{\omega}, \ t\in[\tau_l,\tau_{l+1}),
% % \end{equation}
% %with segment-wise minimum thrust bounds $\underline{T}_k$ defined as in \eqref{cp_omega_Tvec}. 
% %Note that the second statement of Proposition \ref{inclusion_prop} provides a framework for guaranteeing constraints over a local time interval as in \eqref{dang_bounds_lax_local}. 
% % We first constrain each subset of $2$-nd order VCPs:
% By Proposition \ref{inclusion_prop}, if the following conditions on $2$-nd order VCPs hold:
% \begin{equation} \label{vcp_omega_Tmin}
%     \vect{z}_W^T\vect{P}_{j}^{(2)}\geq \zeta_{l-d+1}-g,  \quad j = l-d+2,\dots,l,
% \end{equation}
% then the first inequality in \eqref{dang_bounds_lax_local} is satisfied over time intervals $[\tau_{l},\tau_{l+1})$ with $l=d,\dots,N$. Similarly, by Proposition \ref{inclusion_prop}, if the following condition on $3$-rd order VCPs hold over the same $l$ range:
% \begin{equation}\label{vcp_omega_jerk}
%     \|\vect{P}^{(3)}_j\|_2\leq \zeta_{l-d+1}\overline{\omega},\quad j = l-d+3,\dots,l,
% \end{equation} 
% then the second inequality in \eqref{dang_bounds_lax_local}, which guarantees bounded jerk, is satisfied over time intervals $[\tau_{l},\tau_{l+1})$. Finally, applying the definition of VCPs \eqref{virtual_cp_eq} to \eqref{vcp_omega_Tmin} and \eqref{vcp_omega_jerk} results in \eqref{cp_omega}.
% \end{proof}
% \XX{check the proof. I changed $\underline{T}$ to $\zeta$ to avoid notation confusion.}
% \begin{remark}
% The value $\zeta_i$ in Lemma \ref{lemangular} can be considered as the segment-wise minimum thrust bounds.
% \end{remark}
\subsubsection{Waypoint Constraints}
Waypoint constraints are commonly used for shaping the trajectory.  %Following the relaxed formulation of \cite{lax_wp_paper}, we 
Assume that there are $n_{wp}$ waypoints $\vect{p}_1^{wp},\ldots,\vect{p}_{n_{wp}}^{wp}\in\mathbb{R}^3$ near which the trajectory must pass at certain discrete times $t_i, (i=1,\ldots,n_{wp})$. Specifically, the following constrains are enforced for the position:
\begin{equation}
    \|\vect{r}(t_i)-\vect{p}^{wp}_i\|_2\leq d_{wp}, \quad i=1,\dots,n_{wp}
    \label{const_wp},
\end{equation}
where $d_{wp}\geq 0$ is the desired radius indicating the closeness of the trajectory to the waypoints. Note that the waypoint constraints are exact when $d_{wp}=0$. Recalling that $\vect{r}(t)$ is considered a B-spline curve, constraint \eqref{const_wp} is a SOC constraint in terms of its control points:
\begin{equation}\label{cp_wp}
    \bigg\|\sum_{j=0}^{N}\vect{P}_j\lambda_{j}(t_i) - \vect{p}_{i}^{wp}\bigg\|_2\leq d_{wp}, \quad i = 1,\dots,n_{wp}.
\end{equation}
%The above constraint may be added to \eqref{opt_u} to modify the shape of the trajectory. Note that when $d_{wp}=0$, the problem remains a QP.

\subsection{Safe Trajectory Planning as SOCP}
%\subsection{Minimum-snap Trajectory Without Safety Constraints}

Suppose that the initial (resp. final) conditions of the trajectory are fixed and given as $\vect{p}^{0}_r$ (resp. $\vect{p}^{f}_r$) where $r = 0,\dots,n_0$ (resp. $r=0,\dots,n_f$). For example, the initial (resp. final) position of the trajectory is denoted as $\vect{p}_0^0$ (resp. $\vect{p}_0^f$) for $r=0$, and the initial (resp. final) velocity is denoted as $\vect{p}^{0}_1$ (resp. $\vect{p}^{f}_1$) for $r=1$. 
%In general, we may have initial (resp. final) conditions $\vect{p}^{0}_r$ (resp. $\vect{p}^{f}_r$) for $r = 0,\dots,n_r$. We can encode them with:
These initial/final constraints can be expressed as:
\begin{equation}\label{extremes}
    \sum_{i=0}^N\vect{P}_i\vect{b}_{r,i+1}^T\vect{\Lambda}_{d-r}(t_m) = \vect{p}^m_{r}, \quad r = 0,\dots,n_m,
\end{equation}
where $m\in\{0,f\}$ represents initial and final labels, respectively, and $0\leq n_0,n_f\leq d$. Note that $n_0$ may be different from $n_f$ and that $t_0 = \tau_0$ and $t_f = \tau_v$. 


%Minimum-snap trajectories are desirable for quadcopters because they naturally minimize the actuator effort and are easy to compute \cite{min_snap}. 
%and \XX{more reasons??}\VF{I don't know any other. Minimum jerk will minimize $\vect{u}$ but minimum snap will minimize $\vect{u}_{12} = \mat{T &M_x & M_y & M_z}$ the input to the 12-Dim model.} \cite{min_snap}. 
In this work we choose an objective function that minimizes $\int_{\tau_0}^{\tau_v}\|\vect{r}^{(4)}(t)\|^2_2 \mathrm{d}t$ where $\vect{r}^{(4)}(t)$ denotes the ``snap'' of the trajectory \cite{min_snap}. 
Taking $\vect{r}^{(r)}$ to be a B-spline curve constructed as in \eqref{b-spline_curve}, the (simple) objective function is given as 
%$J_s(\vect{P}_{0},\dots,\vect{P}_{N}) = \int_{t_0}^{t_f}\bigg\|\sum_{i=0}^N\vect{P}_i\vect{b}_{d,4,i+1}^T\vect{\Lambda}_{d-4}(t)\bigg\|_2^2 \mathrm{d}t,$ 
\begin{equation*}%\label{min_snap_obj}
    J_s(\vect{P}_{0},\dots,\vect{P}_{N}) = \int_{t_0}^{t_f}\bigg\|\sum_{i=0}^N\vect{P}_i\vect{b}_{d,4,i+1}^T\vect{\Lambda}_{d-4}(t)\bigg\|_2^2 \mathrm{d}t,
\end{equation*}
where the decision variables are the control points $\vect{P}_0,\ldots,\vect{P}_N$ of the B-spline curve $\vect{r}(t)$. %Other objective functions are also possible; e.g., the minimum path length as in \cite{Flat_trajectory_design}. 
By examining \eqref{dang_bounds_lax_local} we observe that it is desirable to make each entry $\zeta_k$ large to maximize the feasible region. Therefore, we  introduce $\vect{\zeta}$ into the objective function such that it is maximized. Specifically, the modified objective function is chosen to be:
\begin{equation} \label{obj_dang_bounds}
    J(\vect{P}_0,\dots,\vect{P}_{N},\vect{\zeta}) = J_s(\vect{P}_0,\dots,\vect{P}_{N})- \sum_{k=1}^{N-d+1}\zeta_{k}.
\end{equation}

%\XX{are $n_0,n_f$ different?} 
%\XX{Why we should put final position constraint here? Isn't that part of the waypoint constraint?}\VF{If we don't specify a final position for the trajectory, then we are just going to generate a trajectory to 0 to minimize the snap. This is sort of like the boundary value problem (initial and final position given).}

% Trajectory generation without safety constraints is now formulated as the following Quadratic Program (QP):
% \begin{align}    \label{opt_u}
%     \min_{\vect{P}_{0},\dots,\vect{P}_{N}} & \quad J_s(\vect{P}_0,\dots,\vect{P}_N) \\
%     \text{s.t.} &\quad  \eqref{extremes}.
%     \nonumber
% \end{align}

%Any of the constraints described in the previous section can be added to \eqref{opt_u}  


Summarizing results above, the safe trajectory planning problem can  now be expressed as the following  SOCP:
\begin{align} \label{FLAT-SOCP}
\tag{FLAT-SOCP}
    \min_{\vect{P}_{0},\dots,\vect{P}_{N},\vect{\zeta}} \quad &J(\vect{P}_0,\dots,\vect{P}_N,\vect{\zeta}) \\
    \text{s.t.} \quad & \eqref{cp_pos},\eqref{cp_vel},\eqref{cp_ang},\eqref{cp_thrust},\eqref{cp_omega},\eqref{cp_wp},\eqref{extremes}\;\text{hold},\nonumber
\end{align}
where ``FLAT'' means that the solutions parameterize the flat output trajectory $\vect{\sigma}(t) = (\vect{r}(t),0)$. % \XX{finish}\VF{Done I think.}
The trajectory in flat space is defined by control points generated from \eqref{FLAT-SOCP}, and the corresponding state/input trajectory satisfies all the safety constraints. This is summarized in the following theorem.
\begin{theorem}\label{thm1}
Suppose that \eqref{FLAT-SOCP} has a feasible solution $\vect{P}_0,\dots, \vect{P}_N$. Let $\vect{r}(t)$ be a B-spline curve defined by the control points $\vect{P}_0,\dots, \vect{P}_N$. Then, the quadcopter state $\vect{x}(t)$ and input $\vect{u}(t)$ trajectories, which are obtained from the flat map \eqref{flat_map} with $\vect{\sigma}(t) = (\vect{r}(t),0)$,
%(obtained from the flat output trajectory $\vect{\sigma}(t) = (\vect{r}(t),0)$, where $\vect{r}(t)$ is a B-spline curve defined by the control points $\vect{P}_0,\dots, \vect{P}_N$) 
satisfy the safety constraints \eqref{state_input_const} for all $t\in[\tau_0,\tau_v)$.%\XX{Check}\VF{Looks good. No need to use the solution $\vect{\zeta}$ because it is just an adaptive parameter for a constraint. And I assume no need to prove this because of all the previous Lemmas.}
\end{theorem}


Problem \eqref{FLAT-SOCP} can be relaxed to a QP if one replaces the convex regions described by the cone constraints in \eqref{FLAT-SOCP} by their polytopic inner approximations. For example, for the SOC constraint \eqref{cp_Tmax}, if we find a polytopic inner approximation as follows: 
$$
\{\ddot{\vect{r}}\in\mathbb{R}^3:A_T\ddot{\vect{r}}\leq \vect{b}_{T}\}\subset\{\ddot{\vect{r}}\in\mathbb{R}^3:\eqref{T_sphere}\;\mbox{holds}\},
$$
then we can replace the SOC constraint \eqref{cp_Tmax} with the linear constraint 
$$
A_T\vect{P}_j^{(2)}\leq \vect{b}_T,\quad  j=2,\dots,N,
$$ 
in \eqref{FLAT-SOCP}. Other replacements can be done in a similar way for all SOC constraints. 



\begin{remark}
The optimization formulation shown in \eqref{FLAT-SOCP} is based on B-spline parameterization of the flat trajectory, which is more convenient than the parameterization using piece-wise polynomials as in  \cite{min_snap}. By property P1), the smoothness requirement from the flatness map is satisfied automatically by choosing $d\geq4$, without requiring additional smoothness constraints as in \cite{min_snap} or objective function reformulation as in \cite{polynomial}. 
\end{remark}

\begin{remark}
The feasible region of \eqref{FLAT-SOCP} may be expanded at the expense of computational burden by increasing $N$ (number of control points), which must always respect $v = N+d+1$.
\end{remark}

\begin{remark}
In \cite{Flat_trajectory_design}, the problem of ensuring continuous-time safety for thrust and Euler angles is considered via solving a nonlinear optimization problem, for which only local minimum can be found.  In contrast, \eqref{FLAT-SOCP} is a convex program whose optimal solution can be found in polynomial time.
% Furthermore, a 6-dimension quadcopter model is considered in \cite{Flat_trajectory_design}, while a 9-dimension quadcopter model is considered in this paper. %is higher fidelity (9 states instead of 6) than that considered in . %\XX{Quotation marks are `` '', not `` ``}%\XX{Need rephrasing. ``In \cite{Flat_trajectory_design}, ?? problem is also considered. However, ....''} \VF{Rephrased.}
%\XX{Do we have a comparison with \cite{Flat_trajectory_design} and  other works along this line? If not, add it here.} \VF{I added some remarks about this.}
\end{remark}

\begin{remark} \label{Local_cst_rmk}
%The presented constraints were satisfied for all $t\in[\tau_0,\tau_v)$. 
Theorem \ref{thm1} ensures that the trajectory generated respects the safety constraints \eqref{state_input_const} for all $t\in[\tau_0,\tau_v)$. 
However, Proposition \ref{inclusion_prop} provides a framework for satisfying the constraints for $t\in[t_1,t_2]\subseteq[\tau_0,\tau_v)$ as long as we can find two knots $\tau_i, \tau_j$ such that $[t_1,t_2]\subseteq [\tau_i,\tau_j)$. This type of local constraint satisfaction may be of interest in applications where safety is needed only at certain times as well as in obstacle avoidance which will be shown in the next subsection.  %which will be discussed in Section \ref{subavoidance}.
\end{remark}
% \begin{remark}
% The feasible region of \eqref{FLAT-SOCP} may be expanded at the expense of computational burden by increasing $N$ (number of control points). Note that the parameter $N$ must always respect $v = N+d+1$.
% \end{remark}

\begin{remark}
Constraint \eqref{cp_omega_Tmin} is compatible with \eqref{cp_Tmax} and \eqref{cp_Tmin}. When $\underline{T}$ is provided, it follows that $\underline{T} \leq \zeta_k \leq \overline{T}, k = 1,\dots,N-d+1$. In fact, the present formulation may lower-bound $T(t)$ more tightly than required to expand the feasible region for \eqref{dang_bounds_lax_local}, but never so much that it violates the upper bound. 
%\VF{Comment out the rest of this remark if space is needed}
Additionally, while constraint \eqref{cp_omega} increases optimality by expanding the feasible region, it comes at the expense of computational cost. A possible compromise is by replacing the vector $\vect{\zeta}$ with a scalar $\zeta_\omega$ in the objective function \eqref{obj_dang_bounds} and replacing constraint \eqref{cp_omega} with the following global constraint: 
\begin{align*}
    \sum_{i=0}^N\vect{z}_W^T\vect{P}_ib_{2,i+1,j+1}&\geq \zeta_\omega-g, \quad &j= 2,\dots,N,\\
    \bigg\|\sum_{i=0}^N\vect{P}_ib_{3,i+1,j+1}\bigg\|_2 &\leq \overline{\omega}\ \zeta_{\omega}, \quad &j=3,\dots,N.
\end{align*}

\end{remark}
% \begin{remark}
%   The SOCP problem \eqref{FLAT-SOCP} can be relaxed to a quadratic programming by finding polytopic approximations to the convex regions described by the cone constraints. For example, finding polytope $\{\ddot{\vect{r}}\in\mathbb{R}^3:A_T\ddot{\vect{r}}\leq \vect{b}_{T}\}\subset\{\ddot{\vect{r}}\in\mathbb{R}^3:\eqref{T_sphere}\}$ and replacing \eqref{cp_Tmax} with:
%   \begin{equation*}
%     A_T\vect{P}_j^{(2)}\leq \vect{b}_T, \quad j=2,\dots,N.
%   \end{equation*}
% \end{remark}

% \XX{A subsection about how to incorporate obstacle avoidance into our method above, e.g., \cite{stoical2016obstacle}. The problem will be a mixed-integer programming, or something else?}\VF{I Added the two formulations for obs-avoidance.}

\subsection{Obstacle Avoidance with Continuous-time Guarantees via SOCP}\label{subavoidance}


% We present two approaches to guarantee continuous-time obstacle avoidance in \eqref{FLAT-SOCP}. The first method (direct) results in Mixed-Integer constraints and is not suitable for online replanning. The second method (indirect) allows the problem to remain convex by placing the burden of parsing the configuration space on a high-level path finding algorithm.

%\subsubsection{Direct Obstacle Avoidance}

The safe flight space is usually non-convex which makes the obstacle avoidance problem difficult \cite{tang2021real}. Obstacle avoidance with continuous-time safety guarantees was studied in \cite{stoical2016obstacle} where each obstacle is expressed as a polytope and the trajectory planning problem is formulated as a mixed-integer QP, which is an NP-hard problem and renders the online trajectory re-planning computationally infeasible. In this subsection, we show that, by leveraging the locality property of B-splines, \eqref{FLAT-SOCP} can be modified to achieve obstacle avoidance with a rigorous continuous-time safety guarantee. 


% Obstacle avoidance constraints are typically non-convex and methods involving mixed-integer programming have been proposed in \cite{stoical2016obstacle}. We formulate one such method here.

% Let there be $n_{obs}$ obstacles defined as $O_l$ with $l=1,\ldots,n_{obs}$. Each obstacle is the open polytope:
% \begin{equation}
%     O_l \triangleq \{\vect{r} \in\mathbb{R}^3:A^l\vect{r}< \vect{b}^l, \ \vect{b}^{l}\in\mathbb{R}^{c_{l}}\},
% \end{equation}
% with $c_l\in\mathbb{N}, l=1,\dots,n_{obs}$ hyperplane constraints. Then, sufficient conditions to guarantee $\vect{r}(t) \notin O_l$ are:
% \begin{align} \label{cp_obs}
%     {\vect{a}_k^l}^T\vect{P}_j&\geq b_k^l-\alpha_{i,k}^lM, \quad j  = i-d,\ldots,i, \quad k=1,\ldots,c_{l},\nonumber\\
%     \sum_{k=1}^{c_{l}}\alpha_{i,k}^l &\leq c_l-1, \quad \alpha_{i,k}^l \in \{0,1\}, \quad i = d,\ldots,N,
% \end{align}
% where ${\vect{a}_k^l}^T$ is the $k$-th row of $A^l$.  Constraint \eqref{cp_obs} may be added to the optimization problem \eqref{FLAT-SOCP} at the expense of lifting the problem to a Mixed-Integer SOCP.

%\subsubsection{Indirect Obstacle Avoidance} 
% This method is also common in the literature \cite{gao2018online} and relies on the convexification of the generally nonconvex configuration space. Let the obstacle space be the compact set $\mathcal{O}$ such that the available flight space is the open set $\mathcal{F} \triangleq \mathbb{R}^3\setminus\mathcal{O}$. Then, we have the following assumption:
% \begin{assumption} \label{path_exists_assumption}
% There exist a continuous function  $\vect{f}(t)\in \mathcal{F}, \forall t\in[t_0,t_f)$ such that $\vect{f}(t_0) = \vect{p}_0^0$ and $\vect{f}(t_f) = \vect{p}_0^f$.
% \end{assumption}
% The above assumption is quite mild and only states that the initial and final positions are not isolated from each other in the available flight space. We now seek to find such a path $\vect{f}(t)$. Consider the following lemma:
% \begin{lemma} \label{convex_sets_lemma}
% If Assumption \ref{path_exists_assumption} holds, then there exist convex sets $S_l\subset \mathcal{F}, \ l=0,\ldots,N-d$, for sufficiently large $N$, such that:
% \begin{align}
%     \vect{p}_0^0&\in S_0,\quad \vect{p}_0^f\in S_{N-d},\quad\\
%     S_i&\cap S_{i+1}\neq \emptyset, \quad i = 0,\ldots,N-d-1.\nonumber
% \end{align}
% \end{lemma}
% \begin{proof} Suppose Assumption \ref{path_exists_assumption} holds. Then, the ball $\mathcal{B}(t) = \{\vect{r}\in\mathbb{R}^3:\|\vect{r}-\vect{f}(t)\|_2\leq \delta\}$ exists $\forall t\in [t_0,t_f)$ for sufficiently small $\delta$. Then, for some sufficiently large $N$, we can construct finitely-many balls $\mathcal{B}(t_l),\ l =0,\ldots,N-d$ such that:
% \begin{align}
%     \vect{f}(t)&\in\bigcup_{l=0}^{N-d}\mathcal{B}(t_l), \quad \vect{f}(t_0) \in \mathcal{B}(t_0),\quad \vect{f}(t_f) \in \mathcal{B}(t_{N-d}),\nonumber\\ 
%     \mathcal{B}(t_i)&\cap\mathcal{B}(t_{i+1})\neq \emptyset,\quad i=0,\ldots,N-d-1.
% \end{align} 
% Taking these balls as the convex sets from Lemma 1 concludes the proof.
% \end{proof}

% Let the obstacle space be a compact set $\mathcal{O}$ such that the available flight space is the open set $\mathcal{F} \triangleq \mathbb{R}^3\setminus\mathcal{O}$. Although $\mathcal{F}$ is generally nonconvex, mild assumptions about the distribution of the obstacles allow  us to determine the existance of convex subsets that, together, form a ``safe corridor'' in $\mathcal{F}$ for the trajectory (see, for example, \cite{gao2018online}). 

% \XX{Add something related to ``safe corridor'' such as in \cite{gao2018online}.}
%We now propose sufficient conditions to add to \eqref{FLAT-SOCP} so that $\vect{r}(t) \in \mathcal{F} $:

Suppose that the safe flight space $\mathcal{F} \in \mathbb{R}^3$ is a set defined as $\mathcal{F} \triangleq \mathbb{R}^3\setminus\left(\mathcal{O}_1\cup...\cup\mathcal{O}_{n_o}\right)$ where $\mathcal{O}_i\in \mathbb{R}^3$ denotes the overapproximation of the $i$-th obstacle. Although $\mathcal{F}$ is generally nonconvex, mild assumptions about the distribution of the obstacles allow  us to determine the existance of convex subsets that, together, form a ``safe corridor'' in $\mathcal{F}$ for the trajectory (see, for example, \cite{gao2018online}). The union of convex subsets can approximate $\mathcal{F}$ arbitrarily well. In practice, these convex subsets are easily found with sensors such as stereo cameras and lidars \cite{hornung13auro}. The time-optimality of trajectories generated using ``safe corridors'' remains an active research topic  \cite{tang2019time,sun2020fast}. 

% \XX{discuss a bit more about ``safe corridor'', like its property and construction, if possible.}

% \XX{The proposition can be more precise?}
\begin{proposition} \label{convex_sets_prop}
Suppose that there exist $n_s$ overlapping convex sets $\{\mathcal{S}_1,\ldots,\mathcal{S}_{n_s}\}$ inside the safe flight space, i.e., $\mathcal{S}_l\subset \mathcal{F},\ l=1,\ldots,n_s,$ satisfying $\mathcal{S}_i\cap \mathcal{S}_{i+1}\neq \emptyset, \ i = 1,\ldots,n_s-1$ and consider B-spline curve $\vect{r}(t)$ defined as in \eqref{b-spline_curve} with $N=n_s+d-1$. If the following conditions hold:
\begin{equation}\label{eqobstc}
    \vect{P}_{j-1} \in \mathcal{S}_{l},\quad j = l,\ldots,l+d, \quad l = 1,\ldots,n_s,
\end{equation}
then  $\vect{r}(t)\in\mathcal{F}$  for all $t\in[\tau_0,\tau_v)$ and $\vect{r}(t)$ is at least $C^{d-1},\  \forall t\in[\tau_0,\tau_v)$.
\end{proposition}
\begin{proof} Applying Proposition \ref{inclusion_prop} to the inclusion constraint of the $l$-th  group of control points results in $\vect{r}(t)\in \mathcal{S}_l, \ t\in[\tau_{l+d-1},\tau_{l+d})$, which means that the $l$-th segment of the B-spline curve $\vect{r}(t)$ is contained within $\mathcal{S}_l$. Taking now $l=1,\ldots,n_s$ and recalling that we consider a clamped knot vector \eqref{clamped} results in $\vect{r}(t)\in\mathcal{F},\ t\in[\tau_0,\tau_v)$. Further, note that $\mathcal{S}_l\cap \mathcal{S}_{l+1} \neq \emptyset$ allows $\vect{r}(t)$ to be continuous despite the segment-wise constraints. 
\end{proof}
% \begin{proposition} \label{convex_sets_prop}
% Suppose that there exist  overlapping convex sets $\{S_0,S_1,...,S_{N-d}\}$ inside the safe flight space, i.e., $S_l\subset \mathcal{F},\ l=0,\ldots,N-d,$ satisfying $S_i\cap S_{i+1}\neq \emptyset, \ i = 0,\ldots,N-d-1$. If the following condition hold:
% \begin{align}\label{eqobstc}
%     \vect{P}_j \in S_{l},\quad j = l,\ldots,l+d, \quad l = 0,\ldots,N-d.
% \end{align}
% then  $\vect{r}(t)\in\mathcal{F}$ for all $t\geq t_0$.
% \end{proposition}
% \XX{Can we change the number of convex sets to an arbitrary number $n$? Letting it to be $N-d$ is weird. Also, it is not common to write ``sufficient conditions for XX is XX'' in a theorem; using ``if...then'' is more common. Why use $\vect{s}(t)$ not $\vect{r}(t)$?}
% Sufficient conditions for $\vect{s}(t)\in\mathcal{F}$ are $\vect{s}(t)\in S_i, \ t\in[\tau_{i+d},\tau_{i+d+1}),\ i = 0,\ldots,N-d,$
% which states that the $i$-th segment of the B-spline curve should be contained within $S_i$. Note that $S_i\cap S_{i+1} \neq \emptyset$ allows $\vect{s}(t)$ to be continuous despite the segment-wise constraints. Then, applying Proposition \ref{inclusion_prop} for each local constraint results in the conditions of Proposition \ref{convex_sets_prop}.

% \XX{Your proof is always written in an ``reverse order'', which is not good. The proof is not a place to show your thinking process, but to let the reader understand more easily with a good flow. Don't refer to Proposition 2 in the proof of Prop. 2. You even used $s(t)$ in the proof. Doublecheck the proof above to make sure it is correct!}
%

Convex sets $\mathcal{S}_1,\ldots,\mathcal{S}_{n_s}$ in Proposition \ref{convex_sets_prop} play a significant role in the properties of the resulting B-spline curve (such as the magnitude of its derivatives). Finding these sets is considered the task of a high-level planner, and is out of the scope of this work. The reader is referred to \cite{tang2019time,sun2020fast} for more detail. Note that the constraints shown in  \eqref{eqobstc}  can be readily added to \eqref{FLAT-SOCP} without compromising the problem structure if the convex sets $\{\mathcal{S}_l\}$ are restricted to SOC form, as will be demonstrated with experiments in Section \ref{sec:exper}.

% \begin{remark}
% In Proposition \ref{convex_sets_prop}, the sequence of sets $S_l, \ l=0,\ldots,N-d$ plays a significant role in the properties of the resulting B-spline curve (such as the magnitude of its derivatives). The optimal finding of such sets is out of the scope of this work and the reader is referred to \cite{tang2019time,sun2020fast} for more information. 
% \end{remark}

\begin{remark}
In Section \ref{sec:track} we will establish safety guarantees for bounded reference position trajectory tracking in the infinity-norm sense. This bound at the tracking level can be incorporated when finding the sets $\{\mathcal{S}_l\}$ so that obstacle avoidance is guaranteed also during tracking. Alternatively, one can appropriately enlarge the size of obstacles $\{\mathcal{O}_i\}$.
\end{remark}

\section{Trajectory Tracking with Continuous-Time Safety Guarantee} \label{sec:track}
In this section, we propose a trajectory tracking method for quadcopters with position safety guarantees in continuous-time. The idea is to use CBFs as a safety filter to guarantee boundedness between the real trajectory and the nominal safe trajectory $\vect{r}^{\text{ref}}(t)$ generated from \eqref{FLAT-SOCP}. The safe tracking controller is obtained by solving a convex QP online, whose feasibility will be ensured by the control-sharing property of multiple CBFs. %We make use of a simplified dynamics model and provide feasibility guarantees of the safety filter.
\subsection{6-Dim Quadcopter Model \& Feedback Linearization}
We assume now that the angular velocity dynamics are regulated by some high-bandwidth controller and consider the reduced state and input vectors:%XX{replace $\vect{z}$ with another symbol.}
\begin{subequations}
\begin{align}
    \vect{z}&= \mat{x & y & z & \dot{x} & \dot{y} & \dot{z}}^T,\\
    \vect{v} &= \mat{T & \phi & \theta &  \psi}^T,\label{reduced_input}
\end{align}
\end{subequations}
subject to \eqref{dyn_model_acc}. Then, we can consider a set of virtual inputs $\vect{\mu} =(\mu_1,\mu_2,\mu_3)^T\triangleq (\ddot{x},\ddot{y},\ddot{z})^T$ found as in \eqref{dyn_model_acc}:
\begin{equation} \label{def_Psi}
    \vect{\mu} = \Psi(\vect{v}) \triangleq T\vect{z}_B-g\vect{z}_W.
\end{equation}
%\XX{Is $\Psi(\vect{v})$ in the equation above defined for the first time in this paper?}\VF{Map $\Psi$ defines the transformation of the reduced input $\vect{v}$ into virtual inputs $\vect{\mu}$. But it's not really new, $\Psi$ is exactly \eqref{dyn_model_acc} and this transformation is also used in \cite{FMPC,cbf_pos} as cited at the end of the paragraph.} \XX{Then mention where $\Psi$ is defined.} \VF{Defined in \eqref{def_Psi}} 
Note that the map is invertible if we recover the yaw angle such that  $\vect{v} = \Psi^{-1}(\vect{\mu},\psi)$ where  $\Psi^{-1}$ is given as: %\XX{Are the expressions correct?} \VF{Yes, they are just more explicit forms of parts of the map \eqref{flat_map}. Notice that \eqref{flat_map} also depends on $\psi$ because it is a flat output.}
\begin{subequations}
\begin{align}
    T &= \sqrt{\mu_1^2 + \mu_2^2 + (\mu_3+g)^2} = \frac{\mu_3+g}{\cos\phi \cos\theta},\label{inverpsi1}\\
    \phi &= \arctan\Big(\frac{\mu_1\sin\psi-\mu_2\cos\psi}{(\mu_3+g)\cos\theta}\Big),\label{inverpsi2}\\
    \theta&= \arctan\Big(\frac{\mu_1 \cos\psi + \mu_2 \sin\psi}{\mu_3+g}\Big).\label{inverpsi3}
\end{align}
\end{subequations}
% \XX{What does this mean? Why $\psi$ appears in the function domain of $\Psi^{-1}$? no $\psi$ appears in $\vect{\mu} = \Psi(\vect{v})$} \VF{Notice that $\vect{v}\in\mathbb{R}^4$ but $\vect{\mu}\in\mathbb{R}^3$. Then, the transformation $\vect{\mu} =\Psi(\vect{v})$ (where $\Psi$ is unrelated to $\psi$ loses some information ($\psi$) that must be recovered (fourth element of $\vect{v}$ i.e. $v_4$) to invert the map as $\Psi^{-1}$. The process is: Nominal controller gets $\vect{v}$. We then store $\psi=v_4$ and transform to virtual inputs $\vect{\mu} = \Psi(\vect{v})$. We can now solve \eqref{CBF_QP} and use the stored $\psi$ to get the safe control $\vect{v}_s = \Psi^{-1}(\vect{\mu}^*,\psi)$}\XX{It does not make sense. What makes sense to me is $\vect{v} = (\Psi^{-1}(\vect{\mu}),\psi)$ where $\psi$ is defined in (??)}\VF{Here is the expression for $\Psi^{-1}$: \begin{align}
%     T &= \frac{\mu_3+g}{c\phi c\theta},\\
%     \phi &= \arctan(\frac{\mu_1s\psi-\mu_2c\psi}{(\mu_3+g)c\theta}),\\
%     \theta&= \arctan(\frac{\mu_1 c\psi + \mu_2 s\psi}{\mu_3+g}).
% \end{align}
%  You can see that we cannot compute $\phi$ and $\theta$ without having $\psi$. The value of $\psi$ is the one provided by the nominal controller $\pi(\vect{z})$. If you want, instead of $\Psi^{-1}$ since it is not as simple as just the inverse, we can call it a different function altogether: $\Theta(\vect{\mu},\pi(\vect{z}))$ or $\Theta(\vect{\mu},\psi)$}. 
With the virtual inputs, the quadcopter dynamics becomes a double-integrator \cite{FMPC,cbf_pos}: 
\begin{equation} \label{lin_sys}
    \dot{\vect{z}}(t) = f(\vect{z}(t)) + g(\vect{z}(t)) \vect{\mu}(t) =  A\vect{z}(t) + B\vect{\mu}(t)
\end{equation}
where 
%$A = \mat{\vect{0}_3 & I_3\\ \vect{0}_3 & \vect{0}_3}, \quad B=\mat{\vect{0}_3\\ I_3}.$ 
\begin{equation*}
    A = \mat{\vect{0}_3 & I_3\\ \vect{0}_3 & \vect{0}_3}, \quad B=\mat{\vect{0}_3\\ I_3}.
\end{equation*}
Assume that a reference state $\vect{x}^{\text{ref}}(t)$ and a reference input $\vect{u}^{\text{ref}}(t)$ are generated as in Sec. \ref{sec:traj} and note that we can immediately extract $\vect{z}^{\text{ref}}(t)$ and $\vect{v}^{\text{ref}}(t)$ from them. The dynamic feasibility of the reference trajectory will be necessary to establish the forthcoming results. Assume also a \emph{nominal controller} $\vect{v} = \pi(\vect{z})$ is given, which is potentially unsafe. In the following subsection, a CBF-QP-based tracking controller will be designed such that the real trajectory lies within a prescribed tube of the nominal trajectory while the tracking control is close to the nominal controller as much as possible.  
%We will use the CBF-QP-based method to filter the nominal controller such that the real trajectory is in a prescribed tube of the nominal trajectory while the tracking control is close to the nomial controller as much as possible.  % but will be modified to be safe in the following section.
% \XX{Where is $\vect{u}^{\text{ref}}(t)$ used? how is $\vect{v} $ obtained and where is it used? }

\subsection{Safe Trajectory Tracking via CBF-QP Controller}
Consider the following time-varying safe set:
\begin{equation}\label{eqS}
    \mathcal{S}_h(t) = \{\vect{z}\in\mathbb{R}^6: \|\vect{r}-\vect{r}^{\text{ref}}(t) \|_{\infty} \leq \delta\}, \quad t\in[\tau_0, \tau_v),
\end{equation}
%  $$
% S = \{\vect{z}\mid \begin{pmatrix}-\delta_{\underline{x}}\\-\delta_{\underline{y}}\\-\delta_{\underline{z}}\end{pmatrix}\leq \vect{r}(t)-\vect{r}^r(t)\leq \begin{pmatrix}\delta_{\bar{x}}\\\delta_{\bar{y}}\\\delta_{\bar{z}}\end{pmatrix}\}
% $$ 
where $\delta>0$ represents the maximum allowable deviation in each axis. That is, the trajectory is considered as safe if it is within a \emph{safe tube} of radius $\delta$ around the nominal trajectory. Note that it is straightforward to generalize the following results to having different maximum bounds for each axis. 

%In order to ensure $\vect{z}(t)\in\mathcal{S}(t),\ t\in[\tau_0,\tau_v)$, we define candidate CBFs:
%$ h(\vect{z},t)\geq \vect{0}$
% \begin{align} \label{CBF}
%     \vect{h}(\vect{z},t) &= \mat{h_{\overline{x}} & h_{\overline{y}}& h_{\overline{z}} & h_{\underline{x}} &
%     h_{\underline{y}} & h_{\underline{z}}}^T\nonumber\\
%     &=\delta\vect{1}_6 + \mat{(\vect{r}^{\text{ref}}(t)-\vect{r})^T & (\vect{r}-\vect{r}^{\text{ref}}(t))^T}^T,
% \end{align}
To ensure $\vect{z}(t)\in\mathcal{S}_h(t),\; \forall t\in[\tau_0,\tau_v)$, we define the following 6 candidate CBFs for $q\in\{x,y,z\}$:
\begin{equation} \label{CBF}
   %\vect{h}_q(\vect{z},t) = 
   \mat{h_{\overline{q}}\\h_{\underline{q}}} \triangleq\mat{\delta+ q^{\text{ref}}(t)-q(t)\\\delta+q(t)-q^{\text{ref}}(t)}.%, \; q\in\{x,y,z\}.
\end{equation}
Note that all functions $h_{\overline{q}},h_{\underline{q}},\ q\in\{x,y,z\}$ have relative degree 2. % and are $\mathcal{C}^2$ continuous.
%for $d_h \geq 0$ representing the width of the tube around the reference trajectory within which the quadcopter must remain. 
% Note that each CBF in $\vect{h}(\vect{z},t)$ has relative degree $2$.
According to Definition \ref{dfn:cbfhigh}, we choose constants $a_1,a_2$ such that  the roots of $\la^2+a_1\la+a_2$ are reals and negative, and the CBF conditions of relative degree 2 are satisfied for $h_{\overline{q}},h_{\underline{q}},\forall q\in\{x,y,z\}$.  For example,  the CBF condition for $h_{\overline{x}}$ is
\begin{align*}
L_g\bar L_fh_{\overline{x}}\vect{\mu} +\bar L_f^2h_{\overline{x}}+a_1\bar L_fh_{\overline{x}}+a_2h_{\overline{x}} \geq 0,    
\end{align*}
which is equivalent to
\begin{align}\label{cbfx}
\phi_{\overline{x}}^1\vect{\mu}+\phi_{\overline{x}}^2\leq 0,
\end{align}
where
\begin{align*}
    \phi_{\overline{x}}^1&=\vect{x}_W^T,\\
    \phi_{\overline{x}}^2&=-\ddot{x}^{\text{ref}}(t) + a_1(\dot{x}-\dot{x}^{\text{ref}}(t))+a_2(x -x^{\text{ref}}(t)-\delta).
\end{align*}
Putting 6 CBF conditions (for $h_{\overline{x}},h_{\underline{x}},h_{\overline{y}},h_{\underline{y}},h_{\overline{z}},h_{\underline{z}}$), which are all linear constraints on $\vect{\mu}$ for fixed $\vect{z}$, into a QP, and choosing the objective function such that  $\|\Psi(\pi(\vect{z})) - \vect{\mu}\|_2^2$ is minimized, where $\Psi(\pi(\vect{z}))$ yields the nominal virtual inputs, we have the following CBF-QP:
\begin{align} \label{CBF_QP}
\tag{CBF-QP}
    \min_{\vect{\mu}} \quad & \|\Psi(\pi(\vect{z})) - \vect{\mu}\|_2^2\\
    \text{s.t.} \quad &  \phi_{\overline{q}}^1\vect{\mu}+\phi_{\overline{q}}^2\leq 0,\nonumber\\ 
     &\phi_{\underline{q}}^1\vect{\mu}+\phi_{\underline{q}}^2\leq 0,\quad q\in\{x,y,z\},\nonumber
    % \quad & \ddot{h}_i(\vect{z},t,\vect{\mu}) + \vect{k}^T\vect{\eta}_i(\vect{z},t) \geq 0, \quad i\in\{\overline{x},\overline{y},\overline{z},\underline{x},\underline{y},\underline{z}\},\nonumber\\ &\ddot{h}_i(\vect{z},t,\vect{\mu}) = L_f^2h_i(\vect{z}) + L_gL_fh_i(\vect{z})\vect{\mu} + \frac{\partial^2 h_i}{\partial t^2}(t), \nonumber
\end{align}
where the virtual input $\vect{\mu}$ is the decision variable, 
\begin{subequations}
\begin{align}
    \phi_{\overline{q}}^1&=\vect{q}_W^T,\label{phiup1}\\
    \phi_{\overline{q}}^2&=-\ddot{q}^{\text{ref}}(t) + a_1(\dot{q}-\dot{q}^{\text{ref}}(t))+a_2(q-q^{\text{ref}}(t)-\delta),\label{phiup2}\\
    \phi_{\underline{q}}^1&=-\vect{q}^T_W,\label{philow1}\\
    \phi_{\underline{q}}^2&=\ddot{q}^{\text{ref}}(t) + a_1(\dot{q}^{\text{ref}}(t)-\dot{q})+a_2(q^{\text{ref}}(t)-q-\delta),\label{philow2}
\end{align}
\end{subequations}

% \begin{align} \label{CBF_QP}
%     \min_{\vect{\mu}} \quad & \|\Psi(\pi(\vect{z})) - \vect{\mu}\|_2^2\\
%     \text{s.t.} \quad &  \ddot{e}_q(\vect{\mu},t)+a_1\dot{e}_q(\vect{z},t) + a_2e_q(\vect{z},t)\leq a_2 \delta,\nonumber\\ 
%      &\ddot{e}_q(\vect{\mu},t)+a_1\dot{e}_q(\vect{z},t) + a_2e_q(\vect{z},t)\geq -a_2 \delta,\nonumber\\
%     &\quad q\in\{x,y,z\},\nonumber
%     % \quad & \ddot{h}_i(\vect{z},t,\vect{\mu}) + \vect{k}^T\vect{\eta}_i(\vect{z},t) \geq 0, \quad i\in\{\overline{x},\overline{y},\overline{z},\underline{x},\underline{y},\underline{z}\},\nonumber\\ &\ddot{h}_i(\vect{z},t,\vect{\mu}) = L_f^2h_i(\vect{z}) + L_gL_fh_i(\vect{z})\vect{\mu} + \frac{\partial^2 h_i}{\partial t^2}(t), \nonumber
% \end{align}
%where $e_q(\cdot,t) = q(t)-q^{\text{ref}}(t)$, $\Psi(\pi(\vect{z}))$ is ??, 
with $\vect{q}_W$ representing the appropriate world-frame's axis (see Figure \ref{fig:Coord_Frames}). Note that by the dynamic feasibility of the reference trajectory, we can extract $\ddot{q}^{\text{ref}}(t)$ from:
\begin{equation}
    \dot{\vect{z}}^{\text{ref}}(t) = A\vect{z}^{\text{ref}}(t) + B\Psi(\vect{v}^{\text{ref}}(t)), \quad t\in[\tau_0,\tau_v). 
\end{equation}
%and $k_1,k_2 \in \mathbb{R}$ are safety parameters whose choice is described in the following proposition. 
% designed as in \cite{exp_cbf}.
% $\vect{\eta}_i(\vect{z},t) =\mat{h_i(\vect{z},t) & \dot{h}_i(\vect{z},t)}^T \in \mathbb{R}^2$, and Lie derivatives $L_fh_i(\vect{z}) = (\partial h_i/\partial \vect{z})f(\vect{z})$.

The \eqref{CBF_QP} is convex and can be solved in polynomial-time by QP solvers such as MOSEK \cite{mosek} or OSQP \cite{osqp}. Letting $\vect{\mu}^*(t)$ be the solution of \eqref{CBF_QP} for any $(t,\vect{z})\in[\tau_0,\tau_v)\times\mathbb{R}^6$, the safe input to the quadcopter is then  given as 
\begin{equation}\label{safe_input}
    \vect{v}_s(t) = \Psi^{-1}(\vect{\mu}^*(t),\psi(t)),
\end{equation}
where $\psi(t)$ is recovered from $\vect{v}(t) = \pi(\vect{z}(t))$.%\XX{\eqref{safe_input} does not make sense as in the previous page.}
%\XX{$\psi(t)$ is not set to be zero as in the previous section?} \VF{The reference $\psi^{ref}(t)=0$ from the previous section (although this is not required for any result in this paper because the bounds we provide are $\psi$-independent). We allow the nominal controller to specify $\psi(t)\neq0$ which is more general, but this input is unchecked and bypasses the \eqref{CBF_QP} because there is no safety specification/requirement for it. It is completely decoupled. We simply need to use $\psi(t)$ to recover the input $\vect{v}$ from the virtual input $\vect{\mu}$.}

%As QP \eqref{CBF_QP} has multiple CBF constraints, an important question is the feasibility of  \eqref{CBF_QP}. 
The following result shows sufficient conditions that ensure the forward invariance of the set $\mathcal{S}_h(t)$, and feasibility of the \eqref{CBF_QP} with multiple CBF constraints.
% \XX{Remember that $\vect{h}(\vect{z},t)$ is a vector. The constraints in \eqref{CBF_QP} is confusing and needs to be clarified. We can write CBFs for $h_i$ where $i=\{\bar x,...\}$.}


\begin{theorem} \label{CBF_QP_Feasibility}
Consider the system given in \eqref{lin_sys}. If constants $a_1,a_2>0$ are chosen such that  the roots of $\la^2+a_1\la+a_2$ are negative reals $-\lambda_1,-\lambda_2<0$, 
%i.e., the roots are reals and negative, 
%,  the functions $h_{\overline{q}}, h_{\underline{q}}$, $ q\in\{x,y,z\}$ \eqref{CBF} are all CBFs,
%\VF{Is there a way to guarantee that they are CBFs? as in Def 1 of your paper? instead of assuming they are.}, 
%and $k_1,k_2\in\mathbb{R}$ are such that the roots of quadratic polynomial $p(\lambda)=\lambda^2+k_1\lambda+k_2$ are $-\lambda_1,-\lambda_2<0$  and real, and $k_2>0$, 
then \eqref{CBF_QP} is  feasible for all $(t,\vect{z})\in[\tau_0,\tau_v)\times\mathbb{R}^6$. Furthermore, if 
$\|\vect{r}(\tau_0)-\vect{r}^{\text{ref}}(\tau_0) \|_{\infty} \leq \delta$, $\|\dot{\vect{r}}(\tau_0)-\dot{\vect{r}}^{\text{ref}}(\tau_0)+\lambda_1(\vect{r}(\tau_0)-\vect{r}^{\text{ref}}(\tau_0)) \|_{\infty} \leq \lambda_1\delta$,
%$|q(\tau_0)-q^{\text{ref}}(\tau_0)|<\delta$, $|\dot{q}(\tau_0)-\dot{q}^{\text{ref}}(\tau_0)+\lambda_1q(\tau_0)-\lambda_1q^{\text{ref}}(\tau_0)|< \delta\lambda_1,\ \forall  q\in\{x,y,z\}$, 
and $\vect{\mu}^*(t)$ from \eqref{CBF_QP} is locally Lipchitz in $\vect{z}$, then the input $\vect{v}_s(t)$  will render $\|\vect{r}(t)-\vect{r}^{\text{ref}}(t) \|_{\infty} \leq \delta$ for all $t \in [\tau_0,\tau_v)$. Therefore, $\vect{z}(t)\in\mathcal{S}_h(t)$ for all $t\in[\tau_0,\tau_v)$.
% If the pair ($h_{\overline{q}}$, $h_{\underline{q}}$) satisfies the conditions of Theorem 2 \XX{Specific} in \cite{xu2018constrained} for all $q\in$\{x, y, z\}, then the cbf functions grouped in $\vect{h}$ are control-sharing barrier functions. Furthermore, if $\vect{z}(0)\in S(0)$ and $\vect{\dot z}(0)\in S(0)$, then $\vect{v}_s(t)$ will make $\mathcal{S}(t)$ control invariant when $t \in [\tau_0,\tau_v)$.
\end{theorem}
\begin{proof}
The candidate CBFs defined in \eqref{CBF} aim to ensure $q^{\text{ref}}(t)-\delta \triangleq \underline q(t)\leq q(t)\leq \overline q(t)\triangleq q^{\text{ref}}+\delta$, for $q\in \{x,y,z\}$. Recalling that $\vect{\mu} =[\mu_1,\mu_2,\mu_3]^T$, the system model shown in \eqref{lin_sys} is equivalent to $\ddot x=\mu_1$, $\ddot y=\mu_2$, $\ddot z=\mu_3$. Thus, the axes are decoupled which enable us to consider each axis individually. It is easy to see that for $q\in \{x,y,z\}$, $\overline q-\underline q=2\delta$, $\dot{\overline q}-\dot{\underline q}=0$, and $\ddot{\overline q}-\ddot{\underline q}=0$. Since $a_1,a_2>0$, we have $\ddot{\overline q}-\ddot{\underline q}+a_1(\dot{\overline q}-\dot{\underline q})+a_2(\overline q-\underline q)=2a_2\delta>0$. Therefore, condition (15) in Theorem 2 of \cite{xu2018constrained} is satisfied, which implies that $h_{\overline{q}},h_{\underline{q}}$ are control sharing barrier functions. It is easy to see that  functions $h_{\overline{x}},h_{\underline{x}},h_{\overline{y}},h_{\underline{y}},h_{\overline{z}},h_{\underline{z}}$ are all control sharing barrier functions, which means that \eqref{CBF_QP} is  feasible for all $(t,\vect{z})\in[\tau_0,\tau_v)\times\mathbb{R}^6$. Since $\|\vect{r}(\tau_0)-\vect{r}^{\text{ref}}(\tau_0) \|_{\infty} \leq \delta$, $\|\dot{\vect{r}}(\tau_0)-\dot{\vect{r}}^{\text{ref}}(\tau_0)+\lambda_1(\vect{r}(\tau_0)-\vect{r}^{\text{ref}}(\tau_0)) \|_{\infty} \leq \lambda_1\delta$, we have $|q(\tau_0)-q^{\text{ref}}(\tau_0)|\leq\delta$, 
$|\dot{q}(\tau_0)-\dot{q}^{\text{ref}}(\tau_0)+\lambda_1q(\tau_0)-\lambda_1q^{\text{ref}}(\tau_0)|\leq \lambda_1\delta,\ \forall  q\in\{x,y,z\}$, which is equivalent to $h_{\overline{q}}(\tau_0)\geq 0$, $h_{\underline{q}}(\tau_0)\geq0$ and  $\dot h_{\overline{q}}(\tau_0)+\lambda_1 h_{\overline{q}}(\tau_0)\geq0$, $\dot h_{\underline{q}}(\tau_0)+\lambda_1h_{\underline{q}}(\tau_0)\geq0$ for $q\in\{x,y,z\}$. Therefore,  $h_{\overline{q}}(t)\geq 0$, $h_{\underline{q}}(t)\geq 0$ for all $t \in [\tau_0,\tau_v)$ and $q\in\{x,y,z\}$ by \cite{xu2018constrained}, which implies that  $\|\vect{r}(t)-\vect{r}^{\text{ref}}(t) \|_{\infty} \leq \delta$ holds for $t \in [\tau_0,\tau_v)$. The final conclusion follows by the definition of set $\mathcal{S}_h(t)$.
% Suppose the functions $h_{\overline{q}},h_{\underline{q}},\ q\in\{x,y,z\}$ are all CBFs (see Definition 1 in \cite{xu2018constrained}). Then, for any $(\vect{z},t) \in \mathbb{R}^6\times [\tau_0,\tau_v)$, there exist sets $\mathcal{K}_{\overline{q}}\neq \emptyset,\ \mathcal{K}_{\underline{q}}\neq \emptyset,\ q\in\{x,y,z\}$ defined as:
% \begin{align*}
%   \mathcal{K}_{\overline{q}} &\triangleq \{\vect{\mu}:\ddot{e}_q(\vect{\mu},t)+k_1\dot{e}_q(\vect{z},t) + k_2e_q(\vect{z},t)\leq k_2\delta\},\\
%   \mathcal{K}_{\underline{q}} &\triangleq \{\vect{\mu}:-k_2\delta \leq \ddot{e}_q(\vect{\mu},t)+k_1\dot{e}_q(\vect{z},t) + k_2e_q(\vect{z},t)\}.
% \end{align*} 
% Note that $k_2<0 \implies \mathcal{K}_{q} \triangleq \mathcal{K}_{\overline{q}}\cap\mathcal{K}_{\underline{q}} = \emptyset$. Now, taking $k_2>0$ and applying Theorem 1 (iii) of \cite{xu2018constrained}, we determine that the pairs $(h_{\overline{q}},h_{\underline{q}}), \ q\in\{x,y,z\}$ are control-sharing barrier functions (CSBFs). 
% Therefore, $\mathcal{K}_{q} \neq \emptyset, \ q\in\{x,y,z\}, \ \forall (\vect{z},t)\in\mathbb{R}^6\times[\tau_0,\tau_v)$. 
% We now show that $\mathcal{K}_h \triangleq \mathcal{K}_x\cap\mathcal{K}_y\cap \mathcal{K}_z \neq \emptyset$ holds $\forall(\vect{z},t) \in \mathbb{R}^6\times[\tau_0,\tau_v)$ because for any $(\vect{z},t)$, given 3 vectors: $\vect{\alpha}\in\mathcal{K}_{x}$, $\vect{\beta}\in\mathcal{K}_{y}$ and $\vect{\gamma}\in\mathcal{K}_{z}$, we can always construct $\vect{\mu} = \mat{\alpha_1 & \beta_2 & \gamma_3}^T\in\mathcal{K}_h$. Therefore, all the functions $h_{\overline{q}},h_{\underline{q}}, \ q\in \{x,y,z\}$ are CSBFs. Finally, noting that:
% \begin{equation*}
%     h_{\overline{q}}(\vect{z},t) \geq 0, h_{\underline{q}}(\vect{z},t)\geq 0, \ q\in\{x,y,z\} \implies \vect{z}\in\mathcal{S}(t)
% \end{equation*}
% we can apply Theorem 2 in \cite{xu2018constrained} to obtain the final result. 
\end{proof}

% \begin{proof}\XX{Make the proof more clear.}
%   Suppose the pair $(h_{\overline{q}}$, $h_{\underline{q}})$ satisfies the conditions of Theorem 2 in \cite{xu2018constrained}. Then, $h_{\overline{q}}$ and $h_{\underline{q}}$ are CSBFs and allow to define a nonempty set of admissible controls $\mathcal{K}_q$ as in \cite{xu2018constrained}. Consider the decoupled double integrator system \eqref{lin_sys}:
%   \begin{equation}
%       \dot{\vect{q}} = \mat{0 & 1\\ 0 & 0}\vect{q} + \mat{0\\1}\ddot{q},\quad \vect{q} = \mat{q & \dot{q}}^T.
%   \end{equation}
%   It can be shown that the set of admissible controls $\mathcal{K}_q$ imposes restrictions only on input $\ddot{q}$. Therefore, we must have that $\mathcal{K}_x\cap\mathcal{K}_y\cap\mathcal{K}_z\neq \emptyset$ concluding the proof.
% \end{proof}

% \begin{remark}
% In \eqref{CBF_QP}, the input bounds for the decision variable $\vect{\mu}$ is not considered
% \end{remark}

\begin{remark}
The feasibility of \eqref{CBF_QP} can be also seen from the non-emptiness of its admissible input set. 
For $q\in\{x,y,z\}$, define the respective admissible input sets corresponding to CBFs $h_{\overline{q}}$ and $h_{\underline{q}}$ as $\mathcal{K}_{\overline{q}}=\{\vect{\mu}\mid \phi_{\overline{q}}^1\vect{\mu}+\phi_{\overline{q}}^2\leq 0\},\mathcal{K}_{\underline{q}}=\{\vect{\mu}\mid \phi_{\underline{q}}^1\vect{\mu}+\phi_{\underline{q}}^2\leq 0\}$, respectively.  
Note that $-\phi_{\overline{q}}^{2}\geq\phi_{\underline{q}}^2$ holds $\forall q\in\{x,y,z\}$ because $-\phi_{\overline{q}}^2 - \phi_{\underline{q}}^2=2a_2\delta>0$. Since $\mathcal{K}_{x}\triangleq\mathcal{K}_{\overline{x}}\cap \mathcal{K}_{\underline{x}}=\{\vect{\mu}\mid \phi_{\underline{x}}^2\leq \mu_1\leq -\phi_{\overline{x}}^2\}$, we have $\mathcal{K}_{x}\neq\emptyset$. Similarly, 
$\mathcal{K}_{y}\triangleq\mathcal{K}_{\overline{y}}\cup \mathcal{K}_{\underline{y}}=\{\vect{\mu}\mid \phi_{\underline{y}}^2\leq \mu_2\leq -\phi_{\overline{y}}^2\}\neq\emptyset$, and
$\mathcal{K}_{z}\triangleq\mathcal{K}_{\overline{z}}\cup \mathcal{K}_{\underline{z}}=\{\vect{\mu}\mid \phi_{\underline{z}}^2\leq \mu_3\leq -\phi_{\overline{z}}^2\}\neq\emptyset$. It is easy to see that $\mathcal{K}_{x}\cap\mathcal{K}_{y}\cap\mathcal{K}_{z}\neq\emptyset$ as
 $\mathcal{K}_{x},\mathcal{K}_{y},\mathcal{K}_{z}$ impose constraints on $\mu_1,\mu_2,\mu_3$, respectively. Therefore, the admissible input set  of \eqref{CBF_QP} is always non-empty, which guarantees the feasibility of \eqref{CBF_QP}.
\end{remark}

\begin{remark} \label{sdcbf}
The safety guarantee provided by Theorem \ref{CBF_QP_Feasibility} is predicated on the assumption that the control input generated by the \eqref{CBF_QP} can be updated continuously. When a sampled-data controller is implemented in practice, it is still possible to provide the same safety guarantee in continuous time, if we make a slight modification to the CBF conditions of \eqref{CBF_QP}. We refer the readers to our recent work \cite{zhang21interval} that proposes  %computationally efficient algorithms that considers 
sampled-data CBFs suitable for real-time computation, and demonstrates the related algorithms using experiments in real hardware.
%However, implementing the CBF-QP \eqref{CBF_QP} in real hardware is necessarily done in discrete-time. The degradation of the CBF safety guarantees due to discrete-time implementation is an active field of research and addressed in recent works such as \cite{breeden2021control,singletary2020control}.
\end{remark}

For the trajectory obtained from \eqref{CBF_QP}, we can also quantify the upper bounds of $\|\dot{\vect{r}}(t) -\dot{\vect{r}}^{\text{ref}}(t)\|_{\infty}$ and $\|\vect{\mu} - \ddot{\vect{r}}^{\text{ref}}(t)\|_{\infty}$ as shown in the following two corollaries.

\begin{corollary} \label{bounded_vel_corollary}
If the conditions of Theorem \ref{CBF_QP_Feasibility} are satisfied and $\|\dot{\vect{r}}(\tau_0) -\dot{\vect{r}}^{\text{ref}}(\tau_0)\|_{\infty} \leq \frac{2\delta a_2}{a_1}$, then
\begin{equation} \label{cbf_vel_bounds}
    \|\dot{\vect{r}}(t) -\dot{\vect{r}}^{\text{ref}}(t)\|_{\infty} \leq \frac{2\delta a_2}{a_1}
\end{equation}
holds for any $t\in[\tau_0,\tau_v)$. 
\end{corollary}
\begin{proof}
Any feasible solution $\vect{\mu} = (\ddot{x},\ddot{y},\ddot{z})^T$ of \eqref{CBF_QP} satisfies:
\begin{subequations}
\begin{align}
    \ddot{q}&\leq \ddot{q}^{\text{ref}}+a_1(\dot{q}^{\text{ref}}-\dot{q})+a_2(q^{\text{ref}}-q+\delta),\label{veleq1}\\
    \ddot{q}&\geq \ddot{q}^{\text{ref}} + a_1(\dot{q}^{\text{ref}}-\dot{q}) + a_2(q^{\text{ref}} -q - \delta),\label{veleq2}
\end{align}
\end{subequations}
for $q\in\{x,y,z\}$. Since $|q-q^{\text{ref}}|\leq \delta$, \eqref{veleq1}-\eqref{veleq2} are equivalent to:
\begin{subequations}
\begin{align}
    \ddot{q}&\leq \ddot{q}^{\text{ref}}+a_1(\dot{q}^{\text{ref}}-\dot{q})+2a_2\delta,\label{veleq3}\\
    \ddot{q}&\geq \ddot{q}^{\text{ref}} + a_1(\dot{q}^{\text{ref}}-\dot{q}) - 2a_2\delta.\label{veleq4}
\end{align}
\end{subequations}
Define candidate CBFs $\hat{h}_{\overline{q}}$ and $\hat{h}_{\underline{q}}$ as follows:
\begin{align*}
\hat{h}_{\overline{q}}(t)&\triangleq \dot{q}^{\text{ref}}(t)-\dot{q}(t)+ \frac{2\delta a_2}{a_1},\\ \hat{h}_{\underline{q}}(t)&\triangleq \dot{q}(t) - \dot{q}^{\text{ref}}(t) + \frac{2\delta a_2}{a_1}.
\end{align*}
% we have:
% \begin{align*}
%     \ddot{q}&\leq \ddot{q}^{\text{ref}}+a_1(\dot{q}^{\text{ref}}-\dot{q})+2\delta a_2,\\
%     \ddot{q}&\geq \ddot{q}^{\text{ref}} + a_1(\dot{q}^{\text{ref}}-\dot{q}) - 2\delta a_2.
% \end{align*}
% Now 
Then \eqref{veleq3}-\eqref{veleq4} are equivalent to 
$$
\dot{\hat{h}}_{\overline{q}}(t)+a_1\hat{h}_{\overline{q}}(t)\geq 0,\quad \dot{\hat{h}}_{\underline{q}}(t)+a_1\hat{h}_{\underline{q}}(t)\geq 0,
$$
for $q\in\{x,y,z\}$ and any $t\in[\tau_0,\tau_v)$. Therefore,  $\hat{h}_{\overline{q}}$ and $\hat{h}_{\underline{q}}$ satisfy the CBF condition. Furthermore, 
since $\|\dot{\vect{r}}(\tau_0) -\dot{\vect{r}}^{\text{ref}}(\tau_0)\|_{\infty} \leq \frac{2\delta a_2}{a_1}$, we have $\hat{h}_{\overline{q}}(\tau_0)\geq 0$ and $\hat{h}_{\underline{q}}(\tau_0)\geq 0$ for $q\in\{x,y,z\}$. Thus, $\hat{h}_{\overline{q}}(t)\geq 0$ and $\hat{h}_{\underline{q}}(t)\geq 0$ for $t\in[\tau_0,\tau_v)$ by \cite{ames2016control}, which implies that $|\dot{q}^{\text{ref}}(t)-\dot{q}(t)|\leq \frac{2\delta a_2}{a_1}$ for $q\in\{x,y,z\}$ and $t\in[\tau_0,\tau_v)$. The conclusion \eqref{cbf_vel_bounds} follows immediately.
%$\bar{L}_f\bar{h}_{\overline{q}} + L_g\bar{h}_{\overline{q}} \geq -a_1\bar{h}_{\overline{q}}$ and $\bar{L}_f\bar{h}_{\underline{q}} + L_g\bar{h}_{\underline{q}}\geq -a_1\bar{h}_{\underline{q}}$. 
%we can rewrite the inequalities above in the form:
% \begin{align*}
%     \bar{L}_f\bar{h}_{\overline{q}} + L_g\bar{h}_{\overline{q}} &\geq -a_1\bar{h}_{\overline{q}}, \\
%     \bar{L}_f\bar{h}_{\underline{q}} + L_g\bar{h}_{\underline{q}}&\geq -a_1\bar{h}_{\underline{q}}.
% \end{align*}
% where
% \begin{align*}
%     \bar{h}_{\overline{q}}&\triangleq \dot{q}^{\text{ref}}(t)-\dot{q}+ \frac{2\delta a_2}{a_1},\\
%     \bar{h}_{\underline{q}}&\triangleq \dot{q} - \dot{q}^{\text{ref}}(t) + \frac{2\delta a_2}{a_1}.
% \end{align*}
%Finally, we note that $\bar{h}_{\overline{q}}$ and $\bar{h}_{\underline{q}}$ are CBFs by Definition \ref{dfn:cbfhigh} and they certify \eqref{cbf_vel_bounds}.
\end{proof}

%Consider the following corollary bounding the feasible region of \eqref{CBF_QP}:
\begin{corollary} \label{cbf_qp_feas_set_bound_cor}
If the conditions of Theorem \ref{CBF_QP_Feasibility} are satisfied, then the feasible solution $\vect{\mu}$ of \eqref{CBF_QP} is bounded by:
\begin{equation}\label{cbf_qp_feas_set_bound}
    \|\vect{\mu} - \ddot{\vect{r}}^{\text{ref}}(t)\|_{\infty}\leq 4\delta a_2,
\end{equation}
for any $t\in[\tau_0,\tau_v)$.
\end{corollary}
\begin{proof}
Recall that $\|\vect{r}(t)-\vect{r}^{\text{ref}}(t)\|_{\infty}\leq \delta$ by Theorem \ref{CBF_QP_Feasibility} and $\|\dot{\vect{r}}(t) -\dot{\vect{r}}^{\text{ref}}(t)\|_{\infty} \leq 2\delta a_2/a_1$ by Corollary \ref{bounded_vel_corollary}.
These bounds allow us to upper and lower bound the feasible set of \eqref{CBF_QP} by:
\begin{equation*}
    \ddot{q}^{\text{ref}}(t) - 4 \delta a_2\leq \ddot{q}\leq \ddot{q}^{\text{ref}}(t)+4\delta a_2,
\end{equation*}
where $q\in\{x,y,z\}$. Noticing that $\vect{\mu} = (\ddot{x},\ddot{y},\ddot{z})^T$ concludes the proof.
\end{proof}


%By ??, Safety filter \eqref{CBF_QP} guarantees the forward-invariance of set $S$ when $\vect{z}(0)\in S$. 
%Note that the term $\partial^2 h /\partial t^2$ requires finding $\ddot{\vect{r}}^r(t)$. Given the feasibility of the reference trajectory, we can extract this vector from $\dot{\vect{x}}^r(t) = f(\vect{x}^r(t))+g(\vect{x}^r(t))\vect{u}^r(t)$. 
%In order to guarantee the feasibility of \eqref{CBF_QP}, we provide conditions to ensure that candidate CBFs \eqref{CBF} are control-sharing barrier functions \cite{xu2018constrained}.

\subsection{Input Constraints Satisfaction for the CBF-QP Controller}
%\VF{New section}.
%Quadcopters are typically subject to input saturations due to limited actuator effort. If input saturations are not taken into account when designing a controller, the system behavior may become unpredictable. Recent works such as \cite{Flat_trajectory_design} address this issue in the control design process. \XX{Why we talk about ``input saturations are not taken into account''? We already considered it in the previous section. This is not the real motivation for this subsection. The rational is how to  generate a conservative reference trajectory such that when implemented in (CBF-QP) the input saturation will still be respected.}
%In this section, we provide conditions on the trajectory so that controllers filtered by \eqref{CBF_QP} are safe with respect to input saturation. 
%\XX{I changed this paragraph. Please check.} 
The thrust $T$, the roll angle $\phi$, and the pitch angle $\theta$ are related to the decision variable $\vect{\mu}$
% ,which contains the virtual inputs of the double-integrator quadcopter dynamic model, 
in \eqref{CBF_QP} through the relation shown in \eqref{safe_input}. In this subsection, we will consider how the safety constraints on $T,\phi,\theta$ considered in Section \ref{sec:traj} can be respected by the CBF-QP-based tracking controller that filters the nominal controller $\pi(\vect{z})$.
%during tracking with a controller consisting of $\eqref{CBF_QP}$ filtering the nominal controller $\pi(\vect{z})$. 
% the CBF-QP tracking controller, which is filtered by \eqref{CBF_QP} from any nominal  controller. 
Specifically,  we aim to guarantee that $\vect{v}_s(t)\in\mathcal{V}$ where $\vect{v}_s(t)$ is given in \eqref{safe_input} and% \VF{Separated the sets for the legend of Figure \ref{fig:satguard}}: 
\begin{subequations}
\begin{align}
    \mathcal{V}&\triangleq \mathcal{V}_{T}\cap \mathcal{V}_{\xi}, \label{eqV}\\
 \mathcal{V}_{T} &\triangleq  \{\vect{v}\in\mathbb{R}^4\mid 0\leq T\leq \overline{T}\},\\
    \mathcal{V}_{\xi}&\triangleq \{\vect{v}\in\mathbb{R}^4: |\phi|,|\theta|\leq \epsilon<\pi/2\},
\end{align}
\end{subequations}
with $\epsilon$ the upper bound for $|\phi|$ and $|\theta|$, and $\overline{T}$ the upper bound for $T$. Note that the constraint $\vect{v}_s(t)\in\mathcal{V}$ is important to ensure actuator saturation will not happen.
%desirable, for example, to avoid actuator saturation.


% Given saturation bounds $\underline{q},\overline{q}$, the saturation function $\sat(q(t))$ is defined as follows:
% \begin{equation}
%     \sat(q(t))\triangleq \begin{cases} \underline{q}, & q(t)<\underline{q},\\ q(t), & \underline{q} \leq q(t) \leq \overline{q},\\ \overline{q}, & \overline{q} < q(t).
%     \end{cases}
% \end{equation}
% Suppose now the reduced input signal $\vect{v}(t)$ given as in \eqref{reduced_input} is subject to saturations when received by the quadcopter such that it becomes:
% \begin{equation}
%     \bar{\vect{v}}(t) \triangleq \text{sat}(\vect{v}(t)),
%     % \mat{\beta(T(t),0,T_{sat}) \\ \beta(\phi(t),-\epsilon_{sat},\epsilon_{sat}) \\ \beta(\theta(t),-\epsilon_{sat},\epsilon_{sat})\\
%     % \psi(t)},
% \end{equation}
% acting element-wise with $\overline{T} = T_{sat}$, $\underline{T} = 0$, $\overline{\phi} = -\underline{\phi} = \epsilon_{sat}$, $\overline{\theta} = -\underline{\theta} =\epsilon_{sat}$ and $\overline{\psi} = \infty$, $\underline{\psi} = -\infty$.
% % \VF{TODO Saturation function generic and then use bounds as arguments}:
% % \begin{subequations}
% % \begin{align}
% %     s_T(T(t)) &= \begin{cases} 0, & T(t)<0,\\ T(t), & 0 \leq T(t)\leq T_{sat}, \\ T_{sat}, & T_{sat} < T(t),\end{cases}\\
% %     s_{\epsilon}(\rho(t)) &= \begin{cases}-\epsilon_{sat}, & \rho(t)<-\epsilon_{sat},\\ \rho(t), & -\epsilon_{sat} \leq \rho(t) \leq \epsilon_{sat},\\ \epsilon_{sat}, & \epsilon_{sat} < \rho(t),
% %     \end{cases}
% % \end{align}
% % \end{subequations}
% % \XX{the definitions of these two functions are confusing. $T(t)=T$ when $0 \leq T\leq T_{sat}$? $T$ is also a function of $t$}\VF{I added the time dependency to all equations where appropriate. In my previous definition, $T$ was supposed to be the same as $T(t)$. Only $T_{sat}$ and $\epsilon_{sat}$ are constants, the rest are functions.}
% where $T_{sat}, \epsilon_{sat}>0$ are known saturation limits. To avoid unexpected saturation (i.e. to ensure $\bar{\vect{v}}(t) = \vect{v}(t)$), the reduced control input must satisfy $\vect{v}(t) \in \mathcal{V}$ with:
% \begin{equation}
%     \mathcal{V} \triangleq \{\vect{v}\in\mathbb{R}^4: 0\leq T\leq T_{sat}, \ |\phi|,|\theta|\leq \epsilon_{sat}<\pi/2\}.\label{eqV}
% \end{equation}



%We now have conditions on the reference trajectory that safeguard $\vect{v}_s(t)$ against input saturation.

%\XX{the following proposition is a mess.  notations unclear, the proof impossible to follow.} \VF{I changed it to make it more clear, below is the old version commented for reference.}

% \begin{proposition} \VF{OLD} \label{input_safeguard_prop}
% Consider the system given in \eqref{lin_sys} and suppose the conditions of Theorem \ref{CBF_QP_Feasibility} are satisfied. If the state $\vect{z}^{\text{ref}}(t)$ and input $\vect{v}^{\text{ref}}(t)$ trajectories satisfy:
% \begin{subequations} \label{sat_safeguard_cst}
% \begin{align} 
%     \dot{\vect{z}}^{\text{ref}}(t) &= A\vect{z}^{\text{ref}}(t) + B\vect{v}^{\text{ref}}(t), \label{satprop_dyn_feas}\\
%     0\leq T^{\text{ref}}(t) &\leq T_{sat} - 4\sqrt{3}\delta a_2,\label{satprop_T}\\
%     \|A_{\epsilon}^{sat}\ddot{\vect{r}}^{\text{ref}}(t)\|_2&\leq \ddot{z}^{\text{ref}}(t) +g -4\delta a_2(1+\sqrt{2\bar{\epsilon}_{sat}}),\label{satprop_eul}\\
%     A_{\epsilon}^{sat}&=\sqrt{\bar{\epsilon}_{sat}}\mat{1 & 0 & 0\\ 0 & 1 & 0 \\ 0 & 0 & 0},\nonumber\\ \bar{\epsilon}_{sat}&\triangleq \cot^2\epsilon_{sat}\nonumber
%     % (1-\sin^2\epsilon_{sat})/\sin^2\epsilon_{sat}\nonumber
% \end{align}
% \end{subequations}
% then $\vect{v}_s(t)\in\mathcal{V}$ with $\vect{v}_s(t)$ given in \eqref{safe_input} and $\mathcal{V}$ given in \eqref{eqV}.
% \end{proposition}
% \begin{proof} 
% \XX{where is \eqref{satprop_dyn_feas} from? Does it make sense - why it can be satisfied? what does it mean  \eqref{satprop_dyn_feas} is satisfied? $\vect{v}^{\text{ref}}(t)\in\R^4?$}
% Consider first the thrust element of $\vect{v}_s(t)$ given as $T(t) = \|\vect{\mu}^*(t)+g\vect{z}_W\|_2$ \XX{what's $\vect{\mu}^*(t)$}and notice:
% \begin{align*}
%     T(t)-T_{sat} &=\|\vect{\mu}^*(t)+g\vect{z}_W\|_2-T_{sat}\\
%     % &\leq \|\vect{\mu}^*(t)+g\vect{z}_W\|_2-\|\ddot{\vect{r}}^{\text{ref}}(t)+g\vect{z}_W\|_2-4\sqrt{3}\delta a_2\\
%     &\leq \|\vect{\mu}^*(t)\|_2-\|\ddot{\vect{r}}^{\text{ref}}(t)\|_2 - 4\sqrt{3}\delta a_2\\
%     &\leq \|\vect{\mu}^*(t)-\ddot{\vect{r}}^{\text{ref}}(t)\|_2-4\sqrt{3}\delta a_2\\
%     &\leq 0
% \end{align*}
% where \XX{why $\|\vect{\mu}^*(t)+g\vect{z}_W\|_2\leq \|\vect{\mu}^*(t)\|_2-\|\ddot{\vect{r}}^{\text{ref}}(t)\|_2$?}we used \eqref{flat_T_map}, \eqref{satprop_T}, the triangle inequality, and the 2-norm equivalent of Corollary \ref{cbf_qp_feas_set_bound_cor}. Note that $\ddot{\vect{r}}^{\text{ref}}(t)$ can be extracted from $\dot{\vect{z}}^{\text{ref}}(t)$. Next, we examine the roll $\phi(t)$ and pitch $\theta(t)$ elements of $\vect{v}_s(t)$. 
% Begin by defining the set
% \begin{equation*}
%     \mathcal{K}_{\epsilon}^{sat}\triangleq \{\vect{\mu}:\|A_{\epsilon}^{sat}\vect{\mu}\|_2\leq \mu_3+g\},
% \end{equation*}
% and projected signals $\vect{\mu}^*_{2D}(t)\triangleq(\mu^*_1(t),\mu^*_2(t))$ and $\ddot{\vect{r}}^{\text{ref}}_{2D}\triangleq (\ddot{x}^{\text{ref}}(t),\ddot{y}^{\text{ref}}(t))$. Then, by Corollary \ref{cbf_qp_feas_set_bound_cor}, $\|\vect{\mu}^*_{2D}(t)-\ddot{\vect{r}}^{\text{ref}}_{2D}(t)\|_{\infty}\leq 4\delta a_2$, and by norm equivalence:
% \begin{equation*}
%     4\sqrt{2}\delta a_2 \geq \|\vect{\mu}^*_{2D}(t)-\ddot{\vect{r}}^{\text{ref}}_{2D}(t)\|_{2} \geq \|\vect{\mu}^*_{2D}(t)\|_2 - \|\ddot{\vect{r}}^{\text{ref}}_{2D}(t)\|_{2}.
% \end{equation*}
% We can now multiply on both sides by $\sqrt{\bar{\epsilon}_{sat}}\geq 0$ and drop the time-dependence of $\vect{\mu}^*(t)$ for convenience to obtain:
% \begin{align*}
%     4\sqrt{2\bar{\epsilon}_{sat}}\delta a_2 &\geq \|\sqrt{\bar{\epsilon}_{sat}}\vect{\mu}^*_{2D}\|_2 - \|\sqrt{\bar{\epsilon}_{sat}}\ddot{\vect{r}}^{\text{ref}}_{2D}(t)\|_2\\
%     &\geq\|A_{\epsilon}^{sat}\vect{\mu}^*\|_2 - \|A_{\epsilon}^{sat}\ddot{\vect{r}}^{\text{ref}}(t)\|_2\\
%     &\geq \|A_{\epsilon}^{sat}\vect{\mu}^*\|_2 -\ddot{z}^{\text{ref}}(t)-g+4\delta a_2(1+\sqrt{2\bar{\epsilon}_{sat}}).
% \end{align*}
% \XX{how is the last inequality derived?}
% We can finally use $|\mu_3 - \ddot{z}^{\text{ref}}(t)|\leq 4\delta a_2$ to rewrite the inequality as:
% % \begin{equation*}
% %     4\sqrt{2\bar{\epsilon}_{sat}}\delta a_2\geq \|A_{\epsilon}^{sat}\vect{\mu}^*\|_2 -\mu_3-4\delta a_2-g+4\delta a_2(1+\sqrt{2\bar{\epsilon}_{sat}}).
% % \end{equation*}
% \begin{equation*}
%     \|A_{\epsilon}^{sat}\vect{\mu}\|_2\leq \mu_3+g \implies \mu\in\mathcal{K}_{\epsilon}^{sat}.
% \end{equation*}
% \XX{How is the implication above derived?}
%  Then, by Lemma \ref{epsilon_cone_lemma}, \XX{why conditions of Lemma \ref{epsilon_cone_lemma} are satisfied?} we have that $|\phi(t)|,|\theta(t)|\leq \epsilon_{sat}$ for any $t\in[\tau_0,\tau_v)$,
% \end{proof}
%\VF{I revamped the Proposition for input safeguarding below. Left the previous version above commented for reference.}
\begin{proposition} \label{input_safeguard_prop}
Consider the system given in \eqref{lin_sys} and suppose that the conditions of Theorem \ref{CBF_QP_Feasibility} are satisfied. If trajectories of the state $\vect{z}^{\text{ref}}(t)$ and the input $\vect{v}^{\text{ref}}(t)$ are dynamically feasible by satisfying \eqref{dyn_model_acc}  and 
the following conditions hold for all $t\in[\tau_0,\tau_v)$:
\begin{subequations} \label{sat_safeguard_cst}
\begin{align}
0\leq T^{\text{ref}}(t) &\leq \overline{T} - 4\sqrt{3}\delta a_2,\label{satprop_T}\\
\|A_{\epsilon}\ddot{\vect{r}}^{\text{ref}}(t)\|_2&\leq \ddot{z}^{\text{ref}}(t) +g -4\delta a_2(1+|\cot\epsilon|\sqrt{2}),\label{satprop_eul}
\end{align}
\end{subequations}
where $A_{\epsilon} =|\cot\epsilon|\mat{\vect{x}_W & \vect{y}_W & \vect{0}_{3\times 1}}$, then $\vect{v}_s(t)\in\mathcal{V}$ with $\vect{v}_s(t)$ given in \eqref{safe_input} and $\mathcal{V}$ given in \eqref{eqV}.
\end{proposition}
\begin{proof} 
Let $\vect{\mu}^*(t)$ be the solution of \eqref{CBF_QP}. Note that by norm equivalence, \eqref{cbf_qp_feas_set_bound} implies:
\begin{equation}
    \|\vect{\mu}^*(t) - \ddot{\vect{r}}^{\text{ref}}(t)\|_2\leq \sqrt{3}\|\vect{\mu}^*(t)-\ddot{\vect{r}}^{\text{ref}}(t)\|_\infty\leq4\sqrt{3}\delta a_2,\label{eqnprp3eq1}
\end{equation}
for any $t\in[\tau_0,\tau_v)$. 

First, consider the thrust element of $\vect{v}_s(t)$ given as $T(t) = \|\vect{\mu}^*(t)+g\vect{z}_W\|_2$ by the map $\Psi^{-1}$. By $\eqref{satprop_T}$ we have
\begin{equation}
    T^{\text{ref}}(t)+4\sqrt{3}\delta a_2 = \|\ddot{\vect{r}}^{\text{ref}}(t)+g\vect{z}_W\|_2+4\sqrt{3}\delta a_2 \leq \overline{T}.\label{eqnprp3eq2}
\end{equation}
%Let $\tilde{T}(t) \triangleq T(t) - T_{sat}$ and notice that by the reverse triangle inequality and the above relations:
% \begin{align*}
%     \tilde{T}(t) &=\|\vect{\mu}^*(t)+g\vect{z}_W\|_2-T_{sat}\\
%     &\leq \|\vect{\mu}^*(t)+g\vect{z}_W\|_2-\|\ddot{\vect{r}}^{\text{ref}}(t)+g\vect{z}_W\|_2-4\sqrt{3}\delta a_2\\
%     % &\leq \|\vect{\mu}^*(t)\|_2 + \|g\vect{z}_W\|_2-\|\ddot{\vect{r}}^{\text{ref}}(t)\|_2 - \|g\vect{z}_W\|_2- \gamma\\
%     % &\leq \|\vect{\mu}^*(t)\|_2 -\|\ddot{\vect{r}}^{\text{ref}}(t)\|_2 - \gamma\\
%     &\leq \|\vect{\mu}^*(t)-\ddot{\vect{r}}^{\text{ref}}(t)\|_2-4\sqrt{3}\delta a_2 \\
%     % &\leq \|\vect{\mu}-\ddot{\vect{r}}^{\text{ref}}(t)\|_\infty - 4\delta a_2\\
%     &\leq 0.
% \end{align*}
Therefore, 
\begin{align*}
T(t) - \overline{T} &=\|\vect{\mu}^*(t)+g\vect{z}_W\|_2-\overline{T}\\
    &\leq \|\vect{\mu}^*(t)+g\vect{z}_W\|_2-\|\ddot{\vect{r}}^{\text{ref}}(t)+g\vect{z}_W\|_2-4\sqrt{3}\delta a_2\\
    % &\leq \|\vect{\mu}^*(t)\|_2 + \|g\vect{z}_W\|_2-\|\ddot{\vect{r}}^{\text{ref}}(t)\|_2 - \|g\vect{z}_W\|_2- \gamma\\
    % &\leq \|\vect{\mu}^*(t)\|_2 -\|\ddot{\vect{r}}^{\text{ref}}(t)\|_2 - \gamma\\
    &\leq \|\vect{\mu}^*(t)-\ddot{\vect{r}}^{\text{ref}}(t)\|_2-4\sqrt{3}\delta a_2 \\
    % &\leq \|\vect{\mu}-\ddot{\vect{r}}^{\text{ref}}(t)\|_\infty - 4\delta a_2\\
    &\leq 0,
\end{align*}
where the first inequality is from \eqref{eqnprp3eq2}, the second inequality is from the reverse triangle inequality, and the third inequality is from \eqref{eqnprp3eq1}. Thus, $T(t)\leq \overline{T}$ for $t\in[\tau_0,\tau_v)$. Note that $\ddot{\vect{r}}^{\text{ref}}(t)$ can be extracted from $\dot{\vect{z}}^{\text{ref}}(t)$ which is given by dynamic feasibility of the reference trajectory. 

%\XX{I stop here. The notations have not been changed for the rest of the proof. Please check!}\VF{I changed the notation below. Maybe there is a cleaner way to show that $\sqrt{2}$ is sufficient (instead of $\sqrt{3}$) in this case without defining $\ddot{\vect{r}}_{2D}^{\text{ref}}(t).$}
Next, we examine the roll $\phi(t)$ and pitch $\theta(t)$ elements of $\vect{v}_s(t)$. We drop the time-dependence of terms for convenience. Let $\mathcal{K}_{\epsilon}$ be given as in \eqref{ang_bounds_cone}. 
% Define a set
% \begin{equation*}
%     \mathcal{K}_{\epsilon_v}\triangleq \{\vect{\mu}\in\mathbb{R}^3:\|A_{\epsilon_v}\vect{\mu}\|_2\leq \mu_3+g\}.
% \end{equation*}
Note that
\begin{align}
\lv\mat{\mu^*_1\\\mu^*_2}\rv_2&\leq \lv\mat{\mu^*_1\\\mu^*_2}-\mat{\ddot x^{\text{ref}}\\\ddot y^{\text{ref}}}\rv_{2}+\lv\mat{\ddot x^{\text{ref}}\\\ddot y^{\text{ref}}}\rv_2\nonumber\\
&\leq 4\sqrt{2}\delta a_2+\lv\mat{\ddot x^{\text{ref}}\\\ddot y^{\text{ref}}}\rv_2\label{ineqAep}
\end{align}
where the first inequality is from the 
triangle inequality and the second inequality is from norm equivalence and Corollary \ref{cbf_qp_feas_set_bound_cor}.
% and projected signals $\vect{\mu}^*_{2D}(t)\triangleq(\mu^*_1(t),\mu^*_2(t))$ and $\ddot{\vect{r}}^{\text{ref}}_{2D}\triangleq (\ddot{x}^{\text{ref}}(t),\ddot{y}^{\text{ref}}(t))$. Then, by Corollary \ref{cbf_qp_feas_set_bound_cor}, $\|\vect{\mu}^*_{2D}(t)-\ddot{\vect{r}}^{\text{ref}}_{2D}(t)\|_{\infty}\leq 4\delta a_2$, and by norm equivalence and the triangle inequality:
% \begin{equation*}
%     4\sqrt{2}\delta a_2 \geq \|\vect{\mu}^*_{2D}(t)-\ddot{\vect{r}}^{\text{ref}}_{2D}(t)\|_{2} \geq \|\vect{\mu}^*_{2D}(t)\|_2 - \|\ddot{\vect{r}}^{\text{ref}}_{2D}(t)\|_{2}.
% \end{equation*}
Therefore,
\begin{align*}
\|A_{\epsilon}\vect{\mu}^*\|_2 =|\cot\epsilon|\lv\mat{\mu^*_1\\\mu^*_2}\rv_2
&\leq 
4|\cot\epsilon|\sqrt{2}\delta a_2+ \|A_{\epsilon}\ddot{\vect{r}}^{\text{ref}}\|_2  \\
&\leq\ddot{z}^{\text{ref}} -4\delta a_2 +g \\
&\leq \mu_3^*+g
\end{align*}
where the first inequality is from \eqref{ineqAep} and the definition of $A_{\epsilon}$, the second inequality is from \eqref{satprop_eul}, and the third inequality is from the fact that $\ddot{z}^{\text{ref}} \leq \mu_3^*+4\delta a_2$, which can be derived from  $|\mu_3^* - \ddot{z}^{\text{ref}}|\leq 4\delta a_2$ by Corollary \ref{cbf_qp_feas_set_bound_cor}. 
% We can now multiply on both sides by $|\cot\epsilon_v|\geq 0$ and drop the time-dependence of $\vect{\mu}^*(t)$ and $\vect{\mu}^*_{2D}(t)$ for convenience to obtain:
% \begin{align*}
%     4|\cot\epsilon_v|\sqrt{2}\delta a_2 &\geq \||\cot\epsilon_v|\vect{\mu}^*_{2D}\|_2 - \||\cot\epsilon_v|\ddot{\vect{r}}^{\text{ref}}_{2D}(t)\|_2\\
%     &\geq\|A_{\epsilon_v}\vect{\mu}^*\|_2 - \|A_{\epsilon_v}\ddot{\vect{r}}^{\text{ref}}(t)\|_2,
%     % &\geq \|A_{\epsilon_v}\vect{\mu}^*\|_2 -\ddot{z}^{\text{ref}}(t)-g+4\delta a_2(1+|\cot\epsilon_v|\sqrt{2}),
% \end{align*}
% % where the last inequality is obtained from \eqref{satprop_eul}.
% where the last inequality follows from the definition of $A_{\epsilon_v}$. 
% We can now use \eqref{satprop_eul} to write:
% \begin{equation*}
%     \|A_{\epsilon_v}\vect{\mu}^*\|_2\leq \ddot{z}^{\text{ref}}(t) +g -4\delta a_2.
% \end{equation*}
% By Corollary \ref{cbf_qp_feas_set_bound_cor},  $|\mu_3^* - \ddot{z}^{\text{ref}}|\leq 4\delta a_2$, which  implies that  $\ddot{z}^{\text{ref}} \leq \mu_3^*+4\delta a_2$. Therefore, 
% \begin{equation*}
%     \|A_{\epsilon_v}\vect{\mu}^*\|_2\leq \mu_3^*+g \implies \vect{\mu}^*(t)\in\mathcal{K}_{\epsilon_v},
% \end{equation*}
% for any $t\in[\tau_0,\tau_v)$. 
Thus, $\mu^*(t)\in\mathcal{K}_{\epsilon}$ and, by Lemma \ref{epsilon_cone_lemma}, 
%\XX{why conditions of Lemma \ref{epsilon_cone_lemma} are satisfied?}\VF{Changed the Lemma to account for this} 
we have that $|\phi(t)|,|\theta(t)|\leq \epsilon$ for all $t\in[\tau_0,\tau_v)$. This completes the proof.
\end{proof}

\begin{remark}
The conditions \eqref{satprop_T}-\eqref{satprop_eul}  for the reference trajectory can be easily guaranteed in continuous-time with the method presented in Section \ref{sec:traj}, and the resulting \eqref{FLAT-SOCP} remains convex. Intuitively, conditions \eqref{satprop_T}-\eqref{satprop_eul} shrink the cone \eqref{ang_bounds_cone} and the sphere \eqref{T_sphere} constraints on the trajectory acceleration (see Figure \ref{fig:satguard}) enough so that all feasible solutions of \eqref{CBF_QP} satisfy the desired safety constraints when transformed through the map $\Psi^{-1}$. This means that the transformed optimal solution $\vect{v}_s(t) = \Psi^{-1}(\vect{\mu}^*(t),\psi) \in \mathcal{V}$ for any $\psi$ and any $t\in[\tau_0,\tau_v)$.
\begin{figure}[htp]
    \centering
    \includegraphics[width=8cm]{satguard_prop3_shrink.eps}
    \caption{
    %Effect of conditions \eqref{satprop_T}-\eqref{satprop_eul}  on the feasible region of \eqref{FLAT-SOCP}. 
    Conditions \eqref{satprop_T}-\eqref{satprop_eul}  shrink the cone that bounds the Euler angles ($\vect{x}^{\text{ref}}(t)\in\mathcal{X}_{\xi}$) and the sphere that bounds the thrust ($\vect{u}^{\text{ref}}(t)\in\mathcal{U}_{T}$) to guarantee that safe inputs $\vect{v}_{s}(t)\in\mathcal{V}$. This figure shows the shrinkage with typical bounds $\epsilon = 1.0472$, $\overline{T}=2g$ and the same safety parameters  $\delta=0.1$, $a_2=8$ that we use in Examples \ref{example:vcp}, \ref{example:slowzone} and \ref{example:forest} in Sec. \ref{sec:exper}}
    \label{fig:satguard}
\end{figure}
\end{remark}

\begin{remark}\label{remark12}
Instead of  shrinking the cone  and the sphere  constraints for \eqref{FLAT-SOCP}, 
%One might argue that the conditions of Proposition \ref{input_safeguard_prop} are too conservative for certain  $\delta$, $a_2$ parameter values. An
an alternative method to ensuring $\vect{v}_s(t)\in\mathcal{V}$ is to convert the CBF-based QP into a SOCP by adding two additional SOC constraints as follows:
\begin{align*} %\label{CBF-SOCP}
% \tag{CBF-SOCP}
    \min_{\vect{\mu}} \quad & \|\Psi(\pi(\vect{z})) - \vect{\mu}\|_2^2\\
    \text{s.t.} \quad &  \phi_{\overline{q}}^1\vect{\mu}+\phi_{\overline{q}}^2\leq 0,\nonumber\\ 
     &\phi_{\underline{q}}^1\vect{\mu}+\phi_{\underline{q}}^2\leq 0,\quad q\in\{x,y,z\},\nonumber\\
     &\|\vect{\mu}+g\vect{z}_W\|_2\leq \overline{T},\\
     &\|A_{\epsilon}\vect{\mu}\|_2\leq \mu_3+g,
\end{align*}
with $A_{\epsilon}$ given in Proposition \ref{input_safeguard_prop}. Note that the optimization shown above is still a convex program, but its feasibility is no longer guaranteed.
%which can be solved efficiently using convex solvers such as MOSEK, but its feasibility is no longer guaranteed.
%, while this formulation lifts \eqref{CBF_QP} to SOCP and we can no longer guarantee its feasibility, convexity is maintained.

% Proposition \ref{input_safeguard_prop} maintains the feasibility guarantees of \eqref{CBF_QP} provided by Theorem \ref{CBF_QP_Feasibility} by placing the tracking safety constraints in \eqref{FLAT-SOCP} instead.

% While the conditions of Proposition \ref{input_safeguard_prop} guarantee input safeguarding for all solutions of \eqref{CBF_QP}, and maintain the feasibility of the problem, they require designing conservative trajectories. An alternative way to safeguard the solutions of \eqref{CBF_QP} with respect to saturation is to introduce constraints in virtual space that, while lifting the problem to SOCP and potentially compromising its feasibility, still provide input saturation safeguarding in the reduced input space and retain convexity. Namely, we can add SOC constraints:
% \begin{align*}
%     \|\vect{\mu}+g\vect{z}_W\|_2\leq T_{sat},\\
%     \|A^{sat}_{\epsilon}\vect{\mu}\|_2\leq \mu_3+g,
% \end{align*}
% with $A_{\epsilon}^{sat}$ given as in Proposition \ref{input_safeguard_prop}.
\end{remark}

%%%%%%%%%%%%%%%%%%%%%%%%%%%%%%%%%%%
\section{Experiments}\label{sec:exper}
%\VF{There are no simulations in this section.}
We demonstrate the effectiveness and versatility of the proposed framework using three experimental examples. Each example highlights certain different aspects of the framework. For all examples, \eqref{FLAT-SOCP} is solved in MATLAB with Yalmip \cite{Lofberg2004} and solver MOSEK \cite{mosek} at a rate of 30-100 Hz depending on its complexity, and \eqref{CBF_QP} is solved in Python3 with OSQP \cite{osqp} at a rate of 100 Hz. The tracking data is obtained from an OptiTrack motion capture system and a Crazyflie2.1 nano quadcopter. We provide videos of all presented examples in \cite{victor21video}.

% \XX{In the caption of each figure, specify which example it is for.}

\begin{example}\label{example:vcp}


The first example demonstrates rigorous satisfaction of the safety constraints in continuous-time by the reference trajectory obtained from \eqref{FLAT-SOCP} and the bounded tracking performance provided by the CBF-QP controller. % \subsection*{Example 1: Bounded Thrust, Euler Angles and Tracking}
First, we solve \eqref{FLAT-SOCP} to generate a reference trajectory satisfying \eqref{state_input_const} where the following safety constraints are required to be respected for all $t\in[\tau_0,\tau_v)$:
\begin{align*}
    \|\dot{\vect{r}}^{\text{ref}}(t)\|&\leq \overline{v},\\
    |\phi^{\text{ref}}(t)|,|\theta^{\text{ref}}(t)|&\leq \epsilon,\\
    \underline{T}\leq T^{\text{ref}}(t)&\leq \overline{T},\\
    |p^{\text{ref}}(t)|,|q^{\text{ref}}(t)|&\leq \overline{\omega}.
\end{align*}
The parameters of the trajectory and the safety cosntraints are given in Table \ref{tab:sample_trajectory_1}, and the parameters of the waypoints are given in Table \ref{tab:wp_1}. The initial/final conditions are chosen as $\vect{p}^m_0=\cdots =\vect{p}^m_4= (0,0,0)^T, m\in\{0,f\}$.   
%For illustration, we highlight the geometry of the constraints that act on the $2$nd order VCPs to guarantee bounded thrust, roll and pitch. 
%, though other constraints are also enforced. 
%Then, we demonstrate that \eqref{CBF_QP} endows a poorly-designed nominal controller with bounded tracking. 
%\XX{Show or refer what constraints are considered in this example.} \VF{I added reference to \eqref{state_input_const} which are all state/input constraints.} \XX{If we add more constraints, then we have difficulty showing the constraints are satisfied rigorously? Is that the reason only three constraints are considered?} 
%\VF{This trajectory has more than 3 constraints, all constraints for state and input (no obs. avoidance) are active. The three constraints I show geometrically are the most illustrative and are shown for that purpose only.} 
With these parameters, we solve \eqref{FLAT-SOCP}  and show the resulting position trajectory $\vect{r}^{\text{ref}}(t)$ in Figure \ref{fig:vcp_st_r0} along with the waypoints and control points. The geometry of the constraints enforced for the $2$nd order VCPs $\vect{P}_j^{(2)}$ is shown in Figure \ref{fig:vcp_st_r2}:
\begin{itemize}
    \item Inclusion in the conic surface (red) given by $\mathcal{K}_{\epsilon}$ defined in \eqref{ang_bounds_cone} to guarantee $\vect{x}^{\text{ref}}(t)\in \mathcal{X}_{\xi}$.
    \item Inclusion in the sphere (green) given by \eqref{vcp_Tmax} to guarantee upper-bounded thrust $T(t)\leq \overline{T}$.
    \item Restriction above hyperplane (cyan) given by \eqref{vcp_Tmin} to guarantee lower-bounded thrust $\underline{T}\leq T(t)$.
\end{itemize}
The safety constraint satisfaction is verified in Figure \ref{fig:vcp_st_cst} by transforming the flat output trajectory $\vect{\sigma}^{\text{ref}}(t)=(\vect{r}^{\text{ref}}(t),0) $ via \eqref{flat_map} to the state $\vect{x}^{\text{ref}}(t)$ and input $\vect{u}^{\text{ref}}(t)$ trajectories. From  Figure \ref{fig:vcp_st_cst}, it can be observed that all considered safety constraints are satisfied in continuous time. It can be also seen that the values of the vector $\vect{\zeta}=(\zeta_1,...,\zeta_{36})^T$, which is illustrated by the light red line, lower-bound the thrust over local segments as defined in \eqref{cp_omega}. %\XX{Explain the light red line.}

\begin{table} [bt]
\renewcommand{\arraystretch}{1.3}
\caption{Trajectory parameters for Example \ref{example:vcp}}
\label{tab:sample_trajectory_1}
\centering
\begin{tabular}{|c|c||c|c||c|c|}
\hline
$\tau_0$ [s] & 0 &  $n_0, n_f$ [-] & 4 & $\epsilon$ [deg] & 1.75 \\
\hline
$\tau_v$ [s] & 30 & $n_{wp}$ & 8 &  $\underline{T}$ $[m/s^2]$ & 9.7 \\
\hline
N [-] & 40 & $d_{wp}$ [m]& 0.05 & $\overline{T}$  $[m/s^2]$ & 9.9 \\
\hline
d [-] & 5  & $\overline{v}$ $[m/s]$ & 0.5 & $\overline{\omega}$ [deg/s] & 1.5\\
\hline
\end{tabular}
\end{table}

\begin{table} [bt]
\renewcommand{\arraystretch}{1.3}
\caption{Trajectory waypoints for Example \ref{example:vcp}}
\label{tab:wp_1}
\centering
\begin{tabular}{|c|c c c c c c c c|}
\hline
$i$ & 1 & 2 & 3 & 4 & 5 & 6 & 7 & 8\\
\hline
& -0.15 & -0.75 & 0.65 & 0.65 & -0.5 & -0.6 & 0.4 & 0.25\\
$\vect{p}_i^{wp}$ & 0.25 & 0.6 & -0.65 & 0.5 & 0.5 & -0.6 & -0.4 & 0.25\\
& 0.25 & 0.5 & 0.25 & 0.25 & 0.75 & 0.5 & 0.4 & 0.25\\
\hline
$t_i$ & 4.5 & 7.8 & 12.6 & 15.3 & 18 & 21  & 24 & 27\\
\hline
\end{tabular}
\end{table}



\begin{figure}[htp]
    \centering
    \includegraphics[width=8cm]{vcp_st_r0.eps}
    \caption{The position trajectory generated in Example \ref{example:vcp} with the parameters in Tables \ref{tab:sample_trajectory_1} and \ref{tab:wp_1}. We show the position trajectory (dashed black) along with the waypoints (red circles) and the control points (connected, blue circles).}
    \label{fig:vcp_st_r0}
\end{figure}

\begin{figure}[htp]
    \centering
    \includegraphics[width=8cm]{vcp_st_r2.eps}
    \caption{
    %\XX{Do we have other constraints in  FLAT-SOCP? }\VF{Added sentence at end of caption.} 
    Visualization of constraints enforced for the $2$nd order VCPs $\vect{P}_j^{(2)}$ in Example \ref{example:vcp}. Constraint \eqref{cp_ang} represents inclusion within the red cone. Constraint \eqref{cp_Tmax} is the inclusion within the green sphere (only cusp is shown). Constraint \eqref{cp_Tmin} forces the VCPs to lie above the cyan plane (global bound). Note that the local counterpart \eqref{cp_omega_Tmin} is not depicted here. Note also that these are not the only constraints added to \eqref{FLAT-SOCP} in Example \ref{example:vcp}.}
    \label{fig:vcp_st_r2}
\end{figure}

\begin{figure}[htp]
    \centering
    \subfigure[State constraints. The velocity respects $\|\dot{\vect{r}}(t)\|\leq 0.5$ m/s and the Euler angles respect $|\phi(t)|,|\theta(t)|\leq 1.75$ deg for all $t\in[0,30)$.]{
    \includegraphics[width=8cm]{vcp_st_cst_x.eps}
    }
    \subfigure[Input constraints. The thrust respects $9.7\;  \text{m/$s^2$}\leq T(t) \leq 9.9\;  \text{m/$s^2$}$ and the angular velocities respect $|p(t)|,|q(t)|\leq 1.5$ deg/s. The light red line indicates the values of the vector $\vect{\zeta}=(\zeta_1,\ldots,\zeta_{36})^T$ lower-bounding the thrust over local segments as shown in \eqref{cp_omega}.]{
    \includegraphics[width=8cm]{vcp_st_cst_u.eps}
    }
    \caption{Verification that the safety constraints for state and inputs in Example \ref{example:vcp} are satisfied in continuous-time.}
    \label{fig:vcp_st_cst}
\end{figure}

\begin{figure}[htb]
    \centering
    \subfigure[Tracking performance of the nominal controller $\pi(\vect{z})$ with (blue) and without (magenta) the \eqref{CBF_QP}. The data is obtained from a real flight experiment using a Crazyflie2.1 nano quadcopter.]{
    \includegraphics[width=8cm]{vcp_cbf_3D.eps}
    }
    \subfigure[Measured values of  $h_{\overline{x}},h_{\underline{x}},h_{\overline{y}},h_{\underline{y}},h_{\overline{z}},h_{\underline{z}}$ defined in \eqref{CBF} when tracking with the filtered (blue) and unfiltered (magenta) nominal controller $\pi(\vect{z})$. Note that safety is satisfied when the values of CBFs are nonnegative, which means that the real trajectory is within the prescribed tube ($\delta =0.1$) around the nominal  trajectory. Tracking with unfiltered nominal controller (magenta) violates safety as $h_{\overline{x}}(t,\vect{z}(t)) <0$ and $h_{\overline{y}}(t,\vect{z}(t)) <0$ for some $t$. On the other hand, the CBF-QP tracking controller always respects safety as $h_{\overline{x}},h_{\underline{x}},h_{\overline{y}},h_{\underline{y}},h_{\overline{z}},h_{\underline{z}}$ are always non-negative, indicating safety. ]{\includegraphics[width=8cm]{vcp_cbf_h.eps}}
    \caption{Verification of the bounded-tracking property of \eqref{CBF_QP} with experimental data for Example \ref{example:vcp}. 
    %\XX{Use subfigure environment to add caption for the top and bottom plot individually? The top plot is about trajectoy, the bottom is about CBFs, which are totally different.}\VF{I changed it.} 
    }
    \label{fig:cbf_qp}
\end{figure}

Next, we implement in a Crazyflie2.1 nano quadcopter the CBF-QP controller to track the reference trajectory generated by \eqref{FLAT-SOCP} with a poorly-designed nominal controller $\pi(\vect{z})$. We show that, when filtered through the \eqref{CBF_QP} with parameters $a_1 = 6, a_2 =8$, safe trajectory tracking can be achieved with $\delta=0.1$. We compare the performance of the CBF-QP tracking controller and the unfiltered nominal controller in Figure \ref{fig:cbf_qp}. It can be seen that the CBF-QP controller  achieves much better tracking performance than the unfiltered nominal controller. Note that the bounding constraints are satisfied by the  CBF-QP controller, which is demonstrated by the fact that CBFs defined in \eqref{CBF} satisfy $h_{\overline{q}},h_{\underline{q}}\geq 0$ for $q\in\{x,y,z\}$, while these conditions are not true for the unfiltered nominal controller.   
%\XX{Can we make the legend and the axis labels of the figures better? When I zoom in, the quality is not good.}\VF{I exported the figures from MATLAB directly as PNG. If I zoom in the PDF version of the paper, the quality seems reasonable to me. Maybe we can discuss this in our individual meeting.}

\end{example}

\begin{example}\label{example:slowzone}
The second example  demonstrates the proposed framework is suitable for real-time (online) computation.
% low computational burden
% real-time capability
% of the proposed framework.
We let the quadcopter take off from, and land on a moving platform, and fly through a velocity-constrained zone while passing two nearby waypoints. Note that velocity-constrained areas are commonplace  in  safety-critical applications (e.g., when the quadcopter operates near humans). The reference trajectory is recomputed by solving \eqref{FLAT-SOCP} in real time, where the initial and final positions of the reference trajectory are both set to be the position of the moving platform. The tracking controller is obtained by solving \eqref{CBF_QP} online.

The safety specifications for the reference trajectory are: 
\begin{subequations} \label{slowzone_specs}
\begin{align}
    \vect{r}^{\text{ref}}(t) \in \mathcal{S}_v, \quad
    \|\dot{\vect{r}}^{\text{ref}}(t)\| \leq \overline{v} = 0.5, \quad t\in[3,6),\label{slowzone_specs_local}\\
    \vect{r}^{\text{ref}}(0) =  \vect{r}^{\text{ref}}(9) = \vect{p}_{0f}(t), \quad \dot{\vect{r}}^{\text{ref}}(0) = \dot{\vect{r}}^{\text{ref}}(9) = 0,\label{slowzone_specs_ext} \\
    \|\vect{r}^{\text{ref}}(2.5) - \vect{p}_1^{wp}\|_2\leq d_{wp},\|\vect{r}^{\text{ref}}(6.5) - \vect{p}_2^{wp}\|_2\leq d_{wp}\label{slowzone_specs_wp},
\end{align}
\end{subequations}
where $\vect{p}_{0f}(t)$ is the position of the platform at time $t$, $\vect{p}_1^{wp} = (0.75, 0.6, 1.1)^T$ and $\vect{p}_2^{wp} = (-0.75,0.6,1.1)^T$ are the waypoints with desired radius $d_{wp} = 0.2$, and the constrained-velocity zone $\mathcal{S}_v$ is defined as: 
\begin{align*}
    \mathcal{S}_{v}&\triangleq \{\vect{r}\in\mathbb{R}^3: \|A_s\vect{r}+\vect{b}_s\|_2\leq 1\}, \label{ex2_local_specification} \\
    A_s&=\text{diag}(1.33, 13.3, 13.3),\\
    \vect{b}_s &= \mat{0& -10& -14.7}^T,
\end{align*}
which describes an ellipsoid centered at $\vect{p}_h = (0,0.75,1.1)^T$ chosen to ensure the quadcopter clears the \emph{hoop} obstacle.
% \begin{equation*}
%     S_{v}\triangleq \{\vect{r}\in\mathbb{R}^3\mid (-0.5,0.2,0.2)^T\leq\vect{r} \leq (0.5,0.8,0.8)^T\}
% \end{equation*}

%
% and it must pass near ($d_{wp} = 0.2$) waypoints $\vect{p}_1^{wp} = (0.75, 0.6, 1.1)$ and $\vect{p}_2^{wp} = (-0.75,0.6,1.1)$ at times $t_1 = 2.5$ and $t_2 = 6.5$. Notice that \eqref{slowzone_cst} indicates position and velocity bounds over a local time interval, which is commonplace  in  safety-critical applications (e.g., when the quadcopter operates near humans),  and will result in constraints over a subset of the control points and VCPs.
% We assume that there are 3 waypoints $\vect{p}^{wp}_1 = (-0.3, 0.6, 0.35)$, $\vect{p}^{wp}_2 = (0, 0.4, 0.65)$ and $\vect{p}^{wp}_3 = (0.3, 0.6, 0.35)$ with associated times $t_1 = 3.5$, $t_2 = 5$, $t_3 = 6.5$. 

We begin by considering a B-spline curve $\vect{r}^{\text{ref}}(t)$ as defined in \eqref{b-spline_curve} with $N= 45$, $d=5$, $\tau_0 = 0$ and $\tau_v = 9$. Including \eqref{slowzone_specs_ext} and \eqref{slowzone_specs_wp} in \eqref{FLAT-SOCP} is straightforward. However, \eqref{slowzone_specs_local} indicates position and velocity bounds over a local time interval and will result in constraints over a subset of the control points and VCPs.
% The trajectory $\vect{r}^{\text{ref}}(t)$ must satisfy \eqref{const_wp} with $d_{wp} = 0.2$. 
Since $\vect{\tau}$ is clamped and uniform, we have that $[3,6)\in[\tau_{18},\tau_{33}) = [2.85, 6.15)$. Using Proposition \ref{inclusion_prop}, we obtain the following constraints on the control points and $1st$ order VCPs:
\begin{align*}
  &\vect{P}_j \in S_v, \quad j=13,\dots, 32,\\
  \|&\vect{P}_j^{(1)}\|_2\leq \overline{v}, \quad j=14,\dots,32.
\end{align*}
% We suppose there is a dynamic takeoff/landing platform for the quadcopter with center $\vect{p}_{0f}(t)\in\mathcal{C}_v$, where 
% \begin{equation*}
%     \mathcal{C}_v \triangleq \{\vect{r} \in \mathbb{R}^3\mid (-0.5,0,0)^T\leq\vect{r}\leq (0.5,0.2,0)^T\}
% \end{equation*}
% is a subregion of the ground plane and  $\vect{p}_{0f}(t)$ is unknown. 
% The initial and final conditions of the trajectory are given as
% \begin{align*}
%     \vect{r}^{\text{ref}}(\tau_0) &= \vect{r}^{\text{ref}}(\tau_v) = \vect{p}_{0f}(t),\\
%     \dot{\vect{r}}^{\text{ref}}(\tau_0) &= \dot{\vect{r}}^{\text{ref}}(\tau_v) = 0.
% \end{align*}
% \XX{I don't understand why the initial and final positions are the same. What makes sense to me is the platform is constantly changing, and the trajectory for safe landing is computed in real time.} \VF{The initial position in this example is not where the quadcopter is. The initial position is the landing platform. Therefore, as the landing platform moves, the start/end of the trajectory change. This is still online-replanning.}
In the experiment, measurements of $\vect{p}_{0f}(t)$ are obtained at 100 Hz and the convexity of \eqref{FLAT-SOCP} allows us to recompute trajectories $\vect{r}^{\text{ref}}(t)$ satisfying \eqref{slowzone_specs} in real-time. The tracking is done by a nominal controller $\pi(\vect{z})$ filtered by \eqref{CBF_QP} with parameters $\delta = 0.1$, $a_1 = 6$ and $a_2 = 8$.


%at high rates. 
% The above conditions are impossible to satisfy because $\vect{p}_{0f}(t)$ is unpredictable. However, we can achieve close performance by recomputing $\vect{r}(t)$ as measurements of $\vect{p}_{0f}(t)$ are obtained. 
We implement this experiment using a Crazyflie2.1 nano quadcopter and a Turtlebot 2 as the dynamic platform. Figure \ref{fig:slowzone_3D} shows the considered scenario, the quadcopter tracking data and the measurements of $\vect{p}_{0f}(t)$.
% the evolution of the trajectories $\vect{r}^{\text{ref}}(t)$ as they are recomputed by solving \eqref{FLAT-SOCP} with new measurements of $\vect{p}_{0f}(t)$. 
The Crazyflie2.1 takes off near $\vect{p}_{0f}(0)$ and lands near $\vect{p}_{0f}(9)$, by tracking the reference trajectory $\vect{r}^{\text{ref}}(t)$, which is recomputed at 30Hz. Additionally, we show the velocity profiles $\|\dot{\vect{r}}^{\text{ref}}(t)\|_2$ of the reference trajectories in Figure \ref{fig:slowzone_vel}. It can be seen that all computed trajectories respect the constrained-velocity safety specification $\|\dot{\vect{r}}^{\text{ref}}(t)\|_2\leq 0.5$ for all $t\in[3,6)$. We also show the measured velocity magnitude of the quadcopter throughout the experiment.

% shows velocity profiles of computed trajectories when the platform positions $\vect{p}_{0f}$ are randomized. Furthermore, the expected velocity profiles for trajectories with unconstrained speed are also shown for comparison. 
% Additionally, Figure \ref{fig:slowzone} provides velocity profiles for trajectories with and without the enforced speed limit for comparison. 
% We realize a modified version of this example experimentally where the safe zone $S_v$ is an ellipsoid passing through a hoop and there is only one waypoint at its center. The robots are a Crazyflie2.1 nano quadcopter and a Turtlebot 2 robot as the takeoff/landing platform. The Turtlebot 2 is remotely controlled by a human to simulate the unknown $\vect{p}_{0f}(t)$. \XX{This paragraph is impossible to parse. I have no idea what this experiment is about after reading it. Is the experiment video kept the same? Why change the setting/statement of this example?} \VF{I changed the paragraph to make it a bit more clear. I may bring back the ellipsoid $S_v$ to match the video, but the velocity profiles for that case are less exciting than the ones currently in Figure \ref{fig:slowzone}}

\begin{figure}
    \centering
    \subfigure[\label{fig:slowzone_3D}Scenario for Example \ref{example:slowzone} showing the waypoints $\vect{p}^{wp}_1$ and $\vect{p}^{wp}_2$ with radius $d_{wp}=0.2$, the constrained-velocity zone $\mathcal{S}_v$ through the \emph{hoop} obstacle, and the measured trajectories of the dynamic platform (i.e., $\vect{p}_{0f}(t)$) and  Crazyflie2.1. 
    % New trajectories $\vect{r}^{\text{ref}}(t)$ are obtained at $30$Hz. However, we only show every 10th trajectory in the plot to avoid cluttering the figure. 
    The quadcopter takes off near $\vect{p}_{0f}(0)$ and lands near $\vect{p}_{0f}(9)$, which is accomplished by recomputing the reference trajectory $\vect{r}^{\text{ref}}(t)$. 
    % Trajectories are drawn in increasingly darker shades of grey as they are computed.
    % while the platform moved according to $\vect{p}_{0f}(t)$. Measurements o satisfying $\vect{r}(t)\in S_v$ when $t\in[3,7)$ with a randomly drawn start and end point $\vect{p}_{0f}$. We show the case when $\|\dot{\vect{r}}(t)\|_2\leq\overline{v},(t\in[3,7))$ is (blue) and isn't (black) enforced for the same takeoff/landing point\VF{Maybe there is no need for this subfigure and we only show the velocity profiles.}.
    %\XX{I have zero clue about what this figure means, even with the help of the caption. Is it the same experiment as previous described?}
    % \XX{Is (a) from real data? The pattern is too regular, which could mean that the trajectory is not replanning in real time. Moving platform does not necessarily mean hard - if you know its movement in advance, then it is the same as stationary. You should have set the initial/end position randomly. The  zoom-in ellipsoid is impossible to read in a A4 paper. It is hard to tell which end point correspond to which starting point. }
    ]{\includegraphics[width=0.4\textwidth]{slowzone_3D_exp_notraj.eps}}
    \subfigure[\label{fig:slowzone_vel}Velocity profiles of the trajectories in Example \ref{example:slowzone}. Each reference trajectory satisfies the safety specification $\|\dot{\vect{r}}^{\text{ref}}(t)\|_2\leq 0.5$ for $t\in[3,6)$, and its color is drawn in increasingly darker shades of grey as it is computed. Also shown is the speed of the Crazyflie2.1 as it tracks the evolving reference trajectory. ]
    % It can be seen that the limited-velocity trajectories (blue) satisfy the constraint over interval [2.78,7.22) which is more conservative than needed. This effect can be reduced by increasing the number of control points $N$. 
    % \XX{why does this figure make sense? The velocity is constrained not in [4,6) sec.}
    {\includegraphics[width=8cm]{slowzone_vel_exp.eps}}
    \caption{
    % \XX{The figures, especially the top one, need more explanation.}\VF{Added explanation as subfigures.}
    % \XX{Why this comparison makes any sense? Who cares about the comparison? Shouldn't this example about online replanning?}\VF{Yeah maybe there is no need to show the unconstrained velocity, but the example does also sowcase local velocity bounds, and I think seeing that the bound extends to $[2.78,7.22)$ is an interesting point to show. The online replanning in this case is illustrated by the 200 trajectories for random start/endpoint (moving platform) that are computed fast (~110Hz).} 
    % Example \ref{example:slowzone} trajectories computed as the dynamic platform moved and measurements of $\vect{p}_{0f}(t)$ are obtained. The replanning rate was $30$Hz but we show only every 10th computed trajectory to avoid crowding the figure. We also show the tracking data of the Crazyflie2.1 nano quadcopter.
    Visualization of the experimental data for Example \ref{example:slowzone}.
    }
    
    % comparison of trajectories with and without bounded velocity. We simulate an unpredictable, dynamic takeoff/landing platform by drawing random points $\vect{p}_{0f}$.
    % expected for moving takeoff and landing platforms. All trajectories satisfy the safety constraints \eqref{slowzone_cst} and 
    % The average solve time of \eqref{FLAT-SOCP} was 0.0088 seconds. \VF{There are versions of this figure with 1, 10, 20, 100 and 200 trajectories on the left menu, in case you prefer one of those instead.}} %Please see the accompanying video for an implementation of this example with real robots.}
    \label{fig:slowzone}
\end{figure}

\end{example}


% \begin{example} \label{example:slowzone_old}
% % \subsection{Locally-Bounded Velocity and Online Replanning}
% \VF{I am changing this example significantly to address some comments/present things more clearly.} The second example demonstrates bounded velocity enforcement when flying through a safe zone $S_v$ and online replanning by having a moving takeoff/landing platform for the quadcopter. We note that low-speed zones are commonplace in safety-critical applications, e.g., when the quadcopter operates near humans.

% Let the safe zone $S_v$ be defined by the SOC in position $\vect{r}$ coordinates: % \XX{in what coordinate?} \VF{Is that what you meant?}: 
% \begin{align*}
%     S_v &\triangleq \{\vect{r}\in\mathbb{R}^3:\|A_v\vect{r}+\vect{b}_v\|_2\leq 1\},\\
%     % \mat{A_v&\vect{b}_v} &= \mat{0 & -2 & 0 & 2\\ 13.33 & 0 & 0 & 0 \\ 0 & 0 & 13.33 & -13.33}\;
%     \textit{where}\;\;A_v &= \mat{0 & -2 & 0 \\ 13.33 & 0 & 0  \\ 0 & 0 & 13.33 },\;
%     \vect{b}_v = \mat{2\\  0 \\ -13.33}.
% \end{align*}
% %\XX{What's the norm in $S_v$? I found the norm is missing in many places. I corrected some of them, but please double-check others.}
% We consider a B-spline curve $\vect{r}(t)$ as defined in \eqref{b-spline_curve} with $N =40$, $d=5$, $\tau_0=0$ and $\tau_v = 10$. There is a  waypoint $\vect{p}^{wp} = (0,1,1)$ with associated time $t_{wp}=5$ to be enforced as $\vect{r}(t_{wp})=\vect{p}^{wp}$. Suppose that the safety specifications are given as:
% %during times $t\in[4,6)$ we want to guarantee:
% \begin{equation} \label{slowzone_cst_old}
%     \vect{r}(t) \in S_v, \quad
%     \|\dot{\vect{r}}(t)\| \leq \overline{v} = 0.5, \quad t\in[4,6).
% \end{equation}
% Notice that \eqref{slowzone_cst} indicate position and velocity bounds over a local time interval and will result in constraints over a subset of the control points (and VCPs). Since $\vect{\tau}$ is clamped and uniform, we have that $[4,6)\in[\tau_{19},\tau_{27}) = [3.89, 6.11)$. We can now use Proposition \ref{inclusion_prop} to obtain the indices of the control points (and VCPs) that must be constrained, resulting in: %\XX{need more explanation about the constraints below}\VF{I changed the above paragraph to explain a bit more.}
% \begin{align*}
%   &\vect{P}_j \in S_v, \quad j=14,\dots, 26,\\
%   &\bigg\|\sum_{i=0}^N\vect{P}_ib_{1,i+1,j+1}\bigg\|_2\leq \overline{v},\quad j = 15,\dots,26.
% \end{align*}
% The initial and final positions are allowed to change as the takeoff and landing platforms are constantly moving. We realize this experiment with a Crazyflie2.1 nano quadcopter and a Turtlebot 2 robot as the takeoff/landing platform. A number of sample trajectories and corresponding velocity profiles for this example are shown in Figure \ref{fig:slowzone}. \XX{Need to explain fig. \ref{fig:slowzone} more if we were to keep it. I suggest we make it simpler. }







% Inspired by NASA's Mars 2020 mission, we demostrate collaboration between a Crazyflie2.1 nano quadcopter implementing the proposed framework and a tele-operated Turtlebot 2. There are two disjoint sites $S_1$ and $S_3$ that the quadcopter must investigate. However, it must takeoff from and land on the Turtlebot 2. Furthermore, the Turtlebot 2  will provide safe passage between the two disjoint sets (see the provided video and Figure \ref{fig:hitchhiker}).
% \begin{figure}
%     \centering
%     \includegraphics[width=8cm]{experiments/hitchhiker.PNG}
%     \caption{Experiment inspired by NASA's Mars 2020 mission. A Crazyflie2.1 quadcopter collaborates with a Turtlebot 2 to safely explore two disjoint sets.}
%     \label{fig:hitchhiker}
% \end{figure}
% \end{example}

\begin{example} \label{example:forest}
% \subsection{Safe Replanning in Cluttered Environments}
The third example demonstrates safe trajectory planning and tracking in a cluttered environment. The definitions of the start/end zones $\mathcal{C}_1,\ldots,\mathcal{C}_4$, cuboid obstacles $\mathcal{O}_1,\mathcal{O}_2$, \emph{hoop} $\mathcal{O}_{3}$, \emph{desk} $\mathcal{O}_{4}$, \emph{ladder} $\mathcal{O}_{5}$, safe flight space $\mathcal{F}$, and the convex sets  $\mathcal{S}_1,\ldots,\mathcal{S}_6$ (chosen as in Proposition \ref{convex_sets_prop}) are described in Appendix \ref{appendix:forest}. The environment configuration is shown in  Figure \ref{fig:forest}.  We design trajectories that avoid the obstacles and take the quadcopter from any $\vect{p}^0\in\mathcal{C}_i$ to any  $\vect{p}^f\in\mathcal{C}_j$ such that:
\begin{align*} %\label{forest_cst}
    & \vect{r}^{\text{ref}}(t) \in\mathcal{F},\  t\in[t_0,t_f],\\ & \vect{r}^{\text{ref}}(t_0) = \vect{p}^0 , \quad \vect{r}^{\text{ref}}(t_f) = \vect{p}^f. 
\end{align*}
% consider the following problem:
% \begin{problem} \label{prob:forest}
% Given a pair $(i,j), \ i,j\in \{1,\dots,4\}$, find a trajectory satisfying:
% \begin{equation*}
%     \vect{r}^{\text{ref}}(t) \in\mathcal{F}, \quad \vect{r}^{\text{ref}}(t_0) \in \mathcal{C}_i, \quad \vect{r}^{\text{ref}}(t_f)\in \mathcal{C}_j.
% \end{equation*}
% \end{problem}
Because of the choice of convex sets $\{\mathcal{S}_j\}$, given start zone $\mathcal{C}_i$  and end zone $\mathcal{C}_j$, we can always find a finite sequence ($n_{s}<\infty$) of indices $(k_1,k_2,\ldots,k_{n_s}), \ k_l\in\{1,\ldots,6\}, \ l=1,\ldots,n_s$ such that
%\XX{what's $n_s$??}\VF{The number of overlapping convex sets. See Proposition \ref{convex_sets_prop}.}\XX{I know. Need to explain it again here. }
%\XX{ what's $i,j$ below?} \VF{$i$ is the index of the starting set $\mathcal{C}_i$ and $j$ is the index of the ending set $\mathcal{C}_j$. $i\in\{1,\ldots,4\}$.}\XX{I know. The indices are not explained in the equations below.}
\begin{align*}
    &(\mathcal{C}_i \cup \mathcal{S}_{k_1} \cup \mathcal{S}_{k_2} \cup \ldots \cup \mathcal{S}_{k_{n_s}} \cup \mathcal{C}_j)\in\mathcal{F},\\
    &\mathcal{C}_i \cap \mathcal{S}_{k_1} \neq \emptyset, \quad \mathcal{S}_{k_{n_s}}\cap\mathcal{C}_j\neq \emptyset,\\
    &\mathcal{S}_{k_l}\cap \mathcal{S}_{k_{l+1}}\neq \emptyset , \quad l = 1,\ldots,n_s-1,
\end{align*}
where $i,j\in\{1,\ldots,4\}$ are the indices of the start zone $\mathcal{C}_i$ and the end zone $\mathcal{C}_j$, respectively. The sequence of sets $(\mathcal{S}_{k_1},\ldots,\mathcal{S}_{k_{n_s}})$ satisfies the conditions of Proposition \ref{convex_sets_prop} and we can use it to obtain constraints for \eqref{FLAT-SOCP} such that its solution results in a B-spline curve $\vect{r}^{\text{ref}}(t)$ that safely navigates the cluttered environment.
%We describe the setup in detail and show sample trajectories. The accompanying video shows the live experiment with a Crazyflie 2.1 nano quadcopter. 

% We have now satisfied the conditions of Proposition \ref{convex_sets_prop} and can use it to obtain constraints for \eqref{FLAT-SOCP} such that the solution results in a B-spline curve $\vect{r}^{\text{ref}}(t)$ that navigates the cluttered environment as desired. 

In the experiment, we set $\vect{p}^{0}\in\mathcal{C}_i$ as the current position of the quadcopter, and choose a random point $\vect{p}^{f}\in\mathcal{C}_j$ as the next goal position where $\mathcal{C}_j\in\{1,\ldots,4\}$ is also randomly picked. Then, we find a sequence of convex sets $(\mathcal{S}_{k_1},\ldots,\mathcal{S}_{k_{n_s}})$ as above and obtain constraints for \eqref{FLAT-SOCP} from Proposition \ref{convex_sets_prop}. Finally, we solve the optimization problem \eqref{FLAT-SOCP} and generate a reference trajectory. When the quadcopter finishes tracking the trajectory using the CBF-QP-based controller, the process is repeated for a new random $\vect{p}^{f}\in\mathcal{C}_k,k\in\{1,\ldots,4\}$. We repeat this process 6 times, computing each trajectory on-the-fly. The result is a smooth, continuous trajectory that safely navigates the cluttered environment, shown in Figure \ref{fig:forest}. In the experiment, we forced the last $\vect{p}^f$ to match the takeoff location of the quadcopter. The resulting sequence of points visited by the quadcopter was $\vect{p}_1\in\mathcal{C}_2, \vect{p}_2\in\mathcal{C}_3, \vect{p}_3\in\mathcal{C}_2,\vect{p}_4\in\mathcal{C}_1, \vect{p}_5\in\mathcal{C}_3, \vect{p}_6\in\mathcal{C}_4, \vect{p}_1\in\mathcal{C}_2$, and  $\vect{p}_l,l=1,\ldots,6$ are shown in Figure \ref{fig:forest}. 


% \begin{equation*}
%     \mat{\vect{p}_0 &\vect{p}_1 & \vect{p}_2} = \mat{-0.8 & -0.95 & -0.67 & 0.96 & -0.88 & 0.74 \\
%     0.5 & 0.54\\
%     0 & 1.14}
% \end{equation*}

% the quadcopter takes off from $\vect{p}^0 =(-0.8,0.5,0)^T\in\mathcal{C}_2$ and must visit 5   lands at the same point at the end of the flight. We choose 5 other random points position $\vect{p}^f\in\mathcal{C}_j$ is chosen 

% In the experiment, we choose $\vect{p}^0\in\mathcal{C}_i$ as the current position of the quadcopter assuming it starts in $\mathcal{C}_i$,
% % random point $\vect{p}^{0}\in\mathcal{C}_i$ as the initial position where $\mathcal{C}_i,i\in\{1,\ldots,4\}$ is also randomly picked, 
% and choose a random point $\vect{p}^{f}\in\mathcal{C}_j$ as the final position where $\mathcal{C}_j,j\in\{1,\ldots,4\}$ is also randomly picked. We generate a trajectory that begins at $\vect{p}^0$, ends at $\vect{p}^f$ and avoids the obstacles. Then, the trajectory is tracked using the CBF-QP controller. When the end position in set $\mathcal{C}_j$ is reached, the procedure is repeated.
% In this manner, we generate trajectories that randomly explore the cluttered environment. For illustration, we repeat this process 200 times to generate the long continuous trajectory in Figure \ref{fig:forest} as the concatenation of 200 smooth trajectories with an 
% % \XX{what Figure \ref{fig:forest} tells us?} 
% average solve time for each of the 200 segments of 0.0135 seconds (segments may have different duration, but all of them are longer than 5 seconds). In the video experiment, the procedure is repeated only 5 times and the final $\vect{p}^f\in\mathcal{C}_j$ is set to be a designated landing spot. \XX{I thought the trajectory replanning is ``continuous'' in the sense that $\mathcal{C}_j$ is the end point of a trajectory and also the initial point of the next trajectory. It is not how the experiment is done, right? } \XX{Did you actually did 200 experiments? I think this is way too much.}\VF{The trajectory replanning is indeed ``continuous'', I adjusted the paragraph above to reflect this. I did not track 200 trajectories, only 5. I specified this in the paragraph above.} \XX{You generated only 5 trajectories offline, and then track them no matter where the initial conditions are? I don't know what this means. Shouldn't the final position of the trajectory be $p^f$. which is randomly chosen for each trajectory? Why the 200 trajectories are so different, if only 5 reference trajectories exist - tracking performance is so bad? Why do you want to demonstrate using 200 (instead of 20) trajectories as in Example 2? The description of Example 3 is VERY  confusing.}\VF{Figure \ref{example:forest} does not show any tracking data. It is only trajectories. I show many to illustrate that it can be done fast but I guess 20 is probably enough. The trajectories in the actual tracking experiment are not computed offline, they are computed on-the-fly whenever the quadcopter reaches its goal, and they start at its current position and end at a random point. The process is done only 5 times, though.}%Example trajectories with random initial position $\vect{p}^0$ and random final position $\vect{p}^f$ are shown in Figure \ref{fig:forest}. 
%\XX{Need more explanation about the figures here. }\VF{I explained it more in detail.}


\begin{figure}
    \centering
    \includegraphics[width=7.5cm]{forest_exp_top.eps}
    \includegraphics[width=7.5cm]{forest_exp_3D.eps}
    \caption{Experimental setup of the cluttered environment and the trajectories generated  for Example \ref{example:forest}. The designated takeoff/landing position is $\vect{p}_1\in\mathcal{C}_2$. The other points $\vect{p}_l,l=2,\ldots,6$ are chosen randomly from the zones $\mathcal{C}_i,i=1,\ldots,4$. The sequence of points visited by the quadcopter is $\vect{p}_1\in\mathcal{C}_2, \vect{p}_2\in\mathcal{C}_3, \vect{p}_3\in\mathcal{C}_2,\vect{p}_4\in\mathcal{C}_1, \vect{p}_5\in\mathcal{C}_3, \vect{p}_6\in\mathcal{C}_4, \vect{p}_1\in\mathcal{C}_2$. The reference trajectories are shown in black and the real trajectory of the recorded tracking data of the Crazyflie 2.1 quadcopter is shown in blue.  We emphasize that the trajectories were generated on-the-fly as each end point was drawn after completing the tracking of the previous trajectory. The average solve-time of \eqref{FLAT-SOCP} was 0.0135 seconds (see Remark \ref{online_replaning_remk}). The supplementary video is given in \cite{victor21video}.} 
    % The plot shows 200 trajectories starting at random points $\vect{p}^0\in\mathcal{C}_i$ and ending at random points $\vect{p}^f\in\mathcal{C}_j$. The average solve time for each trajectory was 0.0135 seconds.\VF{I added versions of this figure with 20 trajectories instead of 200 to the Figures/forest folder. Note that each sub-trajectory (connecting $\mathcal{C}_i$ and $\mathcal{C}_j$) is plotted in increasingly dark black color. So, the start of the overall trajectory is white/gray and the end is dark black. }\VF{I added also the exact trajectory that the video follows. I will try to get tracking data to show also.}} %Please see the accompanying video for an implementation of this example with a real Crazyflie2.1 nano quadcopter.}
    \label{fig:forest}
\end{figure}
\end{example}

\begin{remark} \label{online_replaning_remk}
For experiments that require online replanning, we use Yalmip's \texttt{optimizer} objects to reduce the overhead compiling time, which greatly improves solve time of parametric problems. Typical compile times for an \texttt{optimizer} object of \eqref{FLAT-SOCP} are about $0.7$ seconds and, as previously stated, any subsequent solving of the parameterized \eqref{FLAT-SOCP} takes between $0.01$ and $0.033$ seconds. Note also that even though the segment-wise constraints of Proposition \ref{convex_sets_prop} are not easily parameterized, we can compile multiple \texttt{optimizer} objects that cover all possible scenarios offline. Then, replanning can still be done online (and easily parallelized) for unstructured scenarios as demonstrated in Example \ref{example:forest}. 
% Since the segment-wise constraints of Proposition \ref{convex_sets_prop} are not easily parameterized, we construct a $4\times4$ matrix $\mathbb{O}$ of \texttt{optimizer} objects covering all possible outcomes of drawing a random $(i,j)$ pair. From here, we have two options, if we know exactly the $(i,j)$ pair, we can use the appropriate \texttt{optimizer} object from $\mathbb{O}$ to solve \eqref{FLAT-SOCP}. If we do not know which \texttt{optimizer} object to use, we can simply solve all of them (this operation can be easily parallelized) and discard the objects that return a certificate of infeasibility. If multiple \texttt{optimizer} objects are feasible, then we can choose the most optimal one.
\end{remark}


\section{Conclusions} \label{sec:Conclusions}
We presented a framework based on convex optimization that achieves rigorous continuous-time safety for the planning and tracking of quadcopter trajectories. We also showed several experimental examples that demonstrate the usefulness of the presented framework. We highlight that the presented trajectory optimization method can easily generalize to other differentially flat systems, especially when geometric interpretation of the flat-space constraints can be inferred. The presented safe tracking approach can also generalize to other control-affine mechanical systems. In future work, we will explore customized solvers for the SOCP problem involved, and incorporate high-level ``safe corridor'' finding algorithms into the proposed framework.
%apply the framework to other cyber-physical systems such as autonomous cars and search the literature for suitable high-level algorithms that can generate the inputs to the planning module to satisfy some high-level goal in continuous-time. 
%\XX{Need to add a bit more stuff in the conclusion, such as future work.} \VF{Added some sentences at the end.}

% \XX{Double-check all the references: clean the format, ensure capital letters are right, get rid of ``available'', etc. }

\bibliographystyle{IEEEtran}
\bibliography{IEEEabrv,trajbib}

\appendix
%\section{Generating the matrix $B_r$} \label{B_mat}
%\XX{can get rid of the appendix if no space.}

\subsection{Construction of Matrix $B_r$}\label{B_mat}
The matrix $B_r\in\R^{(N+1)\times(N+r+1)}$ is only defined when $0\leq r\leq d$, and can be computed from two matrices $M_{d,d-r}$ and $C_r$ such that $B_r = M_{d,d-r}C_r$; see \cite{fajar_thesis} for more details. 

The matrix $M_{d,d-r}\in\R^{(N+1)\times(N-r+1)}$ can be constructed recursively by:
\begin{align*}
    M_{d,d-r} &= \begin{cases} I_{N+1},  & r= 0,\\
    \prod_{i=1}^rf_M(\vect{\tau},d,N,i), & 1\leq r \leq d,
    \end{cases}
\end{align*}
where each iteration in the product is a right-multiplication and the matrix-valued function $f_M: \mathbb{R}^{v+1}\times \mathbb{Z}_{>0}\times\mathbb{N}\times\mathbb{N} \rightarrow \mathbb{R}^{(N-i+2)\times(N-i+1)}$ is defined as:
\begin{align*}
    f_M(\vect{\tau},d,N,i) 
    & = \mat{-a_0 &\dots & 0\\
              a_0 & \ddots & \vdots\\
              \vdots &\ddots&-a_{N-i}\\
              0 & \hdots & a_{N-i}},\\
    a_k &= \frac{d-i+1}{\tau_{k+d+1}-\tau_{k+i}}.
\end{align*}
 
The matrix $C_r\in\R^{(N-r+1)\times(N+r+1)}$ is constructed as
%\begin{equation*}
    $$C_r = \mat{\vect{0}_{(N+1-r)\times r}& I_{N+1-r}& \vect{0}_{(N+1-r)\times r}}
    $$ 
%\end{equation*}
%where the zero matrices do not exist when $r=0$, i.e., 
for $r\in \mathbb{Z}_{>0}$ and $C_0=I_{N+1}$.
%$\vect{0}_r$ is a zero matrix of dimension $(N+1-r)\times (r)$ which means that it is nonexistent when $r=0$ and $C_0$ reduces to $I_{N+1}$.
%\end{itemize}

% \subsection{Properties of B-splines} \label{bsplineprop} 
% Some properties of B-splines that are  used in this paper are listed below:
% \begin{itemize}
%     \item[P1)] \textit{Continuity}. Given a clamped and uniform knot vector satisfying \eqref{clamped_uniform},  $\vect{s}(t)$ is $C^{\infty}$ for any  $t\notin \vect{\tau}$ and $C^{d-1}$ for any $t\in\vect{\tau}$.
%     \item[P2)] \textit{Convexity}. The B-Spline basis functions have ``partition of unity'' property, meaning that: 
%     %(i) $\lambda_{i,d}(t) \geq 0$; (ii) $\sum_{i=0}^{N}\lambda_{i,d}(t)= 1, \ \forall t\in[\tau_0,\tau_v)$.
%     \begin{itemize}
%         \item[(i)] $\lambda_{i,d}(t) \geq 0$;
%         \item[(ii)] $\sum_{i=0}^{N}\lambda_{i,d}(t)= 1, \ \forall t\in[\tau_0,\tau_v)$.
%     \end{itemize}
%     \item[P3)] \textit{Local Support}. Any B-spline basis is only nonzero over a local set of knots: $\lambda_{i,d}(t)\neq 0$ iff $t\in[\tau_i,\tau_{i+d+1})$.
%     \item[P4)] \textit{Strong Convexity}. Segments of the B-spline curve are contained within the convex hull of a limited set of control points. Combining each segment we obtain 
%     %$ \vect{s}(t) \in \bigcup_{i=d}^{N}\text{Conv}\{\vect{P}_{i-d},\dots,\vect{P}_i\}.$
%     \begin{equation*}
%         \vect{s}(t) \in \bigcup_{i=d}^{N}\text{Conv Hull}\{\vect{P}_{i-d},\dots,\vect{P}_i\}.
%     \end{equation*}
% \end{itemize}

\subsection{Configuration of the Environment for Example \ref{example:forest}}\label{appendix:forest}
% \XX{Revise accordingly.}
The start/end zones $\mathcal{C}_1,\mathcal{C}_2,\mathcal{C}_3,\mathcal{C}_4$ are described by 
%$\mathcal{C}_i \triangleq \{\vect{r}\in\mathbb{R}^3\mid ^{c}\vect{r}\leq \vect{b}_{ci}\}$ where
\begin{align*}
    \mathcal{C}_i &\triangleq \{\vect{r}\in\mathbb{R}^3\mid A\vect{r}\leq \vect{b}^c_i\},\quad i=1,2,3,4,\\
    \mbox{where}\;\quad A &= \mat{I_3 & -I_3}^T,\\
    \mat{{\vect{b}^c_1}^T\\{\vect{b}^c_2}^T\\{\vect{b}^c_3}^T\\ {\vect{b}^{c}_4}^T} &= \mat{1&0.3&1.5&-0.6&0.1&0\\ -0.6&0.6&1.5&1&-0.35&0\\
    -0.6&-0.4&1.5&1&1&0\\1&-0.4&1.5&-0.6&1&0
    }.
\end{align*}
The cuboid obstacles $\mathcal{O}_1,\mathcal{O}_2$ are described by:
\begin{align*}
    \mathcal{O}_i &\triangleq \{\vect{r}\in\mathbb{R}^3\mid A\vect{r}\leq \vect{b}^o_i\},\;i=1,2,
    \end{align*}
    where $A$ is the same as above, and 
    \begin{align*}
   % \mbox{where}\;A\qquad A^o &= \mat{I_3 & -I_3}^T,\\
    \mat{{\vect{b}^o_1}^T\\{\vect{b}^o_2}^T} &= \mat{0.5&0.6&0.7&0.1&-0.15&0\\ 0.1&-0.225&1.1&0.1&0.15&0}.
\end{align*}
The other obstacles $\mathcal{O}_3,\mathcal{O}_4,\mathcal{O}_5$ are:
\begin{itemize}
    \item $\mathcal{O}_3$: a \emph{hoop} centered at $(0, 0,1.1)$ through which the quadcopter can pass.
    \item $\mathcal{O}_4$: a \emph{desk} between $\mathcal{C}_2$ and $\mathcal{C}_3$ so that the quadcopter must fly above or below the \emph{desk} to avoid it.
    \item $\mathcal{O}_5$: a \emph{ladder} between $\mathcal{C}_3$ and $\mathcal{C}_4$ so that the quadcopter must fly between two of its rungs.
\end{itemize}
The considered obstacles allow us to define the safe flight space:
\begin{align*}
    \mathcal{F} &= \mathcal{B} \setminus( \mathcal{O}_1 \cup\cdots\cup \mathcal{O}_5),\\
    \mathcal{B}&\triangleq \{\vect{r}\in\mathbb{R}^3\mid \underline{\vect{b}}\leq\vect{r}\leq \overline{\vect{b}}\},\\
    \underline{\vect{b}} &= \mat{-1 & -1 & 0}^T, \quad \overline{\vect{b}} = \mat{1 & 0.6 & 1.5}.
\end{align*}
Note that $\mathcal{C}_i\in\mathcal{F}$ for $i=1,\ldots,4$. We choose 6 convex sets $\mathcal{S}_j\in\mathcal{F},\ j=1,\ldots,6$, defined as:
\begin{align*}
    \mathcal{S}_j &\triangleq \{\vect{r}\in\mathbb{R}^3\mid A_j\vect{r}\leq \vect{b}^s_j\}, \quad j = 1,\ldots,5,
\end{align*}
where $A_1,\ldots,A_5=A$ and 
\begin{align*}
%    A_j^s &= \mat{I_3 & -I_3}^T, \quad j=1,\dots,5,\\
    \mat{{\vect{b}^s_1}^T\\{\vect{b}^s_2}^T\\{\vect{b}^s_3}^T\\{\vect{b}^s_4}^T\\{\vect{b}^s_5}^T} &= \mat{
    -0.55&0.4&0.5&0.9&0.6&-0.2\\
    -0.55&0.4&1.3&0.9&0.6&-0.85\\
    0.8&-0.45&1.3&0.8&0.65&-1.2\\ 0.8&-0.45&0.47&0.8&0.65&-0.37\\ 0.9&0&1.3&-0.7&0.5&-1.15},\\
    \mathcal{S}_6 &\triangleq \{\vect{r}\in\mathbb{R}^3:\|A^s_6\vect{r}+\vect{b}_6^s\|\leq 1\},\\
    A_6^s &= \text{diag}(1.33,13.3,13.3),\\
    \vect{b}_6^s &= \mat{-0.067 & 0 & -14.67}^T.
\end{align*}




%\begin{IEEEbiography}[{\includegraphics[width=1in,height=1.25in,clip,keepaspectratio]{Bio/Freire_Headshot.jpg}}]%
%{Victor Freire}
%received the B.S. degree in mechanical engineering from University of Wisconsin-Madison, in 2020. He is currently pursuing the M.S. degree from the Department of Mechanical Engineering, University of Wisconsin-Madison.
%He joined ARC Lab in 2019. His research interests include nonlinear
%control, UAV systems, and trajectory planning.
%\end{IEEEbiography}
%
%\begin{IEEEbiography}[{\includegraphics[width=1in,height=1.25in,clip,keepaspectratio]{Bio/xiangru.png}}]
%{Xiangru Xu} is an Assistant Professor in the Department of Mechanical Engineering at the University of Wisconsin-Madison, Madison, Wisconsin, USA. He received his B.S. degree from Beijing Normal University, Beijing, and his Ph.D. degree from Chinese Academy of Sciences, Beijing. Before joining UW-Madison, he was a postdoctoral scholar in the Department of Electrical Engineering and Computer Science at the University of Michigan, Ann Arbor, and the Department of Aeronautics \& Astronautics at the University of Washington, Seattle. His research interests include safety-critical control, nonlinear control, and autonomy.
%\end{IEEEbiography}
% \begin{IEEEbiography}[{\includegraphics[width=1in,height=1.25in,clip,keepaspectratio]{a.png}}]%
% {Xiangru Xu}
% received...
% \end{IEEEbiography}


\end{document}
